\input{preamble}

% OK, start here.
%
\begin{document}

\title{Schémas}


\maketitle

\phantomsection
\label{section-phantom}

\tableofcontents

\section{Introduction}
\label{section-introduction}

\noindent
Dans ce document, nous définissons les schémas. Une référence de base est
 \cite{EGA}.









\section{Espaces localement annelés}
\label{section-locally-ringed-spaces}

\noindent
Rappelons que nous avons défini les espaces annelés dans Faisceaux, Section
 \ref{sheaves-section-ringed-spaces}. En bref, un espace annelé est une paire  $(X, \mathcal{O}_X)$  composée d'
un espace topologique  $X$  et d'un faisceau d'anneaux  $\mathcal{O}_X$. Un
morphisme d'espaces annelés  $f : (X, \mathcal{O}_X) \to (Y, \mathcal{O}_Y)$  est donné par une application continue
 $f : X \to Y$  et un  $f$-morphisme de faisceaux d'anneaux  $f^\sharp : \mathcal{O}_Y \to \mathcal{O}_X$.
Vous pouvez considérer  $f^\sharp$  comme une application  $\mathcal{O}_Y \to f_*\mathcal{O}_X$, voir
Faisceaux, Définition  \ref{sheaves-definition-f-map}  et Lemme  \ref{sheaves-lemma-f-map}.

\medskip\noindent
Un bon exemple géométrique à garder à l'esprit est celui des
 $\mathcal{C}^\infty$-variétés et des morphismes de  $\mathcal{C}^\infty$-variétés. À savoir, si
 $M$  est une  $\mathcal{C}^\infty$-variété, alors le faisceau  $\mathcal{C}^\infty_M$  des
fonctions lisses est un faisceau d'anneaux sur  $M$. Et toute application
 $f : M \to N$  de variétés est lisse si et seulement si pour chaque section locale
 $h$  de  $\mathcal{C}^\infty_N$ , la composition  $h \circ f$  est une section locale
de  $\mathcal{C}^\infty_M$. Ainsi, une application lisse  $f$  donne lieu de manière
naturelle à un morphisme d'espaces annelés.
$$
f : (M , \mathcal{C}^\infty_M) \longrightarrow (N, \mathcal{C}^\infty_N)
$$
voir Faisceaux, Exemple  \ref{sheaves-example-continuous-map-ringed}. Il est instructif de considérer ce qui se
passe pour les fibres. À savoir, soit  $m \in M$  avec image  $f(m) = n \in N$.
Rappelons que la fibre  $\mathcal{C}^\infty_{M, m}$  est l'anneau des germes de fonctions lisses en
 $m$, voir Faisceaux, Exemple  \ref{sheaves-example-germs-functions}. L'algèbre des germes de
fonctions sur  $(M, m)$  est un anneau local avec idéal maximal les fonctions
qui s'annulent en  $m$. De même pour  $\mathcal{C}^\infty_{N, n}$. L'application sur les
fibres  $f^\sharp : \mathcal{C}^\infty_{N, n} \to \mathcal{C}^\infty_{M, m}$  envoie l'idéal maximal dans l'idéal maximal, simplement parce
que  $f(m) = n$.

\medskip\noindent
En géométrie algébrique, nous étudions les schémas. Sur un schéma, le faisceau
d'anneaux n'est pas déterminé par une propriété intrinsèque de l'espace. Le
spectre d'un anneau  $R$  (voir Algèbre, Section  \ref{algebra-section-spectrum-ring}) doté d'un
faisceau d'anneaux construit à partir de  $R$  (voir ci-dessous), sera
notre élément de base. Il s'avérera que les fibres de  $\mathcal{O}$  sur
 $\Spec(R)$  sont les anneaux locaux de  $R$  en ses idéaux premiers. Il y
a deux raisons d'introduire des espaces localement annelés dans ce contexte: (1)
Il n'existe généralement aucun mécanisme qui assigne à une application continue
des spectres une application des anneaux correspondants. C'est pourquoi nous
ajoutons comme donnée supplémentaire le morphisme  $f^\sharp$. (2) Si nous
considérons les morphismes de ces spectres dans la catégorie des espaces annelés,
alors les homomorphismes sur les fibres peuvent ne pas être des homomorphismes
locaux. Comme notre intuition géométrique indique que cela devrait être le cas,
nous introduisons les espaces localement annelés comme suit.

\begin{definition}
\label{definition-locally-ringed-space}
Espaces localement annelés.
\begin{enumerate}
\item Un espace localement annelé  {\it   $(X, \mathcal{O}_X)$}  est
une paire composée d'un espace topologique  $X$  et d'un faisceau d'anneaux
 $\mathcal{O}_X$  dont toutes les fibres sont des anneaux locaux.
\item Étant donné un espace localement annelé  $(X, \mathcal{O}_X)$ , nous disons que  $\mathcal{O}_{X, x}$ 
est l'{\it anneau local de  $X$  en
 $x$}. Nous notons  $\mathfrak{m}_{X, x}$  ou simplement  $\mathfrak{m}_x$ 
l'idéal maximal de  $\mathcal{O}_{X, x}$. De plus, le corps résiduel
 {\it  de  $X$  en  $x$}  est le corps
résiduel  $\kappa(x) = \mathcal{O}_{X, x}/\mathfrak{m}_x$.
\item Un  {\it  morphisme d'espaces localement annelés } 
 $(f, f^\sharp) : (X, \mathcal{O}_X) \to (Y, \mathcal{O}_Y)$  est un morphisme d'espaces annelés tel que pour tout  $x \in X$ ,
l'homomorphisme d'anneaux induit  $\mathcal{O}_{Y, f(x)} \to \mathcal{O}_{X, x}$  est un homomorphisme local.
\end{enumerate}
\end{definition}

\noindent
Nous supprimerons généralement le faisceau d'anneaux  $\mathcal{O}_X$  dans la notation
lorsque nous discutons des espaces localement annelés. Nous nous référerons
simplement à « l'espace localement annelé  $X$ ». Par abus de notation,
nous considérerons également  $X$  comme l'espace topologique sous-jacent.
Enfin, nous noterons le faisceau d'anneaux correspondant  $\mathcal{O}_X$  comme le
faisceau structural  {\it  de  $X$}. De plus,
il est d’usage de noter l'idéal maximal de l'anneau local  $\mathcal{O}_{X, x}$  par
 $\mathfrak{m}_{X, x}$  ou simplement  $\mathfrak{m}_x$. Nous dirons « soit  $f : X \to Y$  un
morphisme d'espaces localement annelés » en supprimant encore davantage les
faisceaux structuraux. Dans ce cas, nous penserons par abus de langage à
 $f : X\to Y$  également comme l'application continue sous-jacente des espaces
topologiques. Le  $f$-morphisme de faisceaux correspondant à  $f$  sera
habituellement noté  $f^\sharp$. La condition que  $f$  soit un morphisme
d'espaces localement annelés peut alors s'exprimer en disant que pour chaque
 $x\in X$  l'homomorphisme sur les fibres
$$
f^\sharp_x : \mathcal{O}_{Y, f(x)} \longrightarrow \mathcal{O}_{X, x}
$$
envoie l'idéal maximal  $\mathfrak m_{Y, f(x)}$  dans  $\mathfrak m_{X, x}$.

\medskip\noindent
Utilisons ces conventions de notation pour montrer que la collection des espaces
localement annelés et des morphismes d'espaces localement annelés forme une
catégorie. Pour voir cela, nous devons montrer que la composition des morphismes
d'espaces localement annelés est un morphisme d'espaces localement annelés.
Soient  $f : X \to Y$  et  $g : Y \to Z$  des morphismes d'espaces localement
annelés. La composition de  $f$  et  $g$  est définie dans
Faisceaux, Définition  \ref{sheaves-definition-composition-maps-ringed-spaces}. Soit  $x \in X$. D'après Faisceaux, Lemme
 \ref{sheaves-lemma-compose-f-maps-stalks} , la composition
$$
\mathcal{O}_{Z, g(f(x))}
\xrightarrow{g^\sharp}
\mathcal{O}_{Y, f(x)}
\xrightarrow{f^\sharp}
\mathcal{O}_{X, x}
$$
est l'homomorphisme sur les fibres associé au morphisme  $g \circ f$. Le résultat
s'ensuit puisque la composition d'homomorphismes locaux d'anneaux est un
homomorphisme local d'anneaux.

\medskip\noindent
Une caractéristique agréable de la définition est le fait que le foncteur
$$
\textit{Espaces localement annelés}
\longrightarrow
\textit{Espaces annelés}
$$
reflète les isomorphismes (et plus encore). Voici un énoncé moins abstrait.

\begin{lemma}
\label{lemma-isomorphism-locally-ringed}
\begin{slogan}
Un isomorphisme d'espaces annelés entre des espaces localement annelés est un
isomorphisme d'espaces localement annelés.
\end{slogan}
Soient  $X$,  $Y$  des espaces localement annelés. Si  $f : X \to Y$ 
est un isomorphisme d'espaces annelés, alors  $f$  est un isomorphisme
d'espaces localement annelés.
\end{lemma}

\begin{proof}
Cela découle trivialement du fait correspondant en algèbre: Supposons que
 $A$,  $B$  soient des anneaux locaux. Tout isomorphisme d'anneaux
 $A \to B$  est un isomorphisme d'anneaux locaux.
\end{proof}













\section{Immersions ouvertes d'espaces localement annelés}
\label{section-open-immersion}

\begin{definition}
\label{definition-immersion-locally-ringed-spaces}
Soit  $f : X \to Y$  un morphisme d'espaces localement annelés. Nous disons que
 $f$  est une  {\it  immersion ouverte}  si
 $f$  est un homéomorphisme de  $X$  sur un sous-ensemble ouvert de
 $Y$, et que l'application  $f^{-1}\mathcal{O}_Y \to \mathcal{O}_X$  est un isomorphisme.
\end{definition}

\noindent
La construction suivante est parallèle à Faisceaux, Définition  \ref{sheaves-definition-restriction}  (3).

\begin{example}
\label{example-open-subspace}
Soit  $X$  un espace localement annelé. Soit  $U \subset X$  un ouvert. Soit
 $\mathcal{O}_U = \mathcal{O}_X|_U$  la restriction de  $\mathcal{O}_X$  à  $U$. Pour  $u \in U$ , la
fibre  $\mathcal{O}_{U, u}$  est égale à la fibre  $\mathcal{O}_{X, u}$, et est donc un anneau local.
Ainsi  $(U, \mathcal{O}_U)$  est un espace localement annelé et le morphisme  $j : (U, \mathcal{O}_U) \to (X, \mathcal{O}_X)$ 
est une immersion ouverte.
\end{example}

\begin{definition}
\label{definition-open-subspace}
Soit  $X$  un espace localement annelé. Soit  $U \subset X$  un ouvert.
L'espace localement annelé  $(U, \mathcal{O}_U)$  de l'exemple  \ref{example-open-subspace}  ci-dessus est le
sous-espace ouvert  {\it  de  $X$  associé à
 $U$}.
\end{definition}

\begin{lemma}
\label{lemma-open-immersion}
Soit  $f : X \to Y$  une immersion ouverte d'espaces localement annelés. Soit
 $j : V = f(X) \to Y$  le sous-espace ouvert de  $Y$  associé à l'image de
 $f$. Il existe un isomorphisme unique  $f' : X \cong V$  d'espaces localement
annelés tel que  $f = j \circ f'$.
\end{lemma}

\begin{proof}
Soit  $f'$  l'homéomorphisme entre  $X$  et  $V$  induit par
 $f$. Alors  $f = j \circ f'$  en tant qu'applications d'espaces topologiques.
Comme il existe un isomorphisme de faisceaux  $f^\sharp : f^{-1}(\mathcal{O}_Y) \to \mathcal{O}_X$, il existe un
isomorphisme d'anneaux  $f^\sharp : \Gamma(U, f^{-1}(\mathcal{O}_Y)) \to \Gamma(U, \mathcal{O}_X)$  pour chaque ouvert  $U \subset X$. Puisque
 $\mathcal{O}_V = j^{-1}\mathcal{O}_Y$  et  $f^{-1} = f'^{-1} j^{-1}$  (Faisceaux, Lemme  \ref{sheaves-lemma-pullback-composition}) nous voyons que
 $f^{-1}\mathcal{O}_Y = f'^{-1}\mathcal{O}_V$, donc  $\Gamma(U, f'^{-1}(\mathcal{O}_V)) \to \Gamma(U, f^{-1}(\mathcal{O}_Y))$  est un isomorphisme pour chaque  $U \subset X$ 
ouvert. En composant ces derniers, nous obtenons un isomorphisme d'anneaux
$$
\Gamma(U, f'^{-1}(\mathcal{O}_V)) \to \Gamma(U, \mathcal{O}_X)
$$
pour chaque ouvert  $U \subset X$, et donc un isomorphisme de faisceaux
 $f^{-1}(\mathcal{O}_V) \to \mathcal{O}_X$. En d'autres termes, nous avons un isomorphisme  $f'^{\sharp} : f'^{-1}(\mathcal{O}_V) \to \mathcal{O}_X$  et donc
un isomorphisme d'espaces localement annelés  $(f', f'^{\sharp}) : (X, \mathcal{O}_X) \to (V, \mathcal{O}_V)$  (utilisez le Lemme
 \ref{lemma-isomorphism-locally-ringed}). Notez que  $f = j \circ f'$  en tant que morphismes d'espaces localement
annelés par construction.

\medskip\noindent
Supposons que nous ayons un autre morphisme  $f'' : (X, \mathcal{O}_X) \to (V, \mathcal{O}_V)$  tel que  $f = j \circ f''$. En
tout point  $x \in X$, nous avons  $j(f'(x)) = j(f''(x))$  d'où il découle que  $f'(x) = f''(x)$ 
puisque  $j$  est l'application d'inclusion; par conséquent  $f'$  et
 $f''$  sont les mêmes en tant que morphismes d'espaces topologiques. Sur les
faisceaux structuraux, pour chaque ouvert  $U \subset X$  nous avons un diagramme
commutatif
$$
\xymatrix @R=5em{
\Gamma(U, f^{-1}(\mathcal{O}_Y)) \ar[d]_\cong\ar[r]^\cong &
\Gamma(U, \mathcal{O}_X) \\
\Gamma(U, f'^{-1}(\mathcal{O}_V)) \ar@/^/[ru]^{f'^\sharp}
\ar@/_/[ru]_{f''^\sharp} &
}
$$
d'où nous voyons que  $f'^\sharp$  et  $f''^\sharp$  définissent le même morphisme de
faisceaux.
\end{proof}

\noindent
Dorénavant, nous ne distinguons plus entre les ouverts et leurs sous-espaces
associés.

\begin{lemma}
\label{lemma-restrict-map-to-opens}
Soit  $f : X \to Y$  un morphisme d'espaces localement annelés. Soient  $U \subset X$
et  $V \subset Y$  des ouverts. Supposons que  $f(U) \subset V$. Il existe un unique
morphisme d'espaces localement annelés  $f|_U : U \to V$  tel que le diagramme suivant
soit un carré commutatif d'espaces localement annelés
$$
\xymatrix{
U \ar[d]_{f|_U} \ar[r] & X \ar[d]^f \\
V \ar[r] & Y
}
$$
\end{lemma}

\begin{proof}
Omis.
\end{proof}

\noindent
Dans ce qui suit, nous utiliserons sans autre mention le fait suivant qui découle
du lemme ci-dessus. Étant donné un morphisme  $f : Y \to X$  d'espaces localement
annelés, et tout ouvert  $U \subset X$  tel que  $f(Y) \subset U$, alors il existe un
unique morphisme d'espaces localement annelés  $Y \to U$  tel que la composition
 $Y \to U \to X$  soit égale à  $f$. En fait, nous écrirons même par abus de
langage  $f : Y \to U$  puisque ceci donne rarement lieu à confusion.









\section{Immersions fermées d'espaces localement annelés}
\label{section-closed-immersion}

\noindent
Nous suivons nos conventions introduites dans Modules, Définition  \ref{modules-definition-closed-immersion}.

\begin{definition}
\label{definition-closed-immersion-locally-ringed-spaces}
Soit  $i : Z \to X$  un morphisme d'espaces localement annelés. Nous disons que
 $i$  est une  {\it  immersion fermée}  si:
\begin{enumerate}
\item L'application  $i$  est un homéomorphisme de  $Z$  sur un
sous-ensemble fermé de  $X$.
\item L'application  $\mathcal{O}_X \to i_*\mathcal{O}_Z$  est surjective; notons  $\mathcal{I}$  le noyau.
\item Le  $\mathcal{O}_X$-module  $\mathcal{I}$  est localement engendré par des sections.
\end{enumerate}
\end{definition}

\begin{lemma}
\label{lemma-closed-local-target}
Soit  $f : Z \to X$  un morphisme d'espaces localement annelés. Pour que  $f$ 
soit une immersion fermée, il suffit qu'il existe un recouvrement ouvert
 $X = \bigcup U_i$  tel que chaque  $f : f^{-1}U_i \to U_i$  soit une immersion fermée.
\end{lemma}

\begin{proof}
Omis.
\end{proof}

\begin{example}
\label{example-closed-subspace}
Soit  $X$  un espace localement annelé. Soit  $\mathcal{I} \subset \mathcal{O}_X$  un faisceau
d'idéaux qui est localement engendré par des sections en tant que faisceau de
 $\mathcal{O}_X$-modules. Soit  $Z$  le support du faisceau d'anneaux
 $\mathcal{O}_X/\mathcal{I}$. Il s'agit d'un sous-ensemble fermé de  $X$, d'après Modules,
Lemme  \ref{modules-lemma-support-sheaf-rings-closed}. On note  $i : Z \to X$  l'application d'inclusion. D'après
Modules, Lemme  \ref{modules-lemma-i-star-exact} , il existe un unique faisceau d'anneaux  $\mathcal{O}_Z$ 
sur  $Z$  avec  $i_*\mathcal{O}_Z = \mathcal{O}_X/\mathcal{I}$. Pour tout  $z \in Z$ , la fibre  $\mathcal{O}_{Z, z}$ 
est égale à un quotient  $\mathcal{O}_{X, i(z)}/\mathcal{I}_{i(z)}$  d'un anneau local et non nul, donc un anneau
local. Ainsi,  $i : (Z, \mathcal{O}_Z) \to (X, \mathcal{O}_X)$  est une immersion fermée d'espaces localement annelés.
\end{example}

\begin{definition}
\label{definition-closed-subspace}
Soit  $X$  un espace localement annelé. Soit  $\mathcal{I}$  un faisceau
d'idéaux sur  $X$  qui est localement engendré par des sections. L'espace
localement annelé  $(Z, \mathcal{O}_Z)$  de l'exemple  \ref{example-closed-subspace}  ci-dessus est le sous-espace
fermé  {\it  de  $X$  associé au faisceau d'idéaux
 $\mathcal{I}$}.
\end{definition}

\begin{lemma}
\label{lemma-closed-immersion}
Soit  $f : X \to Y$  une immersion fermée d'espaces localement annelés. Soit
 $\mathcal{I}$  le noyau de l'application  $\mathcal{O}_Y \to f_*\mathcal{O}_X$. Soit  $i : Z \to Y$  le sous-espace
fermé de  $Y$  associé à  $\mathcal{I}$. Il existe un isomorphisme unique
 $f' : X \cong Z$  d'espaces localement annelés tel que  $f = i \circ f'$.
\end{lemma}

\begin{proof}
Omis.
\end{proof}

\begin{lemma}
\label{lemma-characterize-closed-subspace}
Soient  $X$,  $Y$  des espaces localement annelés. Soit
 $\mathcal{I} \subset \mathcal{O}_X$  un faisceau d'idéaux localement engendré par des sections. Soit
 $i : Z \to X$  le sous-espace fermé associé. Un morphisme  $f : Y \to X$  se factorise
par  $Z$  si et seulement si l'application  $f^*\mathcal{I} \to f^*\mathcal{O}_X = \mathcal{O}_Y$  est nulle. Si tel
est le cas, le morphisme  $g : Y \to Z$  tel que  $f = i \circ g$  est unique.
\end{lemma}

\begin{proof}
De toute évidence, si  $f$  se factorise en  $Y \to Z \to X$  alors
l'application  $f^*\mathcal{I} \to \mathcal{O}_Y$  est nulle. Inversement, supposons que  $f^*\mathcal{I} \to \mathcal{O}_Y$  soit
nulle. Choisissez un  $y \in Y$ et considérez l'application
d'anneaux  $f^\sharp_y : \mathcal{O}_{X, f(y)} \to \mathcal{O}_{Y, y}$. Puisque la composition  $\mathcal{I}_{f(y)} \to \mathcal{O}_{X, f(y)} \to \mathcal{O}_{Y, y}$  est nulle par
hypothèse et puisque  $f^\sharp_y(1) = 1$ , nous voyons que  $1 \not \in \mathcal{I}_{f(y)}$, c'est-à-dire
 $\mathcal{I}_{f(y)} \not = \mathcal{O}_{X, f(y)}$. Nous en concluons que  $f(Y) \subset Z = \text{Supp}(\mathcal{O}_X/\mathcal{I})$. Ainsi,  $f = i \circ g$  où
 $g : Y \to Z$  est continue. Considérons l'application  $f^\sharp : \mathcal{O}_X \to f_*\mathcal{O}_Y$. L'hypothèse que
 $f^*\mathcal{I} \to \mathcal{O}_Y$  est nulle implique que la composition  $\mathcal{I} \to \mathcal{O}_X \to f_*\mathcal{O}_Y$  est nulle par
adjonction de  $f_*$  et  $f^*$. En d'autres termes, nous obtenons un
morphisme de faisceaux d'anneaux  $\overline{f^\sharp} : \mathcal{O}_X/\mathcal{I} \to f_*\mathcal{O}_Y$. Notez que  $f_*\mathcal{O}_Y = i_*g_*\mathcal{O}_Y$  et que
 $\mathcal{O}_X/\mathcal{I} = i_*\mathcal{O}_Z$. D'après Faisceaux, Lemme  \ref{sheaves-lemma-equivalence-categories-closed-structures} , nous obtenons un
morphisme unique de faisceaux d'anneaux  $g^\sharp : \mathcal{O}_Z \to g_*\mathcal{O}_Y$  dont l'image directe sous
 $i$  est  $\overline{f^\sharp}$. Nous omettons la vérification que  $(g, g^\sharp)$ 
définit un morphisme d'espaces localement annelés et que  $f = i \circ g$  en tant que
morphisme d'espaces localement annelés. L'unicité de  $(g, g^\sharp)$  a été signalée
ci-dessus.
\end{proof}

\begin{lemma}
\label{lemma-restrict-map-to-closed}
Soit  $f : X \to Y$  un morphisme d'espaces localement annelés. Soit  $\mathcal{I} \subset \mathcal{O}_Y$  un
faisceau d'idéaux qui est localement engendré par des sections. Soit  $i : Z \to Y$ 
le sous-espace fermé associé au faisceau d'idéaux  $\mathcal{I}$. Soit  $\mathcal{J}$ 
l'image de l'application  $f^*\mathcal{I} \to f^*\mathcal{O}_Y = \mathcal{O}_X$. Alors cet idéal est localement engendré par
des sections. De plus, soit  $i' : Z' \to X$  le sous-espace fermé associé de
 $X$. Il existe un unique morphisme d'espaces localement annelés
 $f' : Z' \to Z$  tel que le diagramme suivant soit un carré commutatif d'espaces
localement annelés
$$
\xymatrix{
Z' \ar[d]_{f'} \ar[r]_{i'} & X \ar[d]^f \\
Z \ar[r]^{i} & Y
}
$$
De plus, ce diagramme est un carré cartésien dans la catégorie des espaces
localement annelés.
\end{lemma}

\begin{proof}
L'idéal  $\mathcal{J}$  est localement engendré par des sections d'après Modules, Lemme
 \ref{modules-lemma-pullback-locally-generated}. Le reste du lemme découle de la caractérisation, dans le Lemme
 \ref{lemma-characterize-closed-subspace}  ci-dessus, de ce que signifie qu'un morphisme se factorise à travers
un sous-espace fermé.
\end{proof}














\section{Schémas affines}
\label{section-affine-schemes}

\noindent
Soit  $R$  un anneau. Considérons l'espace topologique  $\Spec(R)$  associé
à  $R$, voir Algèbre, Section  \ref{algebra-section-spectrum-ring}. Nous munirons cet espace d'un
faisceau d'anneaux  $\mathcal{O}_{\Spec(R)}$  et la paire résultante  $(\Spec(R), \mathcal{O}_{\Spec(R)})$  sera un schéma
affine.

\medskip\noindent
Rappelons que  $\Spec(R)$  possède une base d'ouverts  $D(f)$,  $f \in R$ 
que nous appelons ouverts standards, voir Algèbre, Définition  \ref{algebra-definition-Zariski-topology}. De
plus, l'intersection de deux ouverts standards est un autre:  $D(f) \cap D(g) = D(fg)$,
 $f, g\in R$.

\begin{lemma}
\label{lemma-standard-open}
Soit  $R$  un anneau. Soit  $f \in R$.
\begin{enumerate}
\item Si  $g\in R$  et  $D(g) \subset D(f)$, alors
\begin{enumerate}
\item $f$  est inversible dans  $R_g$,
\item $g^e = af$  pour certains  $e \geq 1$  et  $a \in R$,
\item il existe un morphisme d'anneaux canonique  $R_f \to R_g$, et
\item il existe un morphisme de  $R_f$-modules canonique  $M_f \to M_g$  pour tout
 $R$-module  $M$.
\end{enumerate}
\item Tout recouvrement ouvert de  $D(f)$  peut être affiné en un recouvrement
ouvert fini de la forme  $D(f) = \bigcup_{i = 1}^n D(g_i)$.
\item Si  $g_1, \ldots, g_n \in R$, alors  $D(f) \subset \bigcup D(g_i)$  si et seulement si  $g_1, \ldots, g_n$  engendrent
l'idéal unité dans  $R_f$.
\end{enumerate}
\end{lemma}

\begin{proof}
Rappelons que  $D(g) = \Spec(R_g)$  (voir Algèbre, Lemme  \ref{algebra-lemma-standard-open}). Ainsi (a) est
valable parce que l'image de  $f$  dans  $R_g$  n'appartient
à aucun idéal premier, et est donc inversible, voir Algèbre, Lemme
 \ref{algebra-lemma-Zariski-topology}. Écrivons l’inverse de  $f$  dans  $R_g$  comme
 $a/g^d$. Cela signifie que  $g^d - af$  est annihilé par une puissance de
 $g$, d’où (b). Pour (c), l’application  $R_f \to R_g$  existe par (a) à
cause de la propriété universelle de la localisation, ou nous pouvons la définir
en envoyant  $b/f^n$  sur  $a^nb/g^{ne}$. L’égalité  $M_f = M \otimes_R R_f$  peut être utilisée
pour obtenir l’application sur les modules, ou nous pouvons définir  $M_f \to M_g$ 
en envoyant  $x/f^n$  sur  $a^nx/g^{ne}$.

\medskip\noindent
Rappelons que  $D(f)$  est quasi-compact, voir Algèbre, Lemme
 \ref{algebra-lemma-qc-open}. Ainsi, la deuxième affirmation découle directement du fait que les
ouverts standards forment une base pour la topologie.

\medskip\noindent
La troisième affirmation découle directement d'Algèbre, Lemme  \ref{algebra-lemma-Zariski-topology}.
\end{proof}

\noindent
Dans Faisceaux, Section  \ref{sheaves-section-bases} , nous avons défini la notion de faisceau sur
une base, et nous avons montré qu'elle est essentiellement équivalente à la notion
de faisceau sur l'espace, voir Faisceaux, lemmes  \ref{sheaves-lemma-extend-off-basis}  et  \ref{sheaves-lemma-extend-off-basis-structures}. De
plus, nous avons montré dans Faisceaux, Lemme  \ref{sheaves-lemma-cofinal-systems-coverings-standard-case}  qu'il suffit de
vérifier la condition de faisceau sur un système cofinal de recouvrements ouverts
pour chaque ouvert standard. Par le lemme ci-dessus, il suffit de vérifier sur les
recouvrements finis par des ouverts standards.

\begin{definition}
\label{definition-standard-covering}
Soit  $R$  un anneau.
\begin{enumerate}
\item Un  {\it  recouvrement ouvert standard}  de
 $\Spec(R)$  est un recouvrement  $\Spec(R) = \bigcup_{i = 1}^n D(f_i)$, où  $f_1, \ldots, f_n \in R$.
\item Supposons que  $D(f) \subset \Spec(R)$  soit un ouvert standard. Un
 {\it  recouvrement ouvert standard}  de  $D(f)$  est un recouvrement
 $D(f) = \bigcup_{i = 1}^n D(g_i)$, où  $g_1, \ldots, g_n \in R$.
\end{enumerate}
\end{definition}

\noindent
Soit  $R$  un anneau. Soit  $M$  un  $R$-module. Nous
définirons un préfaisceau  $\widetilde M$  sur la base des ouverts standards.
Supposons que  $U \subset \Spec(R)$  soit un ouvert standard. Si  $f, g \in R$  sont tels que
 $D(f) = D(g)$, alors par le Lemme  \ref{lemma-standard-open}  ci-dessus, il existe des
applications canoniques  $M_f \to M_g$  et  $M_g \to M_f$  qui sont mutuellement
inverses. Par conséquent, nous pouvons choisir n'importe quel  $f$  tel que
 $U = D(f)$  et définir
$$
\widetilde M(U) = M_f.
$$
Notez que si  $D(g) \subset D(f)$, alors par le Lemme  \ref{lemma-standard-open}  ci-dessus, nous avons
une application canonique
$$
\widetilde M(D(f)) = M_f \longrightarrow M_g = \widetilde M(D(g)).
$$
Clairement, cela définit un préfaisceau de groupes abéliens sur la base des
ouverts standards. Si  $M = R$, alors  $\widetilde R$  est un préfaisceau
d'anneaux sur la base des ouverts standards.

\medskip\noindent
Calculons la fibre de  $\widetilde M$  en un point  $x \in \Spec(R)$. Supposons que
 $x$  corresponde à l'idéal premier  $\mathfrak p \subset R$. Par définition de la fibre, nous
voyons que
$$
\widetilde M_x = \colim_{f\in R, f\not\in \mathfrak p} M_f
$$
Ici, l'ensemble  $\{f \in R, f \not \in \mathfrak p\}$  est préordonné par la règle  $f \geq f' \Leftrightarrow D(f) \subset D(f')$. Si
 $f_1, f_2 \in R \setminus \mathfrak p$, alors nous avons  $f_1f_2 \geq f_1$  dans cet ordre. Ainsi, par Algèbre,
Lemme  \ref{algebra-lemma-localization-colimit} , nous concluons que
$$
\widetilde M_x = M_{\mathfrak p}.
$$

\medskip\noindent
Ensuite, nous vérifions la condition de faisceau pour les recouvrements ouverts
standards. Si  $D(f) = \bigcup_{i = 1}^n D(g_i)$, alors la condition de faisceau pour ce recouvrement
est équivalente à l’exactitude de la suite
$$
0 \to M_f \to \bigoplus M_{g_i} \to \bigoplus M_{g_ig_j}.
$$
Notez que  $D(g_i) = D(fg_i)$, et donc nous pouvons réécrire cette suite comme la suite
$$
0 \to M_f \to \bigoplus M_{fg_i} \to \bigoplus M_{fg_ig_j}.
$$
De plus, d'après le Lemme  \ref{lemma-standard-open}  ci-dessus, nous voyons que  $g_1, \ldots, g_n$ 
engendrent l'idéal unité dans  $R_f$. Nous pouvons donc appliquer Algèbre, Lemme
 \ref{algebra-lemma-cover-module} , au module  $M_f$  sur  $R_f$  et aux éléments
 $g_1, \ldots, g_n$. Nous en concluons que la suite est exacte. D'après les remarques
faites ci-dessus, nous voyons que  $\widetilde M$  est un faisceau sur la base des
ouverts standards.

\medskip\noindent
Ainsi, nous concluons à partir du matériel de Faisceaux, Section  \ref{sheaves-section-bases} 
qu'il existe un unique faisceau d'anneaux  $\mathcal{O}_{\Spec(R)}$  qui coïncide avec
 $\widetilde R$  sur les ouverts standards. Notez que d'après notre calcul des fibres
ci-dessus, toutes les fibres de ce faisceau d'anneaux sont des anneaux locaux.

\medskip\noindent
De même, pour tout  $R$-module  $M$ , il existe un unique faisceau
de  $\mathcal{O}_{\Spec(R)}$-modules  $\mathcal{F}$  qui coïncide avec  $\widetilde M$  sur les
ouverts standards, voir Faisceaux, Lemme  \ref{sheaves-lemma-extend-off-basis-module}.

\begin{definition}
\label{definition-structure-sheaf}
Soit  $R$  un anneau.
\begin{enumerate}
\item Le  {\it faisceau structural $\mathcal{O}_{\Spec(R)}$ du spectre
de $R$} est l'unique faisceau d'anneaux $\mathcal{O}_{\Spec(R)}$ dont la restriction
à la base des ouverts standards coïncide avec $\widetilde R$.
\item L'espace localement annelé $(\Spec(R), \mathcal{O}_{\Spec(R)})$  est appelé le
 {\it spectre}  de  $R$  et désigné
 $\Spec(R)$.
\item Le faisceau de  $\mathcal{O}_{\Spec(R)}$-modules étendant  $\widetilde M$  à tous les ouverts de
 $\Spec(R)$  est appelé le faisceau de  $\mathcal{O}_{\Spec(R)}$-modules associé à
 $M$. Ce faisceau est également noté  $\widetilde M$ .
\end{enumerate}
\end{definition}

\noindent
Nous résumons les résultats obtenus jusqu'à présent.

\begin{lemma}
\label{lemma-spec-sheaves}
Soit  $R$  un anneau. Soit  $M$  un  $R$-module. Soit
 $\widetilde M$  le faisceau de  $\mathcal{O}_{\Spec(R)}$-modules associé à  $M$.
\begin{enumerate}
\item Nous avons  $\Gamma(\Spec(R), \mathcal{O}_{\Spec(R)}) = R$.
\item Nous avons  $\Gamma(\Spec(R), \widetilde M) = M$  en tant que  $R$-module.
\item Pour chaque  $f \in R$  nous avons  $\Gamma(D(f), \mathcal{O}_{\Spec(R)}) = R_f$.
\item Pour chaque  $f\in R$  nous avons  $\Gamma(D(f), \widetilde M) = M_f$  en tant que  $R_f$-module.
\item Chaque fois que  $D(g) \subset D(f)$ , les applications de restriction sur  $\mathcal{O}_{\Spec(R)}$  et
 $\widetilde M$  sont les applications  $R_f \to R_g$  et  $M_f \to M_g$  du Lemme  \ref{lemma-standard-open}.
\item Soit  $\mathfrak p$  un idéal premier de  $R$, et soit  $x \in \Spec(R)$  le point
correspondant. Nous avons  $\mathcal{O}_{\Spec(R), x} = R_{\mathfrak p}$.
\item Soit  $\mathfrak p$  un idéal premier de  $R$, et soit  $x \in \Spec(R)$  le point
correspondant. Nous avons  $\widetilde M_x = M_{\mathfrak p}$  en tant que  $R_{\mathfrak p}$-module.
\end{enumerate}
De plus, toutes ces identifications sont fonctorielles par rapport au  $R$-module
 $M$. En particulier, le foncteur  $M \mapsto \widetilde M$  est un foncteur exact de la
catégorie des  $R$-modules à la catégorie des  $\mathcal{O}_{\Spec(R)}$-modules.
\end{lemma}

\begin{proof}
Les assertions (1) - (7) sont claires d'après la discussion ci-dessus.
L'exactitude du foncteur  $M \mapsto \widetilde M$  découle du fait que le foncteur  $M \mapsto M_{\mathfrak p}$ 
est exact et du fait que l'exactitude des suites exactes courtes peut être
vérifiée sur les fibres, voir Modules, Lemme  \ref{modules-lemma-abelian}.
\end{proof}

\begin{definition}
\label{definition-affine-scheme}
Un  {\it  schéma affine}  est un espace localement
annelé isomorphe en tant qu'espace localement annelé à  $\Spec(R)$  pour un
certain anneau  $R$. Un  {\it  morphisme de schémas
affines}  est un morphisme dans la catégorie des espaces localement
annelés.
\end{definition}

\noindent
Il s'avère que les schémas affines jouent un rôle particulier parmi tous les
espaces localement annelés, ce dont traite la section suivante.


















\section{La catégorie des schémas affines}
\label{section-category-affine-schemes}

\noindent
Notez que si $Y$ est un schéma affine, alors ses points sont en
bijection canonique (une correspondance  $1-1$) avec les idéaux premiers de $\Gamma(Y, \mathcal{O}_Y)$.

\begin{lemma}
\label{lemma-morphism-into-affine-where-point-goes}
Soit  $X$  un espace localement annelé. Soit  $Y$  un schéma affine.
Soit  $f \in \Mor(X, Y)$  un morphisme d'espaces localement annelés. Étant donné un point
 $x \in X$ , considérons les homomorphismes d'anneaux
$$
\Gamma(Y, \mathcal{O}_Y) \xrightarrow{f^\sharp}
\Gamma(X, \mathcal{O}_X) \to \mathcal{O}_{X, x}
$$
Soit  $\mathfrak p \subset \Gamma(Y, \mathcal{O}_Y)$  l'image réciproque de  $\mathfrak m_x$. Soit  $y \in Y$  le point
correspondant. Alors  $f(x) = y$.
\end{lemma}

\begin{proof}
Considérons le diagramme commutatif
$$
\xymatrix{
\Gamma(X, \mathcal{O}_X) \ar[r] &
\mathcal{O}_{X, x} \\
\Gamma(Y, \mathcal{O}_Y) \ar[r] \ar[u] &
\mathcal{O}_{Y, f(x)} \ar[u]
}
$$
(voir la discussion des  $f$-morphismes qui suit Faisceaux, Définition
 \ref{sheaves-definition-f-map}). Puisque la flèche verticale de droite est locale, nous voyons que
 $\mathfrak m_{f(x)}$  est l'image réciproque de  $\mathfrak m_x$. Le résultat s'ensuit.
\end{proof}

\begin{lemma}
\label{lemma-f-open}
Soit  $X$  un espace localement annelé. Soit  $f \in \Gamma(X, \mathcal{O}_X)$. L'ensemble
$$
D(f) = \{x \in X \mid \text{image de }f \not\in \mathfrak m_x\}
$$
est ouvert. De plus,  $f|_{D(f)}$  a un inverse.
\end{lemma}

\begin{proof}
Ceci est un cas particulier de Modules, Lemme  \ref{modules-lemma-s-open}, mais nous donnons
également une démonstration directe. Supposons que  $U \subset X$  et  $V \subset X$ 
soient deux sous-ensembles ouverts tels que  $f|_U$  ait un inverse  $g$ 
et que  $f|_V$  ait un inverse  $h$. Alors clairement  $g|_{U\cap V} = h|_{U\cap V}$. Il
suffit donc de montrer que  $f$  est inversible dans un voisinage ouvert de
tout  $x \in D(f)$. Cela est clair car  $f \not \in \mathfrak m_x$  implique que  $f \in \mathcal{O}_{X, x}$  a un
inverse  $g \in \mathcal{O}_{X, x}$ , ce qui signifie qu'il existe un certain voisinage ouvert
 $x \in U \subset X$  tel que  $g \in \mathcal{O}_X(U)$  et  $g\cdot f|_U = 1$.
\end{proof}

\begin{lemma}
\label{lemma-f-open-affine}
Dans le Lemme  \ref{lemma-f-open}  ci-dessus, si  $X$  est un schéma affine, alors
l'ouvert  $D(f)$  correspond à l'ouvert standard  $D(f)$  défini
précédemment (en Algèbre, Définition  \ref{algebra-definition-spectrum-ring}).
\end{lemma}

\begin{proof}
Omis.
\end{proof}

\begin{lemma}
\label{lemma-morphism-into-affine}
\begin{reference}
Une référence pour ce fait est  \cite[II, Err 1, Prop. 1.8.1]{EGA}  où il est attribué à J. Tate.
\end{reference}
Soit  $X$  un espace localement annelé. Soit  $Y$  un schéma affine.
L'application
$$
\Mor(X, Y)
\longrightarrow
\Hom(\Gamma(Y, \mathcal{O}_Y), \Gamma(X, \mathcal{O}_X))
$$
qui envoie  $f$  sur  $f^\sharp$  (au niveau des sections globales) est
bijective.
\end{lemma}

\begin{proof}
Étant donné que  $Y$  est affine, nous avons  $(Y, \mathcal{O}_Y) \cong (\Spec(R), \mathcal{O}_{\Spec(R)})$  pour un certain
anneau  $R$. Au cours de la démonstration, nous utiliserons des faits sur
 $Y$  et son faisceau structural qui sont des conséquences directes de
ce que nous savons sur le spectre d'un anneau, voir par exemple \ Lemme
 \ref{lemma-spec-sheaves}.

\medskip\noindent
Motivés par les lemmes ci-dessus, nous construisons l'application inverse. Soit
 $\psi_Y : \Gamma(Y, \mathcal{O}_Y) \to \Gamma(X, \mathcal{O}_X)$  un morphisme d'anneaux. Tout d'abord, nous définissons l'application
correspondante des espaces
$$
\Psi : X \longrightarrow Y
$$
selon la règle du Lemme  \ref{lemma-morphism-into-affine-where-point-goes}. En d'autres termes, étant donné
 $x \in X$ , nous définissons  $\Psi(x)$  comme le point de  $Y$ 
correspondant à l'idéal premier dans  $\Gamma(Y, \mathcal{O}_Y)$  qui est l'image réciproque de
 $\mathfrak m_x$  par la composition  $
\Gamma(Y, \mathcal{O}_Y) \xrightarrow{\psi_Y}
\Gamma(X, \mathcal{O}_X) \to
\mathcal{O}_{X, x}
$.

\medskip\noindent
Nous affirmons que l'application  $\Psi : X \to Y$  est continue. Les ouverts standards
 $D(g)$, pour  $g \in \Gamma(Y, \mathcal{O}_Y)$  sont une base pour la topologie de  $Y$.
Ainsi, il suffit de prouver que  $\Psi^{-1}(D(g))$  est ouvert. Par construction de
 $\Psi$ , l'image inverse  $\Psi^{-1}(D(g))$  est exactement l'ensemble  $D(\psi_Y(g)) \subset X$ 
qui est ouvert par le Lemme  \ref{lemma-f-open}. Par conséquent,  $\Psi$  est
continue.

\medskip\noindent
Ensuite, nous construisons un  $\Psi$-morphisme de faisceaux de  $\mathcal{O}_Y$ 
à  $\mathcal{O}_X$. Par Faisceaux, Lemme  \ref{sheaves-lemma-f-map-basis-below-structures} , il suffit de définir des
homomorphismes d'anneaux  $\psi_{D(g)} : \Gamma(D(g), \mathcal{O}_Y) \to
\Gamma(\Psi^{-1}(D(g)), \mathcal{O}_X)$  compatibles avec les applications de
restriction. Nous avons un isomorphisme canonique  $\Gamma(D(g), \mathcal{O}_Y) = \Gamma(Y, \mathcal{O}_Y)_g$, car  $Y$ 
est un schéma affine. Comme  $\psi_Y(g)$  est inversible sur  $D(\psi_Y(g))$ , nous
voyons qu'il existe une application canonique
$$
\Gamma(Y, \mathcal{O}_Y)_g
\longrightarrow
\Gamma(\Psi^{-1}(D(g)), \mathcal{O}_X)
=
\Gamma(D(\psi_Y(g)), \mathcal{O}_X)
$$
en étendant l'application  $\psi_Y$  par la propriété universelle de la
localisation. Notez qu'il n'y a pas d'autre choix que de prendre l'application
canonique ici ! Après l'identification canonique
 $\Gamma(D(g), \mathcal{O}_Y) = \Gamma(Y, \mathcal{O}_Y)_g$, nous prenons cette application pour définir  $\psi_{D(g)}$. Cela est compatible avec la localisation
puisque les applications de restriction sur les schémas affines sont également
définies en termes des propriétés universelles de la localisation, voir les
lemmes  \ref{lemma-spec-sheaves}  et  \ref{lemma-standard-open}.

\medskip\noindent
Ainsi, nous avons défini un morphisme d'espaces annelés  $(\Psi, \psi) : (X, \mathcal{O}_X) \to (Y, \mathcal{O}_Y)$  induisant
 $\psi_Y$  sur les sections globales. Pour montrer qu'il s'agit d'un morphisme
d'espaces localement annelés, nous devons démontrer que les applications induites
sur les anneaux locaux
$$
\psi_x : \mathcal{O}_{Y, \Psi(x)} \longrightarrow \mathcal{O}_{X, x}
$$
sont locales. Cela découle immédiatement du diagramme commutatif de la
démonstration du Lemme  \ref{lemma-morphism-into-affine-where-point-goes}  et de la définition de  $\Psi$.

\medskip\noindent
Enfin, nous devons montrer que les constructions  $(\Psi, \psi) \mapsto \psi_Y$  et la construction
 $\psi_Y \mapsto (\Psi, \psi)$  sont inverses l'une de l'autre. Clairement,  $\psi_Y \mapsto (\Psi, \psi) \mapsto \psi_Y$. Par
conséquent, la seule chose à prouver est que, étant donné  $\psi_Y$ , il existe
au plus une paire  $(\Psi, \psi)$  donnant lieu à celle-ci. L'unicité de  $\Psi$ 
a été montrée dans le Lemme  \ref{lemma-morphism-into-affine-where-point-goes}  et étant donné l'unicité de  $\Psi$ ,
l'unicité de l'application  $\psi$  a été soulignée au cours de la
démonstration ci-dessus.
\end{proof}

\begin{lemma}
\label{lemma-category-affine-schemes}
La catégorie des schémas affines est équivalente à l'opposée de la catégorie des
anneaux. L'équivalence est donnée par le foncteur qui associe à un schéma affine
les sections globales de son faisceau structural.
\end{lemma}

\begin{proof}
Cela découle maintenant clairement de la Définition  \ref{definition-affine-scheme}  et du Lemme
 \ref{lemma-morphism-into-affine}.
\end{proof}

\begin{lemma}
\label{lemma-standard-open-affine}
Soit  $Y$  un schéma affine. Soit  $f \in \Gamma(Y, \mathcal{O}_Y)$. Le sous-espace ouvert
 $D(f)$  est un schéma affine.
\end{lemma}

\begin{proof}
Nous pouvons supposer que  $Y = \Spec(R)$  et  $f \in R$. Considérons le morphisme
de schémas affines  $\phi : U = \Spec(R_f) \to \Spec(R) = Y$  induit par l'homomorphisme d'anneaux  $R \to R_f$.
D'après Algèbre, Lemme  \ref{algebra-lemma-standard-open} , nous savons que c'est un homéomorphisme
sur  $D(f)$. D'autre part, l'application  $\phi^{-1}\mathcal{O}_Y \to \mathcal{O}_U$  est un isomorphisme sur
les fibres, donc un isomorphisme. Ainsi, nous voyons que  $\phi$  est une
immersion ouverte. Nous concluons que  $D(f)$  est isomorphe à  $U$ 
d'après le Lemme  \ref{lemma-open-immersion}.
\end{proof}

\begin{lemma}
\label{lemma-fibre-product-affine-schemes}
La catégorie des schémas affines a des produits finis et des produits fibrés. En
d'autres termes, elle possède des limites finies. De plus, les produits et
produits fibrés dans la catégorie des schémas affines sont les mêmes que dans la
catégorie des espaces localement annelés. Formellement, nous avons (dans la
catégorie des espaces localement annelés)
$$
\Spec(R) \times \Spec(S) =
\Spec(R \otimes_{\mathbf{Z}} S)
$$
et étant donné des homomorphismes d'anneaux  $R \to A$,  $R \to B$  nous avons
$$
\Spec(A) \times_{\Spec(R)} \Spec(B)
=
\Spec(A \otimes_R B).
$$
\end{lemma}

\begin{proof}
Ceci est simplement une application du Lemme  \ref{lemma-morphism-into-affine}. Tout d'abord, d'après
ce lemme, le schéma affine  $\Spec(\mathbf{Z})$  est l'objet final dans la catégorie des
espaces localement annelés. Ainsi, la première formule affichée découle de la
seconde. Pour prouver la seconde, notez que pour tout espace localement annelé
 $X$ , nous avons
\begin{eqnarray*}
\Mor(X, \Spec(A \otimes_R B))
& = &
\Hom(A \otimes_R B, \mathcal{O}_X(X)) \\
& = &
\Hom(A, \mathcal{O}_X(X))
\times_{\Hom(R, \mathcal{O}_X(X))}
\Hom(B, \mathcal{O}_X(X)) \\
& = &
\Mor(X, \Spec(A))
\times_{\Mor(X, \Spec(R))}
\Mor(X, \Spec(B))
\end{eqnarray*}
Cela prouve la formule. Voir Catégories, Section  \ref{categories-section-fibre-products}  pour les
définitions pertinentes.
\end{proof}

\begin{lemma}
\label{lemma-disjoint-union-affines}
Soit  $X$  un espace localement annelé. Supposons  $X = U \amalg V$  avec
 $U$  et  $V$  ouverts et tels que  $U$,  $V$  soient
des schémas affines. Alors  $X$  est un schéma affine.
\end{lemma}

\begin{proof}
Fixez  $R = \Gamma(X, \mathcal{O}_X)$. Notez que  $R = \mathcal{O}_X(U) \times \mathcal{O}_X(V)$  par la propriété du faisceau. Par le
Lemme  \ref{lemma-morphism-into-affine} , il existe un morphisme canonique d'espaces localement annelés
 $X \to \Spec(R)$. Par Algèbre, Lemme  \ref{algebra-lemma-spec-product} , nous voyons que comme espace
topologique  $\Spec(\mathcal{O}_X(U)) \amalg \Spec(\mathcal{O}_X(V)) =
\Spec(R)$  avec les applications provenant des homomorphismes
d'anneaux  $R \to \mathcal{O}_X(U)$  et  $R \to \mathcal{O}_X(V)$. Cela signifie bien sûr que  $\Spec(R)$ 
est également le coproduit dans la catégorie des espaces localement annelés. Par
hypothèse, le morphisme  $X \to \Spec(R)$  induit un isomorphisme de  $\Spec(\mathcal{O}_X(U))$  avec
 $U$  et de même pour  $V$. Par conséquent,  $X \to \Spec(R)$  est un
isomorphisme.
\end{proof}















\section{Faisceaux quasi-cohérents sur les schémas affines}
\label{section-quasi-coherent-affine}

\noindent
Rappelons que nous avons défini la notion abstraite de faisceau quasi-cohérent
dans Modules, Définition  \ref{modules-definition-quasi-coherent}. Dans cette section, nous montrons que tout
faisceau quasi-cohérent sur un schéma affine  $\Spec(R)$  correspond au faisceau
 $\widetilde M$  associé à un  $R$-module  $M$.

\begin{lemma}
\label{lemma-compare-constructions}
Soit  $(X, \mathcal{O}_X) = (\Spec(R), \mathcal{O}_{\Spec(R)})$  un schéma affine. Soit  $M$  un  $R$-module. Il
existe un isomorphisme canonique entre le faisceau  $\widetilde M$  associé au
 $R$-module  $M$  (Définition  \ref{definition-structure-sheaf}) et le faisceau
 $\mathcal{F}_M$  associé au  $R$-module  $M$  (Modules, Définition
 \ref{modules-definition-sheaf-associated}). Cet isomorphisme est fonctoriel en  $M$. En particulier,
les faisceaux  $\widetilde M$  sont quasi-cohérents. De plus, ils sont caractérisés
par la propriété universelle suivante
$$
\Hom_{\mathcal{O}_X}(\widetilde M, \mathcal{F})
=
\Hom_R(M, \Gamma(X, \mathcal{F}))
$$
pour tout faisceau de  $\mathcal{O}_X$-modules  $\mathcal{F}$. Ici, une application
 $\alpha : \widetilde M \to \mathcal{F}$  correspond à son effet sur les sections globales.
\end{lemma}

\begin{proof}
Par Modules, Lemme  \ref{modules-lemma-construct-quasi-coherent-sheaves} , nous avons un morphisme  $\mathcal{F}_M \to \widetilde M$ 
correspondant à l'application  $M \to \Gamma(X, \widetilde M) = M$. Soit  $x \in X$  correspondant à
l'idéal premier  $\mathfrak p \subset R$. L'homomorphisme induit sur les fibres est le morphisme
 $\mathcal{O}_{X, x} \otimes_R M \to M_{\mathfrak p}$  qui est un isomorphisme car  $R_{\mathfrak p} \otimes_R M = M_{\mathfrak p}$. Par conséquent,
l'application  $\mathcal{F}_M \to \widetilde M$  est un isomorphisme. La propriété universelle découle
de celle des faisceaux  $\mathcal{F}_M$.
\end{proof}

\begin{lemma}
\label{lemma-widetilde-constructions}
Soit  $(X, \mathcal{O}_X) = (\Spec(R), \mathcal{O}_{\Spec(R)})$  un schéma affine. Il existe des isomorphismes canoniques
\begin{enumerate}
\item $
\widetilde{M \otimes_R N}
\cong
\widetilde M \otimes_{\mathcal{O}_X} \widetilde N
$, voir Modules, Section  \ref{modules-section-tensor-product}.
\item $
\widetilde{\text{T}^n(M)}
\cong
\text{T}^n(\widetilde M)
$,  $
\widetilde{\text{Sym}^n(M)}
\cong
\text{Sym}^n(\widetilde M)
$, et  $
\widetilde{\wedge^n(M)}
\cong
\wedge^n(\widetilde M)
$, voir Modules, Section  \ref{modules-section-symmetric-exterior}.
\item si  $M$  est un  $R$-module de présentation finie, alors
 $
\SheafHom_{\mathcal{O}_X}(\widetilde M, \widetilde N)
\cong
\widetilde{\Hom_R(M,  N)}
$, voir Modules, Section  \ref{modules-section-internal-hom}.
\end{enumerate}
\end{lemma}

\begin{proof}[Première démonstration]
En utilisant le Lemme  \ref{lemma-compare-constructions}  et Modules, Lemme  \ref{modules-lemma-construct-quasi-coherent-sheaves} , nous
voyons que le foncteur  $M \mapsto \widetilde M$  peut être vu comme  $\pi^*$  pour un
morphisme  $\pi$  d'espaces annelés. Le foncteur image inverse sur les modules
commute avec les constructions tensorielles d'après Modules, Lemmes  \ref{modules-lemma-tensor-product-pullback}  et
 \ref{modules-lemma-pullback-tensor-algebra}. Le morphisme  $\pi : (X, \mathcal{O}_X) \to (\{*\}, R)$  est plat par exemple parce que les fibres
de  $\mathcal{O}_X$  sont des localisations de  $R$  (Lemme  \ref{lemma-spec-sheaves}) et
donc plates sur  $R$. Ainsi, l'image inverse par  $\pi$  commute
avec le faisceau Hom interne si le premier module est de présentation finie, d'après Modules,
Lemme  \ref{modules-lemma-pullback-internal-hom}.
\end{proof}

\begin{proof}[Deuxième démonstration]
Démonstration de (1). D'après le Lemme  \ref{lemma-compare-constructions} , pour donner une application
 $\widetilde{M \otimes_R N}$  vers  $\widetilde M \otimes_{\mathcal{O}_X} \widetilde N$ , nous devons donner une application sur les
sections globales  $M \otimes_R N \to
\Gamma(X, \widetilde M \otimes_{\mathcal{O}_X} \widetilde N)$  qui existe par définition du produit tensoriel des
faisceaux de modules. Pour voir que cette application est un isomorphisme, il
suffit de vérifier qu'elle est un isomorphisme sur les fibres. Et cela découle de
la description des fibres de  $\widetilde{M}$  (soit dans le Lemme  \ref{lemma-spec-sheaves}  ou dans
Modules, Lemme  \ref{modules-lemma-construct-quasi-coherent-sheaves}), du fait que le produit tensoriel commute avec la
localisation (Algèbre, Lemme  \ref{algebra-lemma-tensor-product-localization}) et Modules, Lemme  \ref{modules-lemma-stalk-tensor-product}.

\medskip\noindent
Démonstration de (2). C'est similaire à la preuve de (1), en utilisant Algèbre, Lemme
 \ref{algebra-lemma-tensor-algebra-localization}  et Modules, Lemme  \ref{modules-lemma-stalk-tensor-algebra}.

\medskip\noindent
Démonstration de (3). Puisque la construction  $M \mapsto \widetilde{M}$  est fonctorielle, il
existe une application  $R$-linéaire  $\Hom_R(M, N) \to \Hom_{\mathcal{O}_X}(\widetilde{M}, \widetilde{N})$. La cible de cette
application est le module des sections globales de  $\SheafHom_{\mathcal{O}_X}(\widetilde M, \widetilde N)$. Ainsi, par le Lemme
 \ref{lemma-compare-constructions} , nous obtenons une application de  $\mathcal{O}_X$-modules  $\widetilde{\Hom_R(M,  N)} \to
\SheafHom_{\mathcal{O}_X}(\widetilde M, \widetilde N)$.
Nous vérifions que c'est un isomorphisme en comparant les fibres. Si  $M$ 
est de présentation finie en tant que  $R$-module, alors  $\widetilde M$  a
une présentation globale finie en tant que  $\mathcal{O}_X$-module. Ainsi, nous
concluons en utilisant Algèbre, Lemme  \ref{algebra-lemma-hom-from-finitely-presented}  et Modules, Lemme  \ref{modules-lemma-stalk-internal-hom}.
\end{proof}

\begin{proof}[ Troisième démonstration de la partie (1)]
Pour tout  $\mathcal{O}_X$-module  $\mathcal{F}$ , nous avons les isomorphismes suivants
fonctoriels en  $M$,  $N$, et  $\mathcal{F}$
\begin{align*}
\Hom_{\mathcal{O}_X}(\widetilde{M} \otimes _{\mathcal{O} _X} \widetilde{N},
\mathcal{F})
& =
\Hom_{\mathcal{O}_X}(\widetilde{M},
\SheafHom_{\mathcal{O} _X} (\widetilde{N}, \mathcal{F})) \\
& =
\Hom_R(M, \Gamma(X,
\SheafHom_{\mathcal{O}_X}(\widetilde{N}, \mathcal{F}))) \\
& =
\Hom_R(M, \Hom_{\mathcal{O}_X}(\widetilde{N}, \mathcal{F})) \\
& =
\Hom_R(M, \Hom_R(N, \Gamma(X,\mathcal{F}))) \\
& =
\Hom_R(M \otimes_R N, \Gamma(X, \mathcal{F})) \\
& =
\Hom_{\mathcal{O}_X}(\widetilde{M \otimes_R N}, \mathcal{F})
\end{align*}
La première égalité est Modules, Lemme  \ref{modules-lemma-internal-hom}. La deuxième égalité est la
propriété universelle de  $\widetilde{M}$, voir Lemme  \ref{lemma-compare-constructions}. La troisième
égalité vaut par définition de  $\SheafHom$. La quatrième égalité est la
propriété universelle de  $\widetilde{N}$. La cinquième égalité est Algèbre, Lemme
 \ref{algebra-lemma-hom-from-tensor-product}. L'égalité finale est la propriété universelle de  $\widetilde{M \otimes_R N}$. Par
le lemme de Yoneda (Catégories, Lemme  \ref{categories-lemma-yoneda}) nous obtenons (1).
\end{proof}

\begin{lemma}
\label{lemma-widetilde-pullback}
Soient  $(X, \mathcal{O}_X) = (\Spec(S), \mathcal{O}_{\Spec(S)})$,  $(Y, \mathcal{O}_Y) = (\Spec(R), \mathcal{O}_{\Spec(R)})$  des schémas affines. Soit  $\psi : (X, \mathcal{O}_X) \to (Y, \mathcal{O}_Y)$  un
morphisme de schémas affines, correspondant à l'homomorphisme d'anneaux  $\psi^\sharp : R \to S$ 
(voir Lemme  \ref{lemma-category-affine-schemes}).
\begin{enumerate}
\item Nous avons  $\psi^* \widetilde M = \widetilde{S \otimes_R M}$  fonctoriellement par rapport au  $R$-module  $M$.
\item Nous avons  $\psi_* \widetilde N = \widetilde{N_R}$  fonctoriellement par rapport au  $S$-module  $N$.
\end{enumerate}
\end{lemma}

\begin{proof}
La première assertion découle de l'identification dans le Lemme  \ref{lemma-compare-constructions}  et du
résultat des Modules, Lemme  \ref{modules-lemma-restrict-quasi-coherent}. La deuxième assertion découle du fait
que  $\psi^{-1}(D(f)) = D(\psi^\sharp(f))$  et donc
$$
\psi_* \widetilde N(D(f)) = \widetilde N(D(\psi^\sharp(f))) =
N_{\psi^\sharp(f)} = (N_R)_f = \widetilde{N_R}(D(f))
$$
comme voulu.
\end{proof}

\noindent
Le Lemme  \ref{lemma-widetilde-pullback}  ci-dessus dit en particulier que si vous restreignez le
faisceau  $\widetilde M$  à un sous-espace ouvert affine standard  $D(f)$, alors
vous obtenez  $\widetilde{M_f}$. Nous utiliserons ceci à partir de maintenant sans autre
mention.

\begin{lemma}
\label{lemma-quasi-coherent-affine}
Soit  $(X, \mathcal{O}_X) = (\Spec(R), \mathcal{O}_{\Spec(R)})$  un schéma affine. Soit  $\mathcal{F}$  un  $\mathcal{O}_X$-module
quasi-cohérent. Alors  $\mathcal{F}$  est isomorphe au faisceau associé au
 $R$-module  $\Gamma(X, \mathcal{F})$.
\end{lemma}

\begin{proof}
Soit  $\mathcal{F}$  un  $\mathcal{O}_X$-module quasi-cohérent. Comme chaque ouvert
standard  $D(f)$  est quasi-compact, nous voyons que  $X$  est
localement quasi-compact, c’est-à-dire que chaque point possède un système
fondamental de voisinages quasi-compacts, voir Topologie, Définition  \ref{topology-definition-locally-quasi-compact}.
Ainsi, d’après Modules, Lemme  \ref{modules-lemma-quasi-coherent-module} , pour chaque idéal premier  $\mathfrak p \subset R$ 
correspondant à  $x \in X$ , il existe un voisinage ouvert  $x \in U \subset X$  tel que
 $\mathcal{F}|_U$  soit isomorphe au faisceau quasi-cohérent associé à un certain
 $\mathcal{O}_X(U)$-module  $M$. En d’autres termes, nous obtenons un
recouvrement ouvert par des  $U$ possédant cette propriété. Par le Lemme
 \ref{lemma-standard-open} , par exemple, nous pouvons affiner ce recouvrement en un recouvrement
ouvert standard. Ainsi, nous obtenons un recouvrement  $\Spec(R) = \bigcup D(f_i)$  et des
 $R_{f_i}$-modules  $M_i$  et des isomorphismes  $\varphi_i : \mathcal{F}|_{D(f_i)} \to \mathcal{F}_{M_i}$  pour un
certain  $R_{f_i}$-module  $M_i$. Sur les intersections, nous obtenons des
isomorphismes
$$
\xymatrix{
\mathcal{F}_{M_i}|_{D(f_if_j)}
\ar[rr]^{\varphi_i^{-1}|_{D(f_if_j)}}
& &
\mathcal{F}|_{D(f_if_j)}
\ar[rr]^{\varphi_j|_{D(f_if_j)}}
& &
\mathcal{F}_{M_j}|_{D(f_if_j)}.
}
$$
Notons-les  $\psi_{ij}$. Il est clair que nous avons la condition de cocycle
$$
\psi_{jk}|_{D(f_if_jf_k)}
\circ
\psi_{ij}|_{D(f_if_jf_k)}
=
\psi_{ik}|_{D(f_if_jf_k)}
$$
sur les triples intersections.

\medskip\noindent
Rappelons que chacun des sous-espaces ouverts  $D(f_i)$,  $D(f_if_j)$,
 $D(f_if_jf_k)$  est un schéma affine. Par conséquent, les faisceaux  $\mathcal{F}_{M_i}$  sont
isomorphes aux faisceaux  $\widetilde M_i$  d'après le Lemme  \ref{lemma-compare-constructions}  ci-dessus. En
particulier, nous voyons que  $\mathcal{F}_{M_i}(D(f_if_j)) = (M_i)_{f_j}$, etc. De plus, d'après le Lemme
 \ref{lemma-compare-constructions}  ci-dessus, nous voyons que  $\psi_{ij}$  correspond à un unique
isomorphisme de  $R_{f_if_j}$-modules
$$
\psi_{ij} : (M_i)_{f_j} \longrightarrow (M_j)_{f_i}
$$
à savoir, l'effet de  $\psi_{ij}$  sur les sections au-dessus de  $D(f_if_j)$. De
plus, celles-ci satisfont alors la condition de cocycle selon laquelle
$$
\xymatrix{
(M_i)_{f_jf_k}
\ar[rd]_{\psi_{ij}}
\ar[rr]^{\psi_{ik}}
& &
(M_k)_{f_if_j} \\
&
(M_j)_{f_if_k} \ar[ru]_{\psi_{jk}}
}
$$
commute (pour tout triple  $i, j, k$).

\medskip\noindent
Algèbre, Lemme  \ref{algebra-lemma-glue-modules}  montre qu'il existe un
 $R$-module  $M$  et des identifications  $M_i = M_{f_i}$  compatibles avec les
morphismes  $\psi_{ij}$. Considérons  $\mathcal{F}_M = \widetilde M$. À ce stade, c'est une
formalité de montrer que  $\widetilde M$  est isomorphe au faisceau quasi-cohérent
 $\mathcal{F}$  avec lequel nous avons commencé. À savoir, les faisceaux  $\mathcal{F}$ 
et  $\widetilde M$  donnent lieu à des ensembles de données de recollement de
faisceaux de  $\mathcal{O}_X$-modules isomorphes par rapport au recouvrement
 $X = \bigcup D(f_i)$, voir Faisceaux, Section  \ref{sheaves-section-glueing-sheaves}  et en particulier le Lemme
 \ref{sheaves-lemma-mapping-property-glue}. Explicitement, dans la situation actuelle, cela se ramène à
l'argument suivant: Construisons un homomorphisme de  $R$-modules
$$
M \longrightarrow \Gamma(X, \mathcal{F}).
$$
À savoir, étant donné  $m \in M$ , nous obtenons  $m_i = m/1 \in M_{f_i} = M_i$  par la construction
de  $M$. Par la construction de  $M_i$ , cela correspond à une
section  $s_i \in \mathcal{F}(U_i)$. (À savoir,  $\varphi^{-1}_i(m_i)$.) Nous affirmons que  $s_i|_{D(f_if_j)} = s_j|_{D(f_if_j)}$.
Cela résulte de la construction de  $M$, qui donne  $\psi_{ij}(m_i) = m_j$,
et de celle des  $\psi_{ij}$. Par la condition de faisceau de
 $\mathcal{F}$ , cette collection de sections donne lieu à une section unique
 $s$  de  $\mathcal{F}$  sur  $X$. Nous laissons au lecteur le soin de
montrer que  $m \mapsto s$  est un morphisme de  $R$-modules. Par le Lemme
 \ref{lemma-compare-constructions} , nous obtenons un morphisme de  $\mathcal{O}_X$-modules associé
$$
\widetilde M \longrightarrow \mathcal{F}.
$$
Par construction, ce morphisme se réduit aux isomorphismes  $\varphi_i^{-1}$  sur
chaque  $D(f_i)$  et est donc un isomorphisme.
\end{proof}

\begin{lemma}
\label{lemma-equivalence-quasi-coherent}
Soit  $(X, \mathcal{O}_X) = (\Spec(R), \mathcal{O}_{\Spec(R)})$  un schéma affine. Les foncteurs  $M \mapsto \widetilde M$  et  $\mathcal{F} \mapsto \Gamma(X, \mathcal{F})$ 
définissent des équivalences quasi-inverses de catégories
$$
\xymatrix{
\QCoh(\mathcal{O}_X)
\ar@<1ex>[r]
&
\text{Mod}_R
\ar@<1ex>[l]
}
$$
entre la catégorie des  $\mathcal{O}_X$-modules quasi-cohérents et la catégorie des
 $R$-modules.
\end{lemma}

\begin{proof}
Voir les lemmes  \ref{lemma-compare-constructions}  et  \ref{lemma-quasi-coherent-affine}  ci-dessus.
\end{proof}

\noindent
Dorénavant, nous ne distinguerons pas entre les faisceaux quasi-cohérents sur les
schémas affines et les faisceaux de la forme  $\widetilde M$.

\begin{lemma}
\label{lemma-kernel-cokernel-quasi-coherent}
Soit  $X = \Spec(R)$  un schéma affine. Les noyaux et conoyaux des morphismes de
 $\mathcal{O}_X$-modules quasi-cohérents sont quasi-cohérents.
\end{lemma}

\begin{proof}
Cela découle de l'exactitude du foncteur  $\widetilde{\ }$  puisque d'après le lemme
 \ref{lemma-compare-constructions} , nous savons que tout morphisme  $\psi : \widetilde{M} \to \widetilde{N}$  provient d'un morphisme
de  $R$-modules  $\varphi : M \to N$. (Ainsi nous avons  $\Ker(\psi) = \widetilde{\Ker(\varphi)}$  et
 $\Coker(\psi) = \widetilde{\Coker(\varphi)}$.)
\end{proof}

\begin{lemma}
\label{lemma-colimit-quasi-coherent}
Soit  $X = \Spec(R)$  un schéma affine. La somme directe d'une collection quelconque
de faisceaux quasi-cohérents sur  $X$  est quasi-cohérente. Il en va de
même pour les colimites.
\end{lemma}

\begin{proof}
Supposons que les  $\mathcal{F}_i$,  $i \in I$,  forment une collection de faisceaux
quasi-cohérents sur  $X$. D'après le Lemme  \ref{lemma-equivalence-quasi-coherent}  ci-dessus, nous
pouvons écrire  $\mathcal{F}_i = \widetilde{M_i}$  pour un certain  $R$-module  $M_i$.
Posons  $M = \bigoplus M_i$. Considérons le faisceau  $\widetilde{M}$. Pour chaque ouvert
standard  $D(f)$ , nous avons
$$
\widetilde{M}(D(f)) = M_f =
\left(\bigoplus M_i\right)_f =
\bigoplus M_{i, f}.
$$
Ainsi nous voyons que le  $\mathcal{O}_X$-module quasi-cohérent  $\widetilde{M}$  est la
somme directe des faisceaux  $\mathcal{F}_i$. Un argument similaire fonctionne pour
les colimites générales.
\end{proof}

\begin{lemma}
\label{lemma-extension-quasi-coherent}
Soit  $(X, \mathcal{O}_X) = (\Spec(R), \mathcal{O}_{\Spec(R)})$  un schéma affine. Supposons que
$$
0 \to
\mathcal{F}_1 \to
\mathcal{F}_2 \to
\mathcal{F}_3 \to
0
$$
soit une suite exacte courte de faisceaux de  $\mathcal{O}_X$-modules. Si deux sur
trois sont quasi-cohérents, alors le troisième l'est aussi.
\end{lemma}

\begin{proof}
Cela est clair dans le cas où à la fois  $\mathcal{F}_1$  et  $\mathcal{F}_2$  sont
quasi-cohérents car le foncteur  $M \mapsto \widetilde M$  est exact, voir le Lemme
 \ref{lemma-spec-sheaves}. De même dans le cas où à la fois  $\mathcal{F}_2$  et  $\mathcal{F}_3$  sont
quasi-cohérents. Maintenant, supposons que  $\mathcal{F}_1 = \widetilde M_1$  et  $\mathcal{F}_3 = \widetilde M_3$  soient
quasi-cohérents. Posons  $M_2 = \Gamma(X, \mathcal{F}_2)$. Nous affirmons qu'il suffit de montrer que
la suite
$$
0 \to M_1 \to M_2 \to M_3 \to 0
$$
est exacte. En effet, si tel est le cas, alors (en utilisant la
propriété universelle du Lemme  \ref{lemma-compare-constructions}) nous obtenons un diagramme commutatif
$$
\xymatrix{
0 \ar[r] &
\widetilde M_1 \ar[r] \ar[d] &
\widetilde M_2 \ar[r] \ar[d] &
\widetilde M_3 \ar[r] \ar[d] &
0 \\
0 \ar[r] &
\mathcal{F}_1 \ar[r] &
\mathcal{F}_2 \ar[r] &
\mathcal{F}_3 \ar[r] &
0
}
$$
et le lemme du serpent permet alors de conclure.

\medskip\noindent
L'argument ``correct'' ici serait de montrer d'abord que  $H^1(X, \mathcal{F}) = 0$  pour tout
faisceau quasi-cohérent  $\mathcal{F}$. Cela n'est en fait pas si difficile, mais il
est peut-être préférable de le reporter à plus tard. À la place, nous utilisons
une petite astuce.

\medskip\noindent
Choisissez  $m \in M_3 = \Gamma(X, \mathcal{F}_3)$. Considérons l'ensemble suivant
$$
I = \{ f \in R \mid \text{l'élément }fm\text{ provient de }M_2\}.
$$
Clairement, il s'agit d'un idéal. Il suffit de montrer  $1 \in I$. Par
conséquent, il suffit de montrer que pour tout idéal premier  $\mathfrak p$ , il existe un
 $f \in I$,  $f \not\in \mathfrak p$. Soit  $x \in X$  le point correspondant à
 $\mathfrak p$. Comme la surjectivité peut être vérifiée sur les fibres, il existe
un voisinage ouvert  $U$  de  $x$  tel que  $m|_U$  provienne
d'une section locale  $s \in \mathcal{F}_2(U)$. En fait, nous pouvons supposer que  $U = D(f)$ 
est un ouvert standard, c’est-à-dire  $f \in R$,  $f \not \in \mathfrak p$. Nous allons
montrer que pour un certain  $N \gg 0$  nous avons  $f^N \in I$, ce qui terminera
la preuve.

\medskip\noindent
Prenons un point quelconque  $z \in V(f)$, disons correspondant à l'idéal premier
 $\mathfrak q \subset R$. Nous pouvons également trouver un  $g \in R$,  $g \not \in \mathfrak q$  tel que
 $m|_{D(g)}$  se relève en un certain  $s' \in \mathcal{F}_2(D(g))$. Considérons la différence
 $s|_{D(fg)} - s'|_{D(fg)}$. C'est un élément  $m'$  de  $\mathcal{F}_1(D(fg)) = (M_1)_{fg}$. Pour un certain
entier  $n = n(z)$ , l'élément  $f^n m'$  provient d'un certain  $m'_1 \in (M_1)_g$.
Nous voyons que  $f^n s$  s'étend à une section  $\sigma$  de  $\mathcal{F}_2$ 
sur  $D(f) \cup D(g)$  parce qu'elle coïncide avec la restriction de  $f^n s' + m'_1$  sur
 $D(f) \cap D(g) = D(fg)$. De plus, l'image de  $\sigma$  est la restriction de  $f^n m$  à
 $D(f) \cup D(g)$.

\medskip\noindent
Étant donné que  $V(f)$  est quasi-compact, il existe une liste finie
d'éléments  $g_1, \ldots, g_m \in R$  telle que  $V(f) \subset \bigcup D(g_j)$, un entier  $n > 0$  et des
sections  $\sigma_j \in \mathcal{F}_2(D(f) \cup D(g_j))$  telles que  $\sigma_j|_{D(f)} = f^n s$  et que chaque  $\sigma_j$  ait pour image la
section  $f^nm|_{D(f) \cup D(g_j)}$  de  $\mathcal{F}_3$. Considérons les différences
$$
\sigma_j|_{D(f) \cup D(g_jg_k)}
-
\sigma_k|_{D(f) \cup D(g_jg_k)}.
$$
Cela correspond à des sections de  $\mathcal{F}_1$  sur  $D(f) \cup D(g_jg_k)$  qui sont nulles
sur  $D(f)$. En particulier, leurs images dans  $\mathcal{F}_1(D(g_jg_k)) = (M_1)_{g_jg_k}$  sont nulles dans
 $(M_1)_{g_jg_kf}$. Ainsi, une certaine puissance élevée de  $f$  annule chacune
d'entre elles. En d'autres termes, les éléments  $f^N \sigma_j$, pour un certain
 $N \gg 0$ , satisfont à la condition de recollement du faisceau et
donnent lieu à une section  $\sigma $  de  $\mathcal{F}_2$  sur  $\bigcup (D(f) \cup D(g_j)) = X$  comme
souhaité.
\end{proof}







\section{Sous-espaces fermés des schémas affines}
\label{section-closed-in-affine}

\begin{example}
\label{example-closed-immersion-affines}
Soit  $R$  un anneau. Soit  $I \subset R$  un idéal. Considérons le morphisme
de schémas affines  $i : Z = \Spec(R/I) \to \Spec(R) = X$. Selon le Lemme d'Algèbre  \ref{algebra-lemma-spec-closed} , il s'agit
d'un homéomorphisme de  $Z$  sur un sous-ensemble fermé de  $X$. De
plus, si  $I \subset \mathfrak p \subset R$  est un idéal premier correspondant à un point  $x = i(z)$,
 $x \in X$,  $z \in Z$, alors sur les fibres, on obtient l'application
$$
\mathcal{O}_{X, x} = R_{\mathfrak p}
\longrightarrow
R_{\mathfrak p}/IR_{\mathfrak p} = \mathcal{O}_{Z, z}
$$
Ainsi, nous voyons que  $i$  est une immersion fermée d'espaces localement
annelés, voir la Définition  \ref{definition-closed-immersion-locally-ringed-spaces}. Clairement, cela est (isomorphe) à
le sous-espace fermé associé au faisceau quasi-cohérent d'idéaux  $\widetilde I$, comme
dans l'Exemple  \ref{example-closed-subspace}.
\end{example}

\begin{lemma}
\label{lemma-closed-immersion-affine-case}
\begin{slogan}
Pour les schémas affines, les immersions fermées correspondent aux idéaux.
\end{slogan}
Soit  $(X, \mathcal{O}_X) = (\Spec(R), \mathcal{O}_{\Spec(R)})$  un schéma affine. Soit  $i : Z \to X$  une quelconque immersion
fermée d'espaces localement annelés. Alors il existe un unique idéal  $I \subset R$ 
tel que le morphisme  $i : Z \to X$  puisse être identifié avec l'immersion fermée
 $\Spec(R/I) \to \Spec(R)$  construite dans l'Exemple  \ref{example-closed-immersion-affines}  ci-dessus.
\end{lemma}

\begin{proof}
La vérification est immédiate : d'après le Lemme  \ref{lemma-closed-immersion} , nous pouvons
identifier  $Z \to X$  avec le sous-espace fermé associé à un faisceau d'idéaux
 $\mathcal{I} \subset \mathcal{O}_X$  comme dans la Définition  \ref{definition-closed-subspace}  et l'Exemple  \ref{example-closed-subspace}.
Selon nos conventions, ce faisceau d'idéaux est localement engendré par des
sections en tant que faisceau de  $\mathcal{O}_X$-modules. Par conséquent, le faisceau
quotient  $\mathcal{O}_X / \mathcal{I}$  est localement sur  $X$  le conoyau d'une application
 $\bigoplus_{j \in J} \mathcal{O}_U \to \mathcal{O}_U$. Ainsi, par définition,  $\mathcal{O}_X / \mathcal{I}$  est quasi-cohérent. D'après nos
résultats de la Section  \ref{section-quasi-coherent-affine} , il est de la forme  $\widetilde S$  pour un
certain  $R$-module  $S$. De plus, comme  $\mathcal{O}_X = \widetilde R \to \widetilde S$  est
surjectif, nous voyons d'après le Lemme  \ref{lemma-extension-quasi-coherent}  que  $\mathcal{I}$ 
est également quasi-cohérent, disons  $\mathcal{I} = \widetilde I$. Bien sûr,  $I \subset R$  et  $S = R/I$  et
tout est clair.
\end{proof}













\section{Schémas}
\label{section-schemes}

\begin{definition}
\label{definition-scheme}
\begin{history}
Dans  \cite{EGA1}  ce que nous appelons un schéma était appelé un
``pré-sch\'éma'' et le nom ``sch\'éma'' était réservé pour ce
qui est un schéma séparé dans le projet Stacks. Dans la deuxième édition
 \cite{EGA1-second}  la terminologie a été remplacée par celle qui est désormais
standard. Cependant, on peut parfois rencontrer la terminologie ``préschéma'', par
exemple dans  \cite{Murre-lectures}.
\end{history}
Un  {\it  schéma}  est un espace localement annelé
avec la propriété que chaque point possède un voisinage ouvert qui est un schéma
affine. Un  {\it  morphisme de schémas}  est un
morphisme d'espaces localement annelés. La catégorie des schémas sera notée
 $\Sch$.
\end{definition}

\noindent
Soit $X$ un schéma. Nous utiliserons le (très léger) abus de langage
suivant. Nous dirons que $U \subset X$  est une  {\it partie ouverte
affine}, ou un  {\it ouvert affine}, si le
sous-espace ouvert $U$  est un schéma affine. Nous écrirons souvent
 $U = \Spec(R)$ pour indiquer que $U$  est isomorphe à  $\Spec(R)$ et de plus
que nous identifierons (temporairement) $U$  et  $\Spec(R)$.

\begin{lemma}
\label{lemma-open-subspace-scheme}
Soit  $X$  un schéma. Soit  $j : U \to X$  une immersion ouverte d'espaces
localement annelés. Alors  $U$  est un schéma. En particulier, tout
sous-espace ouvert de  $X$  est un schéma.
\end{lemma}

\begin{proof}
Soit  $U \subset X$. Soit  $u \in U$. Choisissez un voisinage ouvert affine
 $u \in V \subset X$. Comme les ouverts standards de  $V$  forment une base de la
topologie sur  $V$ , nous voyons qu'il existe un  $f\in \mathcal{O}_V(V)$  tel que
 $u \in D(f) \subset U$. Et  $D(f)$  est un schéma affine d'après le Lemme  \ref{lemma-standard-open-affine}.
Cela prouve que chaque point de  $U$  possède un voisinage ouvert qui est
affine.
\end{proof}

\noindent
De toute évidence, le lemme (ou sa preuve) montre que tout schéma  $X$  a
une base (voir Topologie, Section  \ref{topology-section-bases}) pour la topologie composée
d'ouverts affines.

\begin{example}
\label{example-not-affine}
Soit  $k$  un corps. Un exemple de schéma qui n'est pas affine est donné
par le sous-espace ouvert  $U = \Spec(k[x, y]) \setminus \{ (x, y)\}$  du schéma affine  $X =\Spec(k[x, y])$. Il est
recouvert par deux ouverts affines, à savoir  $D(x) = \Spec(k[x, y, 1/x])$  et  $D(y) = \Spec(k[x, y, 1/y])$  dont
l'intersection est  $D(xy) = \Spec(k[x, y, 1/xy])$. Par la propriété du faisceau pour  $\mathcal{O}_U$ ,
il existe une suite exacte
$$
0 \to
\Gamma(U, \mathcal{O}_U) \to
k[x, y, 1/x] \times k[x, y, 1/y] \to
k[x, y, 1/xy]
$$
Nous concluons que l'application  $k[x, y] \to \Gamma(U, \mathcal{O}_U)$  (provenant du morphisme
 $U \to X$) est un isomorphisme. Par conséquent,  $U$  ne peut pas être
affine puisqu'alors, d'après le Lemme  \ref{lemma-category-affine-schemes} , nous aurions  $U \cong X$.
\end{example}











\section{Immersions de schémas}
\label{section-immersions}

\noindent
Dans le Lemme  \ref{lemma-open-subspace-scheme} , nous avons vu que tout sous-espace ouvert d'un schéma
est un schéma. Ci-dessous, nous prouverons que la même chose vaut pour un
sous-espace fermé d'un schéma.

\medskip\noindent
Notez que la notion de faisceau quasi-cohérent de  $\mathcal{O}_X$-modules est définie
pour tout espace annelé  $X$ , en particulier lorsque  $X$  est un
schéma. Grâce à nos efforts dans la Section  \ref{section-quasi-coherent-affine} , nous savons qu'un tel
faisceau sur tout ouvert affine  $U \subset X$  est de la forme  $\widetilde M$  pour un
certain  $\mathcal{O}_X(U)$-module  $M$.

\begin{lemma}
\label{lemma-closed-subspace-scheme}
Soit  $X$  un schéma. Soit  $i : Z \to X$  une immersion fermée d'espaces
localement annelés.
\begin{enumerate}
\item L'espace localement annelé  $Z$  est un schéma,
\item Le noyau  $\mathcal{I}$  de l'application  $\mathcal{O}_X \to i_*\mathcal{O}_Z$  est un faisceau quasi-cohérent
d'idéaux,
\item pour tout ouvert affine  $U = \Spec(R)$  de  $X$  le morphisme  $i^{-1}(U) \to U$ 
peut être identifié à  $\Spec(R/I) \to \Spec(R)$  pour un certain idéal  $I \subset R$, et
\item nous avons  $\mathcal{I}|_U = \widetilde I$.
\end{enumerate}
En particulier, tout faisceau d'idéaux localement engendré par des sections est un
faisceau quasi-cohérent d'idéaux (et réciproquement), et tout sous-espace fermé de
 $X$  est un schéma.
\end{lemma}

\begin{proof}
Soit  $i : Z \to X$  une immersion fermée. Soit  $z \in Z$  un point. Choisissez un
voisinage ouvert affine  $i(z) \in U \subset X$ quelconque. Disons  $U = \Spec(R)$. D'après le
Lemme  \ref{lemma-closed-immersion-affine-case} , nous savons que  $i^{-1}(U) \to U$  peut être identifié au morphisme
de schémas affines  $\Spec(R/I) \to \Spec(R)$. Tout d'abord, cela implique que  $z \in i^{-1}(U) \subset Z$  est
un voisinage affine de  $z$. Ainsi,  $Z$  est un schéma.
Deuxièmement, cela implique que  $\mathcal{I}|_U$  est  $\widetilde I$. En d'autres termes,
pour chaque point  $x \in i(Z)$ , il existe un voisinage ouvert tel que  $\mathcal{I}$ 
soit quasi-cohérent dans ce voisinage. Notez que  $\mathcal{I}|_{X \setminus i(Z)}
\cong \mathcal{O}_{X \setminus i(Z)}$. Ainsi, la
restriction du faisceau d'idéaux est quasi-cohérente sur  $X \setminus i(Z)$  également.
Nous concluons que  $\mathcal{I}$  est quasi-cohérent.
\end{proof}

\begin{definition}
\label{definition-immersion}
Soit  $X$  un schéma.
\begin{enumerate}
\item Un morphisme de schémas est appelé une  {\it  immersion
ouverte}  s'il s'agit d'une immersion ouverte d'espaces localement
annelés (voir la Définition  \ref{definition-immersion-locally-ringed-spaces}).
\item Un  {\it  sous-schéma ouvert}  de  $X$  est un
sous-espace ouvert de  $X$  au sens de la Définition  \ref{definition-open-subspace}; un
sous-schéma ouvert de  $X$  est un schéma d'après le Lemme  \ref{lemma-open-subspace-scheme}.
\item Un morphisme de schémas est appelé une  {\it  immersion
fermée}  s'il s'agit d'une immersion fermée d'espaces localement annelés
(voir la Définition  \ref{definition-closed-immersion-locally-ringed-spaces}).
\item Un  {\it  sous-schéma fermé}  de  $X$  est un
sous-espace fermé de  $X$  au sens de la Définition  \ref{definition-closed-subspace}; un
sous-schéma fermé est un schéma d'après le Lemme  \ref{lemma-closed-subspace-scheme}.
\item Un morphisme de schémas $f : X \to Y$  est appelé une
 {\it immersion}, ou une
 {\it immersion localement fermée} s'il peut se
factoriser sous la forme $j \circ i$ où $i$  est une immersion fermée et
 $j$ est une immersion ouverte.
\end{enumerate}
\end{definition}

\noindent
Il résulte des lemmes des sections  \ref{section-open-immersion}  et  \ref{section-closed-immersion}  que toute
immersion ouverte (resp. \  fermée) de schémas est isomorphe à
l'inclusion d'un sous-schéma ouvert (resp. \  fermé) de la cible.

\medskip\noindent
Notre définition d'une immersion fermée se situe à mi-chemin entre Hartshorne et
l'EGA. Hartshorne définit une immersion fermée comme un morphisme $f : X \to Y$ de
schémas qui induit un homéomorphisme de $X$  sur un sous-ensemble fermé de
 $Y$ et tel que l'application $f^\# : \mathcal{O}_Y \to f_*\mathcal{O}_X$ soit surjective, voir \cite[Page 85]{H}. Nous montrerons
que cela est équivalent à notre notion dans le Lemme \ref{lemma-characterize-closed-immersions}. Dans
 \cite{EGA}, Grothendieck et Dieudonn\'e définissent d'abord les
sous-schémas fermés via la construction de l'Exemple \ref{example-closed-subspace} en utilisant des
faisceaux quasi-cohérents d'idéaux, puis définissent une immersion fermée comme un
morphisme $f : X \to Y$ qui induit un isomorphisme avec un sous-schéma fermé. Il
découle du Lemme  \ref{lemma-closed-subspace-scheme}  que cela correspond à notre notion.

\medskip\noindent
D'un point de vue pédagogique, la définition ci-dessus est un désastre/cauchemar.
Lors de l'enseignement de ce matériel aux étudiants, nous avons souvent trouvé
pratique de définir une immersion fermée comme un morphisme affine  $f : X \to Y$  de
schémas tel que  $f^\# : \mathcal{O}_Y \to f_*\mathcal{O}_X$  soit surjectif. En effet, il s’avère que la notion de
morphisme affine (Morphismes, Section  \ref{morphisms-section-affine}) est assez naturelle et facile
à comprendre.

\medskip\noindent
Pour plus d'informations sur les immersions fermées, nous suggérons au lecteur de
consulter Morphismes, Sections  \ref{morphisms-section-closed-immersions}  et  \ref{morphisms-section-closed-immersions-quasi-coherent}.

\medskip\noindent
Nous discuterons des sous-schémas localement fermés et des immersions à la fin de
cette section.

\begin{remark}
\label{remark-not-reverse-open-closed}
Si  $f : X \to Y$  est une immersion de schémas, il n'est généralement pas possible
de factoriser  $f$  en une immersion ouverte suivie d'une immersion fermée.
Voir Morphismes, Exemple  \ref{morphisms-example-thibaut}.
\end{remark}

\begin{lemma}
\label{lemma-immersion-when-closed}
Soit  $f : Y \to X$  une immersion de schémas. Alors  $f$  est une immersion
fermée si et seulement si  $f(Y) \subset X$  est un sous-ensemble fermé.
\end{lemma}

\begin{proof}
Si  $f$  est une immersion fermée, alors  $f(Y)$  est fermé par
définition. Réciproquement, supposons que  $f(Y)$  soit fermé. Par définition,
il existe un sous-schéma ouvert  $U \subset X$  tel que  $f$  soit la
composition d'une immersion fermée  $i : Y \to U$  et de l'immersion ouverte
 $j : U \to X$. Soit  $\mathcal{I} \subset \mathcal{O}_U$  le faisceau d'idéaux quasi-cohérent associé à
l'immersion fermée  $i$. Notez  $\mathcal{I}|_{U \setminus i(Y)}
= \mathcal{O}_{U \setminus i(Y)}
= \mathcal{O}_{X \setminus i(Y)}|_{U \setminus i(Y)}$. Ainsi, nous pouvons recoller
(voir Faisceaux, Section  \ref{sheaves-section-glueing-sheaves})  $\mathcal{I}$  et  $\mathcal{O}_{X \setminus i(Y)}$  en un faisceau
d'idéaux  $\mathcal{J} \subset \mathcal{O}_X$. Puisque chaque point de  $X$  possède un voisinage
où  $\mathcal{J}$  est quasi-cohérent, nous voyons que  $\mathcal{J}$  est
quasi-cohérent (en particulier localement engendré par des sections). Par
construction,  $\mathcal{O}_X/\mathcal{J}$  est supporté sur  $U$  et sa restriction à cet ouvert est égale à
 $\mathcal{O}_U/\mathcal{I}$. Ainsi, nous voyons que les sous-espaces fermés associés à
 $\mathcal{I}$  et  $\mathcal{J}$  sont canoniquement isomorphes, voir l'Exemple
 \ref{example-closed-subspace}. En particulier, le sous-espace fermé de  $U$  associé à
 $\mathcal{I}$  est isomorphe à un sous-espace fermé de  $X$. Puisque
 $Y \to U$  est identifié avec le sous-espace fermé associé à  $\mathcal{I}$, voir
le Lemme  \ref{lemma-closed-immersion}, nous concluons que  $Y \to U \to X$  est une immersion fermée.
\end{proof}

\noindent
Soit  $f : Y \to X$  une immersion. Soit  $Z = \overline{f(Y)} \setminus f(Y)$  qui est un sous-ensemble fermé
de  $X$. Soit  $U = X \setminus Z$. Le lemme implique que  $U$  est le plus
grand sous-espace ouvert de  $X$  tel que  $f : Y \to X$  se factorise par une
immersion fermée dans  $U$. Nous définissons un  {\it 
sous-schéma localement fermé de  $X$}  comme une paire  $(Z, U)$ 
constituée d'un sous-schéma fermé  $Z$  d'un sous-schéma ouvert
 $U$  de  $X$  tel qu'en outre  $\overline{Z} \cup U = X$. Nous disons
généralement simplement « soit  $Z$  un sous-schéma localement fermé de
 $X$ » puisque nous pouvons récupérer  $U$  à partir du morphisme
 $Z \to X$. Ce qui précède montre alors que toute immersion  $f : Y \to X$  se
factorise de manière unique en  $Y \to Z \to X$  où  $Z$  est un sous-schéma
localement fermé de  $X$  et  $Y \to Z$  est un isomorphisme.

\medskip\noindent
L'intérêt de ceci est que la collection de sous-schémas localement fermés de
 $X$  forme un ensemble. Nous pouvons définir un ordre partiel sur cet
ensemble, que nous appelons inclusion pour des raisons évidentes. Pour être
précis, si  $Z \to X$  et  $Z' \to X$  sont deux sous-schémas localement fermés
de  $X$, alors nous disons que  {\it   $Z$  est
contenu dans  $Z'$}  simplement si le morphisme  $Z \to X$  se
factorise à travers  $Z'$. Si c’est le cas, alors bien sûr  $Z$  est
identifié à un sous-schéma localement fermé unique de  $Z'$, et ainsi de
suite.

\medskip\noindent
Pour plus d’informations sur les immersions, nous renvoyons le lecteur à
Morphismes, Section  \ref{morphisms-section-immersions}.










\section{Topologie de Zariski des schémas}
\label{section-topology}

\noindent
Voir Topologie, Section  \ref{topology-section-introduction}  pour quelques notions de base en topologie
adaptées à la topologie de Zariski des schémas.

\begin{lemma}
\label{lemma-scheme-sober}
Soit  $X$  un schéma. Tout sous-ensemble fermé irréductible de  $X$ 
a un point générique unique. En d'autres termes,  $X$  est un espace
topologique sobre, voir Topologie, Définition  \ref{topology-definition-generic-point}.
\end{lemma}

\begin{proof}
Soit  $Z \subset X$  un sous-ensemble fermé irréductible. Pour chaque ouvert affine
 $U \subset X$,  $U = \Spec(R)$ , nous savons que  $Z \cap U = V(I)$  pour un idéal radical
unique  $I \subset R$. Notez que  $Z \cap U$  est soit vide, soit irréductible. Dans
le second cas (qui se produit pour au moins un  $U$), nous voyons que
 $I = \mathfrak p$  est un idéal premier, qui est un point générique  $\xi$  de
 $Z \cap U$. Il s'ensuit que  $Z = \overline{\{\xi\}}$, en d'autres termes  $\xi$  est un
point générique de  $Z$. Si  $\xi'$  était un deuxième point
générique, alors  $\xi' \in Z \cap U$  et il en découle immédiatement que  $\xi' = \xi$.
\end{proof}

\begin{lemma}
\label{lemma-basis-affine-opens}
Soit  $X$  un schéma. La collection des ouverts affines de  $X$ 
forme une base pour la topologie sur  $X$.
\end{lemma}

\begin{proof}
Cela découle de la discussion sur les sous-schémas ouverts dans la Section
 \ref{section-schemes}.
\end{proof}

\begin{remark}
\label{remark-intersection-affine-opens}
En général, l'intersection de deux ouverts affines dans  $X$  n'est pas un
ouvert affine. Voir l'Exemple  \ref{example-affine-space-zero-doubled}.
\end{remark}

\begin{lemma}
\label{lemma-locally-quasi-compact}
L'espace topologique sous-jacent de tout schéma est localement quasi-compact, voir
Topologie, Définition  \ref{topology-definition-locally-quasi-compact}.
\end{lemma}

\begin{proof}
Cela découle du Lemme  \ref{lemma-basis-affine-opens}  ci-dessus et du fait que le spectre d'un anneau
est quasi-compact, voir Algèbre, Lemme  \ref{algebra-lemma-quasi-compact}.
\end{proof}

\begin{lemma}
\label{lemma-standard-open-two-affines}
Soit  $X$  un schéma. Soient  $U, V$  des ouverts affines de
 $X$, et soit  $x \in U \cap V$. Il existe un voisinage affine ouvert
 $W$  de  $x$  tel que  $W$  soit un ouvert standard à la
fois de  $U$  et de  $V$.
\end{lemma}

\begin{proof}
Posons  $U = \Spec(A)$  et  $V = \Spec(B)$. Supposons que  $x$  corresponde aux idéaux premiers
 $\mathfrak p \subset A$  et  $\mathfrak q \subset B$. Nous pouvons choisir un  $f \in A$,
 $f \not \in \mathfrak p$  tel que  $D(f) \subset U \cap V$. Remarquons que tout ouvert standard de
 $D(f)$  est un ouvert standard de  $\Spec(A) = U$. Ainsi, nous pouvons
supposer que  $U \subset V$. En d'autres termes, maintenant nous pouvons considérer
 $U$  comme un ouvert affine de  $V$. Ensuite, nous choisissons
un  $g \in B$,  $g \not \in \mathfrak q$  tel que  $D(g) \subset U$. Dans ce cas, nous voyons que
 $D(g) = D(g_A)$  où  $g_A \in A$  désigne l'image de  $g$  par l'application
 $B \to A$. Ainsi le lemme est prouvé.
\end{proof}

\begin{lemma}
\label{lemma-good-subcover}
Soit  $X$  un schéma. Soit  $X = \bigcup_i U_i$  un recouvrement ouvert affine. Soit
 $V \subset X$  un ouvert affine. Il existe un recouvrement ouvert standard
 $V = \bigcup_{j = 1, \ldots, m} V_j$  (voir Définition  \ref{definition-standard-covering}) tel que chaque  $V_j$  soit
un ouvert standard de l'un des  $U_i$.
\end{lemma}

\begin{proof}
Choisissez  $v \in V$. Alors  $v \in U_i$  pour un certain  $i$. D'après le
Lemme  \ref{lemma-standard-open-two-affines}  ci-dessus, il existe un ouvert  $v \in W_v \subset V \cap U_i$  tel que
 $W_v$  soit un ouvert standard à la fois dans  $V$  et  $U_i$.
Comme  $V$  est quasi-compact, le lemme s'ensuit.
\end{proof}

\begin{lemma}
\label{lemma-sheaf-on-affines}
Soit  $X$  un schéma. Soit  $\mathcal{B}$  l'ensemble des ouverts affines de
 $X$. Soit  $\mathcal{F}$  un préfaisceau d'ensembles sur  $\mathcal{B}$, voir
Faisceaux, Définition  \ref{sheaves-definition-presheaf-basis}. Les énoncés suivants sont équivalents
\begin{enumerate}
\item $\mathcal{F}$  est la restriction d'un faisceau sur  $X$  à  $\mathcal{B}$,
\item $\mathcal{F}$  est un faisceau sur  $\mathcal{B}$, et
\item $\mathcal{F}(\emptyset)$  est un singleton et chaque fois que  $U = V \cup W$  avec  $U, V, W \in \mathcal{B}$  et
 $V, W \subset U$  ouverts standards (Algèbre, Définition  \ref{algebra-definition-Zariski-topology}) l'application
$$
\mathcal{F}(U) \longrightarrow \mathcal{F}(V) \times \mathcal{F}(W)
$$
est injective avec pour image l'ensemble des paires  $(s, t)$  telles que
 $s|_{V \cap W} = t|_{V \cap W}$.
\end{enumerate}
\end{lemma}

\begin{proof}
L'équivalence de (1) et (2) est Faisceaux, Lemme  \ref{sheaves-lemma-restrict-basis-equivalence}. Il est clair
que (2) implique (3). Par conséquent, il suffit de prouver que (3) implique (2).
D'après Faisceaux, Lemme  \ref{sheaves-lemma-cofinal-systems-coverings-standard-case} , et le Lemme  \ref{lemma-standard-open}  ci-dessus, il suffit de
prouver que la condition de faisceau est satisfaite pour les recouvrements ouverts
standards (Définition  \ref{definition-standard-covering}) des éléments de  $\mathcal{B}$. Soit
 $U = U_1 \cup \ldots \cup U_n$  un recouvrement ouvert standard avec  $U \subset X$  ouvert affine. Nous
prouverons la condition de faisceau pour ce recouvrement par récurrence sur
 $n$. Si  $n = 0$, alors  $U$  est vide et nous obtenons la
condition de faisceau par hypothèse. Si  $n = 1$, alors il n'y a rien à
prouver. Si  $n = 2$, alors il s'agit de l'hypothèse (3). Si  $n > 2$,
alors nous écrivons  $U_i = D(f_i)$  pour  $f_i \in A = \mathcal{O}_X(U)$. Supposons que les  $s_i \in \mathcal{F}(U_i)$ 
soient des sections telles que  $s_i|_{U_i \cap U_j} = s_j|_{U_i \cap U_j}$  pour tout  $1 \leq i < j \leq n$. Comme
 $U = U_1 \cup \ldots \cup U_n$ , nous avons  $1 = \sum_{i = 1, \ldots, n} a_i f_i$  dans  $A$  pour certains  $a_i \in A$,
voir Algèbre, Lemme  \ref{algebra-lemma-Zariski-topology}. Posons  $g = \sum_{i = 1, \ldots, n - 1} a_if_i$. Alors  $U = D(g) \cup D(f_n)$.
Remarquons que  $D(g) = D(gf_1) \cup \ldots \cup D(gf_{n - 1})$  est un recouvrement ouvert standard. Par récurrence, il
existe une section unique  $s' \in \mathcal{F}(D(g))$  qui concorde avec  $s_i|_{D(gfi)}$  pour
 $i = 1, \ldots, n - 1$. Nous affirmons que  $s'$  et  $s_n$  ont la même
restriction à  $D(gf_n)$. Cela est vrai par récurrence, en utilisant le recouvrement
 $D(gf_n) = D(gf_nf_1) \cup \ldots \cup D(gf_nf_{n - 1})$. Ainsi, il existe une section unique  $s \in \mathcal{F}(U)$  dont la restriction
à  $D(g)$  est  $s'$  et dont la restriction à  $D(f_n)$  est
 $s_n$. Nous omettons la vérification que  $s$  se restreint à
 $s_i$  sur  $D(f_i)$  pour  $i = 1, \ldots, n - 1$  et nous omettons la vérification
que  $s$  est unique.
\end{proof}

\begin{lemma}
\label{lemma-scheme-finite-discrete-affine}
Soit  $X$  un schéma dont l'espace topologique sous-jacent est un ensemble
discret fini. Alors  $X$  est affine.
\end{lemma}

\begin{proof}
Prenons  $X = \{x_1, \ldots, x_n\}$. Alors  $U_i = \{x_i\}$  est un voisinage ouvert de  $x_i$.
Par le Lemme  \ref{lemma-basis-affine-opens}  il est affine. Par conséquent,  $X$  est une
union disjointe finie de schémas affines, et est donc affine par le Lemme
 \ref{lemma-disjoint-union-affines}.
\end{proof}

\begin{example}
\label{example-scheme-without-closed-points}
Il existe un schéma sans points fermés. À savoir, soit  $R$  un anneau
local intègre dont le spectre ressemble à  $(0) = \mathfrak p_0 \subset \mathfrak p_1 \subset \mathfrak p_2
\subset \ldots \subset \mathfrak m$. Alors le sous-schéma ouvert
 $\Spec(R) \setminus \{\mathfrak m\}$  n'a pas de point fermé. Pour voir qu'un tel anneau  $R$ 
existe, nous utilisons le fait que pour tout groupe totalement ordonné
 $(\Gamma, \geq)$ , il existe un anneau de valuation  $A$  dont le groupe de
valuation est  $(\Gamma, \geq)$, voir  \cite{Krull}. Voir Algèbre, Section  \ref{algebra-section-valuation-rings}  pour
la notation. Nous prenons  $\Gamma = \mathbf{Z}x_1 \oplus \mathbf{Z}x_2 \oplus \mathbf{Z}x_3 \oplus
\ldots$  et nous définissons  $\sum_i a_i x_i \geq 0$  si et
seulement si le premier  $a_i$  non nul est  $> 0$, ou si tous les coefficients satisfont
 $a_i = 0$. Donc  $x_1 \geq x_2 \geq x_3 \geq  \ldots \geq 0$. Les sous-ensembles  $x_i + \Gamma_{\geq 0}$  sont des idéaux
premiers de  $(\Gamma, \geq)$, voir l'algèbre, notation au-dessus du Lemme
 \ref{algebra-lemma-ideals-valuation-ring}. Ceux-ci, ainsi que  $\emptyset$  et  $\Gamma_{\geq 0}$ , sont les seuls
idéaux premiers. Par conséquent,  $A$  est un exemple d'un anneau avec la
structure donnée de son spectre, selon Algèbre, Lemme  \ref{algebra-lemma-ideals-valuation-ring}.
\end{example}



















\section{Schémas réduits}
\label{section-reduced}

\begin{definition}
\label{definition-reduced}
Soit  $X$  un schéma. Nous disons que  $X$  est
 {\it  réduit}  si chaque anneau local  $\mathcal{O}_{X, x}$ 
est réduit.
\end{definition}

\begin{lemma}
\label{lemma-reduced}
Un schéma  $X$  est réduit si et seulement si  $\mathcal{O}_X(U)$  est un anneau
réduit pour tout  $U \subset X$  ouvert.
\end{lemma}

\begin{proof}
Supposons que  $X$  soit réduit. Soit  $f \in \mathcal{O}_X(U)$  une section telle que
 $f^n = 0$. Alors l'image de  $f$  dans  $\mathcal{O}_{U, u}$  est nulle pour tout
 $u \in U$. Par conséquent,  $f$  est nul, voir Faisceaux, Lemme
 \ref{sheaves-lemma-sheaf-subset-stalks}. Inversement, supposons que  $\mathcal{O}_X(U)$  soit réduit pour tous les
ouverts  $U$. Prenez un élément non nul  $f \in \mathcal{O}_{X, x}$. Tout représentant
 $(U, f \in \mathcal{O}(U))$  de  $f$  est non nul et donc non nilpotent. Par conséquent,
 $f$  n'est pas nilpotent dans  $\mathcal{O}_{X, x}$.
\end{proof}

\begin{lemma}
\label{lemma-affine-reduced}
Un schéma affine  $\Spec(R)$  est réduit si et seulement si  $R$  est
réduit.
\end{lemma}

\begin{proof}
L'implication directe suit immédiatement du Lemme  \ref{lemma-reduced}  ci-dessus. Dans
l'autre sens, cela découle du fait que toute localisation d'un anneau réduit est
réduite, et en particulier que les anneaux locaux d'un anneau réduit sont réduits.
\end{proof}

\begin{lemma}
\label{lemma-reduced-closed-subscheme}
Soit  $X$  un schéma. Soit  $T \subset X$  un sous-ensemble fermé. Il existe
un sous-schéma fermé unique  $Z \subset X$  possédant les propriétés suivantes: (a)
l'espace topologique sous-jacent de  $Z$  est égal à  $T$, et (b)
 $Z$  est réduit.
\end{lemma}

\begin{proof}
Soit  $\mathcal{I} \subset \mathcal{O}_X$  le sous-préfaisceau défini par la règle
$$
\mathcal{I}(U) = \{f \in \mathcal{O}_X(U) \mid
f(t) = 0\text{ pour tout }t \in T\cap U\}
$$
Ici, nous utilisons  $f(t)$  pour indiquer l'image de  $f$  dans le
corps résiduel  $\kappa(t)$  de  $X$  en  $t$. En raison de la
nature locale de la condition, il est clair que  $\mathcal{I}$  est un faisceau
d'idéaux. De plus, soit  $U = \Spec(R)$  un ouvert affine. Nous pouvons écrire
 $T \cap U = V(I)$  pour un idéal radical unique  $I \subset R$. Étant donné un
 $\mathfrak p \in V(I)$  premier correspondant à  $t \in T \cap U$  et un élément  $f \in R$ , nous
avons  $f(t) = 0 \Leftrightarrow f \in \mathfrak p$. Ainsi  $\mathcal{I}(U) = \bigcap_{\mathfrak p \in V(I)} \mathfrak p
= I$  par Algèbre, Lemme  \ref{algebra-lemma-Zariski-topology}. De plus,
pour tout ouvert standard  $D(g) \subset \Spec(R) = U$ , nous avons  $\mathcal{I}(D(g)) = I_g$  par le même
raisonnement. Ainsi  $\widetilde I$  et  $\mathcal{I}|_U$  coïncident (en tant qu'idéaux)
sur une base d'ouverts et sont donc égaux. Par conséquent,  $\mathcal{I}$  est un
faisceau quasi-cohérent d'idéaux.

\medskip\noindent
À ce stade, nous pouvons définir  $Z$  comme le sous-espace fermé associé
au faisceau d'idéaux  $\mathcal{I}$. Pour chaque ouvert affine  $U = \Spec(R)$  de
 $X$ , nous voyons que  $Z \cap U = \Spec(R/I)$  où  $I$  est un idéal radical et
donc  $Z$  est réduit (d'après le Lemme  \ref{lemma-affine-reduced}  ci-dessus). Par
construction, le sous-ensemble fermé sous-jacent de  $Z$  est  $T$.
Par conséquent, nous avons trouvé un sous-schéma fermé avec les propriétés (a) et
(b).

\medskip\noindent
Soit  $Z' \subset X$  un second sous-schéma fermé avec les propriétés (a) et (b). Pour
chaque ouvert affine  $U = \Spec(R)$  de  $X$ , nous voyons que  $Z' \cap U = \Spec(R/I')$ 
pour un certain idéal  $I' \subset R$. D'après le Lemme  \ref{lemma-affine-reduced} , l'anneau
 $R/I'$  est réduit et donc  $I'$  est radical. Puisque  $V(I') = T \cap U = V(I)$ ,
nous en déduisons que  $I = I'$  par le Lemme d'Algèbre  \ref{algebra-lemma-Zariski-topology}. Par
conséquent,  $Z'$  et  $Z$  sont définis par le même faisceau
d'idéaux et donc sont égaux.
\end{proof}

\begin{definition}
\label{definition-reduced-induced-scheme}
Soit $X$ un schéma. Soit $Z \subset X$ un sous-ensemble fermé.
Une {\it structure de schéma sur $Z$}  est donnée
par un sous-schéma fermé  $Z'$  de  $X$ dont l'ensemble sous-jacent
est égal à $Z$. Nous disons souvent « soit  $(Z, \mathcal{O}_Z)$ une
structure de schéma sur $Z$ » pour indiquer cela. La
 {\it structure de schéma réduite
induite} sur $Z$  est celle construite dans le Lemme  \ref{lemma-reduced-closed-subscheme}.
La  {\it réduction $X_{red}$  de  $X$} est la
structure de schéma réduite induite sur $X$  lui-même.
\end{definition}

\noindent
Souvent, lorsque nous disons « soit  $Z \subset X$  une composante irréductible de
 $X$ », nous pensons à  $Z$  comme à un sous-schéma fermé réduit de
 $X$  en utilisant la structure de schéma réduite induite.

\begin{remark}
\label{remark-reduced-induced-locally-closed}
Soit  $X$  un schéma. Soit  $T \subset X$  un sous-ensemble localement fermé.
Dans cette situation, nous utilisons parfois aussi l'expression « structure de
schéma réduite induite sur  $T$ ». Cela se réfère à la structure de schéma
réduite induite de la Définition  \ref{definition-reduced-induced-scheme}  lorsque nous considérons  $T$ 
comme un sous-ensemble fermé du sous-schéma ouvert  $X \setminus \partial T$  de  $X$.
Ici,  $\partial T = \overline{T} \setminus T$  est la « frontière » de  $T$  dans l'espace topologique
de  $X$.
\end{remark}

\begin{lemma}
\label{lemma-map-into-reduction}
Soit  $X$  un schéma. Soit  $Z \subset X$  un sous-schéma fermé. Soit
 $Y$  un schéma réduit. Un morphisme  $f : Y \to X$  se factorise à travers
 $Z$  si et seulement si  $f(Y) \subset Z$  (au sens des ensembles). En
particulier, tout morphisme  $Y \to X$  se factorise comme  $Y \to X_{red} \to X$.
\end{lemma}

\begin{proof}
Supposons  $f(Y) \subset Z$  (au sens des ensembles). Soit  $\mathcal{I} \subset \mathcal{O}_X$  le faisceau
d'idéaux de  $Z$. Pour tous ouverts affines  $U \subset X$  et  $\Spec(B) = V \subset Y$  tels que
 $f(V) \subset U$ , et pour tout  $g \in \mathcal{I}(U)$ , l'image inverse  $b = f^\sharp(g) \in \Gamma(V, \mathcal{O}_Y) = B$  a une image
nulle dans le corps résiduel de tout  $y \in V$. En d'autres termes  $b \in \bigcap_{\mathfrak p \subset B} \mathfrak p$.
Cela implique  $b = 0$  puisque  $B$  est réduit (lemme  \ref{lemma-reduced} et
Algèbre, lemme  \ref{algebra-lemma-Zariski-topology}). Ainsi  $f$  se factorise à travers
 $Z$  selon le lemme  \ref{lemma-characterize-closed-subspace}.
\end{proof}










\section{Points des schémas}
\label{section-points}

\noindent
Étant donné un schéma  $X$ , nous pouvons définir un foncteur
$$
h_X : \Sch^{opp}
\longrightarrow
\textit{Ensembles}, \quad
T \longmapsto \Mor(T, X).
$$
Voir Catégories, Exemple  \ref{categories-example-hom-functor}. On l'appelle le foncteur des points
 {\it  de  $X$}. Une partie amusante de la
théorie des schémas est de trouver des descriptions de la géométrie interne de
 $X$  en termes de ce foncteur  $h_X$. Dans cette section, nous
trouvons une manière simple de décrire les points de  $X$.

\medskip\noindent
Soit  $X$  un schéma. Soit  $R$  un anneau local avec un idéal
maximal  $\mathfrak m \subset R$. Supposons que  $f : \Spec(R) \to X$  soit un morphisme de schémas.
Soit  $x \in X$  l'image du point fermé  $\mathfrak m \in \Spec(R)$. Alors nous obtenons un
homomorphisme local d'anneaux
$$
f^\sharp :
\mathcal{O}_{X, x}
\longrightarrow
\mathcal{O}_{\Spec(R), \mathfrak m} = R.
$$

\begin{lemma}
\label{lemma-morphism-from-spec-local-ring}
Soit  $X$  un schéma. Soit  $R$  un anneau local. La construction
ci-dessus donne une correspondance bijective entre les morphismes  $\Spec(R) \to X$  et
les paires  $(x, \varphi)$  constituées d'un point  $x \in X$  et d'un homomorphisme
local d'anneaux  $\varphi : \mathcal{O}_{X, x} \to R$.
\end{lemma}

\begin{proof}
Soit  $A$  un anneau. Pour tout homomorphisme d'anneaux  $\psi : A \to R$ , il
existe un idéal premier unique  $\mathfrak p \subset A$  et une factorisation  $A \to A_{\mathfrak p} \to R$  où
la dernière application est un homomorphisme local d'anneaux. À savoir,
 $\mathfrak p = \psi^{-1}(\mathfrak m)$. Via le Lemme  \ref{lemma-morphism-into-affine} , cela prouve que le lemme est vrai si
 $X$  est un schéma affine.

\medskip\noindent
Soit  $X$  un schéma général. Tout  $x \in X$  est contenu dans un ouvert
affine  $U \subset X$. Par le cas affine, nous concluons que chaque paire
 $(x, \varphi)$  résulte de la construction qui précède le lemme.

\medskip\noindent
Pour terminer la preuve, il suffit de montrer que tout morphisme  $f : \Spec(R) \to X$  a
son image contenue dans tout ouvert affine contenant l'image  $x$  du point
fermé de  $\Spec(R)$. En fait, soit  $x \in V \subset X$  un voisinage ouvert quelconque
contenant  $x$. Alors  $f^{-1}(V) \subset \Spec(R)$  est un ouvert contenant le point fermé
unique et donc égal à  $\Spec(R)$.
\end{proof}

\noindent
Comme cas particulier du lemme ci-dessus, nous obtenons pour tout point
 $x$  d'un schéma  $X$  un morphisme canonique
\begin{equation}
\label{equation-canonical-morphism}
\Spec(\mathcal{O}_{X, x}) \longrightarrow X
\end{equation}
correspondant à l'application identité sur l'anneau local de  $X$  en
 $x$. On peut reformuler le lemme ci-dessus en disant que pour tout
morphisme  $f : \Spec(R) \to X$ , il existe un point unique  $x \in X$  tel que
 $f$  se factorise comme  $\Spec(R) \to \Spec(\mathcal{O}_{X, x}) \to X$  où la première application provient
d'un homomorphisme local  $\mathcal{O}_{X, x} \to R$.

\medskip\noindent
Dans le cas où nous avons un morphisme de schémas  $f : X \to S$, et un point
 $x$  dont l'image est un point  $s \in S$ , nous obtenons un diagramme
commutatif
$$
\xymatrix{
\Spec(\mathcal{O}_{X, x}) \ar[r] \ar[d] & X \ar[d] \\
\Spec(\mathcal{O}_{S, s}) \ar[r] & S
}
$$
où l'application verticale gauche correspond à l'homomorphisme local
 $f^\sharp_x : \mathcal{O}_{S, s} \to \mathcal{O}_{X, x}$.

\begin{lemma}
\label{lemma-specialize-points}
Soit  $X$  un schéma. Soient  $x, x' \in X$  des points de  $X$. Alors
 $x' \in X$  est une générisation de  $x$  si et seulement si
 $x'$  est dans l'image du morphisme canonique  $\Spec(\mathcal{O}_{X, x}) \to X$.
\end{lemma}

\begin{proof}
Une application continue préserve la relation de spécialisation/générisation.
Puisque chaque point de  $\Spec(\mathcal{O}_{X, x})$  est une générisation du point fermé, nous
voyons que chaque point dans l'image de  $\Spec(\mathcal{O}_{X, x}) \to X$  est une générisation de
 $x$. Inversement, supposons que  $x'$  soit une générisation de
 $x$. Choisissez un voisinage ouvert affine  $U = \Spec(R)$  de  $x$.
Alors  $x' \in U$. Disons que  $\mathfrak p \subset R$  et  $\mathfrak p' \subset R$  sont les idéaux
premiers correspondant à  $x$  et  $x'$. Puisque  $x'$  est
une générisation de  $x$ , nous voyons que  $\mathfrak p' \subset \mathfrak p$. Cela signifie
que  $\mathfrak p'$  est dans l'image du morphisme  $\Spec(\mathcal{O}_{X, x}) = \Spec(R_{\mathfrak p})
\to \Spec(R) = U \subset X$  comme voulu.
\end{proof}

\noindent
Maintenant, discutons des morphismes provenant des spectres des corps. Soit
 $(R, \mathfrak m, \kappa)$  un anneau local avec un idéal maximal  $\mathfrak m$  et un corps
résiduel  $\kappa$. Soit  $K$  un corps. Un homomorphisme local
 $R \to K$ , par définition, se factorise comme  $R \to \kappa \to K$, c'est-à-dire qu'il
revient à un morphisme  $\kappa \to K$. Ainsi, nous voyons que les
morphismes
$$
\Spec(K) \longrightarrow X
$$
correspondent à des paires  $(x, \kappa(x) \to K)$. Nous pouvons définir un préordre sur les
morphismes des spectres des corps vers  $X$  en disant que  $\Spec(K) \to X$ 
domine  $\Spec(L) \to X$  si  $\Spec(K) \to X$  se factorise par  $\Spec(L) \to X$. Cela suggère la
notion suivante: Disons temporairement que deux morphismes  $p : \Spec(K) \to X$  et
 $q : \Spec(L) \to X$  sont  {\it  équivalents }  s'il existe un
troisième corps  $\Omega$  et un diagramme commutatif
$$
\xymatrix{
\Spec(\Omega) \ar[r] \ar[d] &
\Spec(L) \ar[d]^q \\
\Spec(K) \ar[r]^p &
X
}
$$
Bien sûr, cela implique immédiatement que les points uniques des trois schémas
 $\Spec(K)$,  $\Spec(L)$ et  $\Spec(\Omega)$  se projettent sur le même  $x \in X$.
Ainsi, un diagramme (d'après les remarques ci-dessus) correspond à un point
 $x \in X$  et à un diagramme commutatif
$$
\xymatrix{
\Omega &
L \ar[l] \\
K \ar[u] &
\kappa(x) \ar[l] \ar[u]
}
$$
de corps. Cela définit une relation d'équivalence, car étant donné un ensemble
d'extensions de corps  $K_i/\kappa$ , il existe une extension de corps  $\Omega/\kappa$ 
telle que toutes les extensions de corps  $K_i$  soient contenues dans
l'extension  $\Omega$.

\begin{lemma}
\label{lemma-characterize-points}
Soit  $X$  un schéma. Les points de  $X$  correspondent
bijectivement aux classes d'équivalence de morphismes des spectres de corps dans
 $X$. De plus, chaque classe d'équivalence contient un élément minimal
(unique à isomorphisme unique près)  $\Spec(\kappa(x)) \to X$.
\end{lemma}

\begin{proof}
Cela découle de la discussion ci-dessus.
\end{proof}

\noindent
Bien sûr, les morphismes  $\Spec(\kappa(x)) \to X$  se factorisent par les morphismes canoniques
 $\Spec(\mathcal{O}_{X, x}) \to X$. Et le contenu du Lemme  \ref{lemma-specialize-points}  est, dans ce cadre, que le
morphisme  $\Spec(\kappa(x')) \to X$  se factorise comme  $\Spec(\kappa(x')) \to \Spec(\mathcal{O}_{X, x}) \to X$  chaque fois que
 $x'$  est une générisation de  $x$. Dans le cas où nous avons un
morphisme de schémas  $f : X \to S$, et un point  $x$  dont l'image est un point
 $s \in S$ , nous obtenons un diagramme commutatif
$$
\xymatrix{
\Spec(\kappa(x)) \ar[r] \ar[d] &
\Spec(\mathcal{O}_{X, x}) \ar[r] \ar[d] &
X \ar[d] \\
\Spec(\kappa(s)) \ar[r] &
\Spec(\mathcal{O}_{S, s}) \ar[r] &
S.
}
$$










\section{Recollement de schémas}
\label{section-glueing-schemes}

\noindent
Soit  $I$  un ensemble. Pour chaque  $i \in I$ , soit  $(X_i, \mathcal{O}_i)$  un
espace localement annelé. (En réalité, la construction qui suit fonctionne tout
aussi bien pour les espaces annelés.) Pour chaque paire  $i, j \in I$  soit
 $U_{ij} \subset X_i$  un sous-espace ouvert. Pour chaque paire  $i, j \in I$, soit
$$
\varphi_{ij} : U_{ij} \to U_{ji}
$$
un isomorphisme d'espaces localement annelés. Pour plus de commodité, nous
supposons que  $U_{ii} = X_i$. Pour chaque triple  $i, j, k \in I$ , supposons que
\begin{enumerate}
\item nous avons  $\varphi_{ij}^{-1}(U_{ji} \cap U_{jk}) =  U_{ij} \cap U_{ik}$, et
\item le diagramme
$$
\xymatrix{
U_{ij} \cap U_{ik} \ar[rr]_{\varphi_{ik}} \ar[rd]_{\varphi_{ij}} & &
U_{ki} \cap U_{kj} \\
& U_{ji} \cap U_{jk} \ar[ru]_{\varphi_{jk}}
}
$$
est commutatif.
\end{enumerate}
(En particulier, il en découle que  $\varphi_{ii} = \text{id}_{X_i}$.) Appelons une collection
 $(I, (X_i)_{i\in I}, (U_{ij})_{i, j\in I}, (\varphi_{ij})_{i, j\in I})$  satisfaisant les conditions ci-dessus des données de recollement.

\begin{lemma}
\label{lemma-glue}
\begin{slogan}
Si deux espaces localement annelés possèdent des sous-espaces isomorphes, alors
on peut les recoller pour former un
grand espace localement annelé.
\end{slogan}
Étant donnée une famille de données de recollement d'espaces localement annelés, il existe un
espace localement annelé  $X$  et des sous-espaces ouverts  $U_i \subset X$ 
avec des isomorphismes  $\varphi_i : X_i \to U_i$  d'espaces localement annelés tels que
\begin{enumerate}
\item $X=\bigcup_{i\in I} U_i$,
\item $\varphi_i(U_{ij}) = U_i \cap U_j$, et
\item $\varphi_{ij} =
\varphi_j^{-1}|_{U_i \cap U_j} \circ \varphi_i|_{U_{ij}}$.
\end{enumerate}
L'espace localement annelé  $X$  est caractérisé par les propriétés
universelles suivantes : étant donné un espace localement annelé  $Y$ , nous
avons
\begin{eqnarray*}
\Mor(X, Y) & = & \{ (f_i)_{i\in I} \mid
f_i : X_i \to Y, \ f_j \circ \varphi_{ij} = f_i|_{U_{ij}}\} \\
f & \mapsto & (f|_{U_i} \circ \varphi_i)_{i \in I} \\
\Mor(Y, X) & = &
\left\{
\begin{matrix}
\text{recouvrement ouvert }Y = \bigcup\nolimits_{i \in I} V_i\text{ et }
(g_i : V_i \to X_i)_{i \in I}
\text{ tels que}\\
g_i^{-1}(U_{ij}) = V_i \cap V_j
\text{ et }
g_j|_{V_i \cap V_j} = \varphi_{ij} \circ g_i|_{V_i \cap V_j}
\end{matrix}
\right\} \\
g & \mapsto &
V_i = g^{-1}(U_i), \ g_i = \varphi_i^{-1} \circ g|_{V_i}
\end{eqnarray*}
\end{lemma}

\begin{proof}
Nous construisons  $X$  en plusieurs étapes. Comme ensemble, nous prenons
$$
X = (\coprod X_i) / \sim.
$$
Ici, étant donné  $x \in X_i$  et  $x' \in X_j$ , nous disons  $x \sim x'$  si et
seulement si  $x \in U_{ij}$,  $x' \in U_{ji}$  et  $\varphi_{ij}(x) = x'$. Il s'agit d'une relation
d'équivalence car si  $x \in X_i$,  $x' \in X_j$,  $x'' \in X_k$,  $x \sim x'$  et
 $x' \sim x''$, alors  $x' \in U_{ji} \cap U_{jk}$, donc, d'après la condition (1) des
données de recollement,  $x \in U_{ij} \cap U_{ik}$  et  $x'' \in U_{ki} \cap U_{kj}$ ; d'après la condition (2), nous voyons
que  $\varphi_{ik}(x) = x''$. (La réflexivité et la symétrie découlent de nos hypothèses que
 $U_{ii} = X_i$  et  $\varphi_{ii} = \text{id}_{X_i}$.) On note  $\varphi_i : X_i \to X$  les applications naturelles.
On note  $U_i = \varphi_i(X_i) \subset X$. Notez que  $\varphi_i : X_i \to U_i$  est une bijection.

\medskip\noindent
La topologie sur  $X$  est définie par la règle selon laquelle  $U \subset X$ 
est ouvert si et seulement si  $\varphi_i^{-1}(U)$  est ouvert pour tout  $i$. Nous
laissons au lecteur le soin de vérifier que cela définit effectivement une
topologie. Notez qu'en particulier  $U_i$  est ouvert puisque  $\varphi_j^{-1}(U_i)
= U_{ji}$ 
qui est ouvert dans  $X_j$  pour tout  $j$. De plus, pour tout
ensemble ouvert  $W \subset X_i$ , l'image  $\varphi_i(W) \subset U_i$  est ouverte parce que
 $\varphi_j^{-1}(\varphi_i(W)) = \varphi_{ji}^{-1}(W \cap U_{ij})$. Par conséquent,  $\varphi_i : X_i \to U_i$  est un homéomorphisme.

\medskip\noindent
Pour obtenir un espace localement annelé, nous devons construire le faisceau
d'anneaux  $\mathcal{O}_X$. Nous le faisons en recollant les faisceaux d'anneaux
 $\mathcal{O}_{U_i} := \varphi_{i, *} \mathcal{O}_i$. Plus précisément, dans le diagramme commutatif
$$
\xymatrix{
U_{ij} \ar[rr]_{\varphi_{ij}} \ar[rd]_{\varphi_i|_{U_{ij}}}
& &
U_{ji} \ar[ld]^{\varphi_j|_{U_{ji}}} \\
& U_i \cap U_j &
}
$$
la flèche du haut est un isomorphisme d'espaces annelés et, par conséquent, nous
obtenons des isomorphismes uniques de faisceaux d'anneaux
$$
\mathcal{O}_{U_i}|_{U_i \cap U_j}
\longrightarrow
\mathcal{O}_{U_j}|_{U_i \cap U_j}.
$$
Ces isomorphismes satisfont une condition de cocycle comme dans Faisceaux, Section
 \ref{sheaves-section-glueing-sheaves}. D'après les résultats de cette section, nous obtenons un faisceau
d'anneaux  $\mathcal{O}_X$  sur  $X$  tel que  $\mathcal{O}_X|_{U_i}$  soit isomorphe à
 $\mathcal{O}_{U_i}$  de manière compatible avec les morphismes de recollement affichés
ci-dessus. En particulier,  $(X, \mathcal{O}_X)$  est un espace localement annelé puisque
les fibres de  $\mathcal{O}_X$  sont égales aux fibres de  $\mathcal{O}_i$  aux points
correspondants.

\medskip\noindent
La démonstration des propriétés universelles est omise.
\end{proof}

\begin{lemma}
\label{lemma-glue-schemes}
\begin{slogan}
Les schémas peuvent être recollés pour donner de nouveaux schémas.
\end{slogan}
Dans le Lemme  \ref{lemma-glue}  ci-dessus, supposons que tous les  $X_i$  soient des
schémas. Alors l'espace localement annelé résultant  $X$  est un schéma.
\end{lemma}

\begin{proof}
Cela est clair puisque chacun des  $U_i$  est un schéma et donc chaque
 $x \in X$  possède un voisinage affine.
\end{proof}

\noindent
Il est d'usage de considérer  $X_i$  comme un sous-espace ouvert de
 $X$  via les isomorphismes  $\varphi_i$. Nous ferons cela dans les deux
exemples suivants.

\begin{example}[ Espace affine à origine dédoublée]
\label{example-affine-space-zero-doubled}
Soit  $k$  un corps. Soit  $n \geq 1$. Soit  $X_1 = \Spec(k[x_1, \ldots, x_n])$, soit
 $X_2 = \Spec(k[y_1, \ldots, y_n])$. Soit  $0_1 \in X_1$  le point correspondant à l'idéal maximal
 $(x_1, \ldots, x_n) \subset k[x_1, \ldots, x_n]$. Soit  $0_2 \in X_2$  le point correspondant à l'idéal maximal
 $(y_1, \ldots, y_n) \subset k[y_1, \ldots, y_n]$. Soit  $U_{12} = X_1 \setminus \{0_1\}$  et soit  $U_{21} = X_2 \setminus \{0_2\}$. Soit  $\varphi_{12} : U_{12} \to U_{21}$ 
l'isomorphisme provenant de l'isomorphisme des  $k$-algèbres  $k[y_1, \ldots, y_n] \to k[x_1, \ldots, x_n]$ 
envoyant  $y_i$  sur  $x_i$  (ce qui induit  $X_1 \cong X_2$  envoyant
 $0_1$  sur  $0_2$). Soit  $X$  le schéma obtenu à partir des
données de recollement  $(X_1, X_2, U_{12}, U_{21}, \varphi_{12},
\varphi_{21} = \varphi_{12}^{-1})$. Par le léger abus de notation introduit ci-dessus
dans l'exemple, nous considérons  $X_1, X_2 \subset X$  comme des sous-schémas ouverts. Il
existe un morphisme  $f : X \to \Spec(k[t_1, \ldots, t_n])$  qui sur  $X_1$  (resp.
\ $X_2$) correspond à l'homomorphisme de  $k$-algèbres
 $k[t_1, \ldots, t_n] \to k[x_1, \ldots, x_n]$  (resp. \ $k[t_1, \ldots, t_n] \to k[y_1, \ldots, y_n]$) envoyant  $t_i$  vers
 $x_i$  (resp. \   $t_i$  vers  $y_i$). Il est facile de
voir que ce morphisme identifie  $k[t_1, \ldots, t_n]$  avec  $\Gamma(X, \mathcal{O}_X)$. Puisque
 $f(0_1) = f(0_2)$ , nous voyons que  $X$  n'est pas affine.

\medskip\noindent
Notez que  $X_1$  et  $X_2$  sont des ouverts affines de  $X$.
Mais, si  $n = 2$, alors  $X_1 \cap X_2$  est le schéma décrit dans l'Exemple
 \ref{example-not-affine}  et donc pas affine. Ainsi, en général, l'intersection des ouverts
affines d'un schéma n'est pas affine. (Ce fait vaut plus généralement pour tout
 $n > 1$.)

\medskip\noindent
Une autre caractéristique curieuse de cet exemple est la suivante. Si  $n > 1$ 
il y a de nombreux sous-ensembles fermés irréductibles  $T \subset X$  (prenez par
exemple l'adhérence de tout point non fermé dans  $X_1$ ). Mais à
moins que  $T = \{0_1\}$, ou  $T = \{0_2\}$  nous avons  $0_1 \in T \Leftrightarrow 0_2 \in T$. Démonstration omise.
\end{example}

\begin{example}[ Droite projective]
\label{example-projective-line}
Soit  $k$  un corps. Soit  $X_1 = \Spec(k[x])$, soit  $X_2 = \Spec(k[y])$. Soit
 $0 \in X_1$  le point correspondant à l'idéal maximal  $(x) \subset k[x]$. Soit
 $\infty \in X_2$  le point correspondant à l'idéal maximal  $(y) \subset k[y]$. Soit
 $U_{12} = X_1 \setminus \{0\} = D(x) = \Spec(k[x, 1/x])$  et soit  $U_{21} = X_2 \setminus \{\infty\} = D(y) = \Spec(k[y, 1/y])$. Soit  $\varphi_{12} : U_{12} \to U_{21}$  l'isomorphisme provenant de
l'isomorphisme des  $k$-algèbres  $k[y, 1/y] \to k[x, 1/x]$  envoyant  $y$  sur
 $1/x$. Soit  $\mathbf{P}^1_k$  le schéma obtenu à partir des données de
recollement  $(X_1, X_2, U_{12}, U_{21}, \varphi_{12},
\varphi_{21} = \varphi_{12}^{-1})$. Via le léger abus de notation introduit ci-dessus dans
l'exemple, nous considérons  $X_i \subset \mathbf{P}^1_k$  comme des sous-schémas ouverts. Dans ce
cas, nous voyons que  $\Gamma(\mathbf{P}^1_k, \mathcal{O}) = k$  car les seuls polynômes  $g(x)$  en
 $x$  tels que  $g(1/y)$  soit également un polynôme en  $y$ 
sont des polynômes constants. Puisque  $\mathbf{P}^1_k$  est infini, nous voyons que
 $\mathbf{P}^1_k$  n'est pas affine.

\medskip\noindent
Nous affirmons qu'il existe un ouvert affine  $U \subset \mathbf{P}^1_k$  qui contient à la fois
 $0$  et  $\infty$. À savoir, soit  $U = \mathbf{P}^1_k \setminus \{1\}$, où  $1$  est le
point de  $X_1$  correspondant à l'idéal maximal  $(x - 1)$  et également le
point de  $X_2$  correspondant à l'idéal maximal  $(y - 1)$. Il est alors
facile de voir que  $s = 1/(x - 1) = y/(1 - y) \in \Gamma(U, \mathcal{O}_U)$. En fait, vous pouvez montrer que  $\Gamma(U, \mathcal{O}_U)$  est
égal à l'anneau de polynômes  $k[s]$  et que le morphisme correspondant
 $U \to \Spec(k[s])$  est un isomorphisme de schémas. Détails omis.
\end{example}




\section{Un critère de représentabilité}
\label{section-representable}

\noindent
Dans cette section, nous reformulons le lemme de recollement de la Section
 \ref{section-glueing-schemes}  en termes de foncteurs. Nous rappelons une partie du matériel de
Catégories, Section  \ref{categories-section-opposite}. Rappelons qu'étant donné un schéma  $X$ ,
nous pouvons définir un foncteur
$$
h_X : \Sch^{opp}
\longrightarrow
\textit{Ensembles}, \quad
T \longmapsto \Mor(T, X).
$$
On l'appelle le foncteur des points  {\it  de
 $X$}.

\medskip\noindent
Soit  $F$  un foncteur contravariant de la catégorie des schémas vers la
catégorie des ensembles. Dans une formule
$$
F : \Sch^{opp}
\longrightarrow
\textit{Ensembles}.
$$
Nous utiliserons la même terminologie que dans Sites, Section  \ref{sites-section-presheaves}. À
savoir, étant donné un schéma  $T$, un élément  $\xi \in F(T)$ et un
morphisme  $f : T' \to T$ , nous noterons  $f^*\xi$  l'élément  $F(f)(\xi)$, et
parfois nous utiliserons même la notation  $\xi|_{T'}$

\begin{definition}
\label{definition-representable-functor}
(Voir Catégories, Définition  \ref{categories-definition-representable-functor}.) Soit  $F$  un foncteur
contravariant de la catégorie des schémas vers la catégorie des ensembles (comme
ci-dessus). Nous disons que  $F$  est  {\it 
représentable par un schéma }  ou  {\it  représentable
}  s'il existe un schéma  $X$  tel que  $h_X \cong F$.
\end{definition}

\noindent
Supposons que  $F$  soit représentable par le schéma  $X$  et que
 $s : h_X \to F$  soit un isomorphisme. D'après les Catégories, le lemme de Yoneda
 \ref{categories-lemma-yoneda} , la paire  $(X, s : h_X \to F)$  est unique à isomorphisme unique près si elle
existe. De plus, le lemme de Yoneda affirme que pour tout foncteur contravariant
 $F$  comme ci-dessus et tout schéma  $Y$, nous avons une bijection
$$
\Mor_{\text{Fun}(\Sch^{opp}, \textit{Ensembles})} (h_Y, F)
\longrightarrow
F(Y), \quad
s \longmapsto s(\text{id}_Y).
$$
Voici la construction inverse. Étant donné un  $\xi \in F(Y)$ , la transformation de
foncteurs  $s_\xi : h_Y \to F$  associe à tout morphisme  $f : T \to Y$  l'élément
 $f^*\xi \in F(T)$.

\medskip\noindent
En particulier, dans le cas où  $F$  est représentable, il existe un schéma
 $X$  et un élément  $\xi \in F(X)$  tels que le morphisme correspondant
 $h_X \to F$  soit un isomorphisme. Dans ce cas, nous disons aussi
 {\it  que la paire  $(X, \xi)$  représente
 $F$}. L'élément  $\xi \in F(X)$  est souvent appelé la
 {\it  ``famille universelle'' }  pour des raisons qui
deviendront plus claires lorsque nous parlerons des champs algébriques (insérer
référence future ici). Pour le moment, observons simplement que, si la paire
 $(X, \xi)$  représente  $F$, alors tout élément  $\xi' \in F(T)$ 
pour tout  $T$  est de la forme  $\xi' = f^*\xi$  pour un morphisme unique
 $f : T \to X$.

\begin{example}
\label{example-global-sections}
Considérons la règle qui associe à chaque schéma  $T$  l'ensemble
 $F(T) = \Gamma(T, \mathcal{O}_T)$. Nous pouvons en faire un foncteur contravariant en utilisant, pour
un morphisme  $f : T' \to T$ , l'homomorphisme d'image inverse  $f^\sharp : \Gamma(T, \mathcal{O}_T) \to \Gamma(T', \mathcal{O}_{T'})$. Étant
donné un anneau  $R$  et un élément  $t \in R$ , il existe un unique
homomorphisme d'anneaux  $\mathbf{Z}[x] \to R$  qui envoie  $x$  sur  $t$.
Ainsi, en utilisant le Lemme  \ref{lemma-morphism-into-affine}, nous constatons que
$$
\Mor(T, \Spec(\mathbf{Z}[x])) =
\Hom(\mathbf{Z}[x], \Gamma(T, \mathcal{O}_T)) =
\Gamma(T, \mathcal{O}_T).
$$
Cela donne effectivement un isomorphisme  $h_{\Spec(\mathbf{Z}[x])} \to F$. Quelle est la « famille
universelle »  $\xi$? Pour l'obtenir, nous devons appliquer les
identifications ci-dessus à  $\text{id}_{\Spec(\mathbf{Z}[x])}$. Clairement, sous les identifications
ci-dessus, cela donne que  $\xi = x \in \Gamma(\Spec(\mathbf{Z}[x]),
\mathcal{O}_{\Spec(\mathbf{Z}[x])}) = \mathbf{Z}[x]$  comme prévu.
\end{example}

\begin{definition}
\label{definition-representable-by-open-immersions}
Soit  $F$  un foncteur contravariant sur la catégorie des schémas à valeurs
dans les ensembles.
\begin{enumerate}
\item On dit que  $F$   {\it  satisfait la propriété de
faisceau pour la topologie de Zariski}  si pour tout schéma  $T$ 
et tout recouvrement ouvert  $T = \bigcup_{i \in I} U_i$, et pour toute collection d'éléments
 $\xi_i \in F(U_i)$  telle que  $\xi_i|_{U_i \cap U_j} =
\xi_j|_{U_i \cap U_j}$ , il existe un élément unique  $\xi \in F(T)$  tel
que  $\xi_i = \xi|_{U_i}$  dans  $F(U_i)$.
\item Un  {\it  sous-foncteur  $H \subset F$}  est une règle
qui associe à chaque schéma  $T$  un sous-ensemble  $H(T) \subset F(T)$  tel que les
applications  $F(f) : F(T) \to F(T')$  envoient  $H(T)$  dans  $H(T')$  pour tous les
morphismes de schémas  $f : T' \to T$.
\item Soit  $H \subset F$  un sous-foncteur. On dit que  $H \subset F$  est
 {\it  représentable par des immersions ouvertes}  si
pour tous les couples  $(T, \xi)$, où  $T$  est un schéma et  $\xi \in F(T)$ ,
il existe un sous-schéma ouvert  $U_\xi \subset T$  avec la propriété suivante:
\begin{itemize}
\item[(*)] Un morphisme  $f : T' \to T$  se factorise à travers  $U_\xi$  si et seulement si
 $f^*\xi \in H(T')$.
\end{itemize}
\item Soit  $I$  un ensemble. Pour chaque  $i \in I$ , soit  $H_i \subset F$  un
sous-foncteur. Nous disons que la collection  $(H_i)_{i \in I}$   {\it 
recouvre  $F$}  si et seulement si pour chaque  $\xi \in F(T)$ , il
existe un recouvrement ouvert  $T = \bigcup U_i$  tel que  $\xi|_{U_i} \in H_i(U_i)$.
\end{enumerate}
\end{definition}

\noindent
Dans la condition (4), si  $H_i \subset F$  est représentable par des immersions
ouvertes pour tous les  $i$, alors pour vérifier que  $(H_i)_{i \in I}$  recouvre
 $F$, il suffit de vérifier  $F(T) = \bigcup H_i(T)$  chaque fois que  $T$  est
le spectre d'un corps.

\begin{lemma}
\label{lemma-glue-functors}
Soit  $F$  un foncteur contravariant sur la catégorie des schémas avec des
valeurs dans la catégorie des ensembles. Supposons que
\begin{enumerate}
\item $F$  satisfasse la propriété de faisceau pour la topologie de Zariski.
\item il existe un ensemble  $I$  et une collection de sous-foncteurs
 $F_i \subset F$  tels que
\begin{enumerate}
\item chaque  $F_i$  est représentable,
\item chaque  $F_i \subset F$  est représentable par des immersions ouvertes, et
\item la collection  $(F_i)_{i \in I}$  recouvre  $F$.
\end{enumerate}
\end{enumerate}
Alors  $F$  est représentable.
\end{lemma}

\begin{proof}
Soit  $X_i$  un schéma représentant  $F_i$  et soit  $\xi_i \in F_i(X_i) \subset F(X_i)$  la
``famille universelle''. Puisque  $F_j \subset F$  est représentable par des immersions
ouvertes, il existe un ouvert  $U_{ij} \subset X_i$  tel que  $T \to X_i$  se factorise par
 $U_{ij}$  si et seulement si  $\xi_i|_T \in F_j(T)$. En particulier  $\xi_i|_{U_{ij}} \in F_j(U_{ij})$  et donc
nous obtenons un morphisme canonique  $\varphi_{ij} : U_{ij} \to X_j$  tel que  $\varphi_{ij}^*\xi_j = \xi_i|_{U_{ij}}$. Par
définition de  $U_{ji}$ , cela implique que  $\varphi_{ij}$  se factorise par
 $U_{ji}$. Comme  $(\varphi_{ij} \circ \varphi_{ji})^*\xi_j
=\varphi_{ji}^*(\varphi_{ij}^*\xi_j) =
\varphi_{ji}^*\xi_i = \xi_j$ , nous concluons que  $\varphi_{ij} \circ \varphi_{ji} = \text{id}_{U_{ji}}$  parce que le
couple  $(X_j, \xi_j)$  représente  $F_j$. En particulier, les applications
 $\varphi_{ij} : U_{ij} \to U_{ji}$  sont des isomorphismes de schémas. Ensuite, nous devons montrer que
 $\varphi_{ij}^{-1}(U_{ji} \cap U_{jk}) = U_{ij} \cap U_{ik}$. Cela est vrai parce que (a)  $U_{ji} \cap U_{jk}$  est le plus grand ouvert
de  $U_{ji}$  tel que  $\xi_j$  se restreigne à un élément de  $F_k$,
(b)  $U_{ij} \cap U_{ik}$  est le plus grand ouvert de  $U_{ij}$  tel que  $\xi_i$ 
se restreigne à un élément de  $F_k$, et (c)  $\varphi_{ij}^*\xi_j = \xi_i$. De plus, la
condition de cocycle dans la Section  \ref{section-glueing-schemes}  suit parce que  $\varphi_{jk}|_{U_{ji} \cap U_{jk}} \circ
\varphi_{ij}|_{U_{ij} \cap U_{ik}}$  et
 $\varphi_{ik}|_{U_{ij} \cap U_{ik}}$  ont tous deux pour image inverse de  $\xi_k$  l'élément  $\xi_i$. Ainsi, nous pouvons
appliquer le Lemme  \ref{lemma-glue-schemes}  pour obtenir un schéma  $X$  avec un
recouvrement ouvert  $X = \bigcup U_i$  et des isomorphismes  $\varphi_i : X_i \to U_i$  avec des
propriétés comme dans le Lemme  \ref{lemma-glue}. Soit  $\xi_i' = (\varphi_i^{-1})^* \xi_i$. Les conditions du
Lemme  \ref{lemma-glue}  impliquent que  $\xi_i'|_{U_i \cap U_j} = \xi_j'|_{U_i \cap U_j}$. Par conséquent, d'après la
condition que  $F$  satisfait la condition de faisceau dans la topologie de
Zariski, nous voyons qu'il existe un élément  $\xi' \in F(X)$  tel que  $\xi_i = \varphi_i^*\xi'|_{U_i}$ 
pour tout  $i$. Puisque  $\varphi_i$  est un isomorphisme, nous obtenons
également que  $(U_i, \xi'|_{U_i})$  représente le foncteur  $F_i$.

\medskip\noindent
Nous affirmons que la paire  $(X, \xi')$  représente le foncteur  $F$. Pour
le montrer, soit  $T$  un schéma et soit  $\xi \in F(T)$. Nous allons
construire un morphisme unique  $g : T \to X$  tel que  $g^*\xi' = \xi$. À savoir,
d’après la condition que les sous-foncteurs  $F_i$  couvrent  $F$ , il
existe un recouvrement ouvert  $T = \bigcup V_i$  tel que pour chaque  $i$ , la
restriction  $\xi|_{V_i} \in F_i(V_i)$. De plus, comme chacune des inclusions  $F_i \subset F$  est
représentable par des immersions ouvertes, nous pouvons supposer que chaque
 $V_i \subset T$  est le plus grand ouvert possédant cette propriété. Puisque
 $(U_i, \xi'|_{U_i})$  représente le foncteur  $F_i$ , nous obtenons un morphisme
unique  $g_i : V_i \to U_i$  tel que  $g_i^*\xi'|_{U_i} = \xi|_{V_i}$. Sur les intersections  $V_i \cap V_j$ , les
morphismes  $g_i$  et  $g_j$  coïncident, par exemple parce qu'ils
ont tous deux pour image inverse de  $\xi'|_{U_i \cap U_j} \in F_i(U_i \cap U_j)$  le même élément. Ainsi les morphismes
 $g_i$  se recollent en un unique morphisme  $T \to X$  comme
souhaité.
\end{proof}

\begin{remark}
\label{remark-representable-locally-ringed}
Supposons que le foncteur  $F$  soit défini sur tous les espaces localement
annelés, et que les conditions du Lemme  \ref{lemma-glue-functors}  soient remplacées par ce qui
suit:
\begin{enumerate}
\item $F$  satisfait la propriété de faisceau sur la catégorie des espaces
localement annelés,
\item il existe un ensemble  $I$  et une collection de sous-foncteurs
 $F_i \subset F$  tels que
\begin{enumerate}
\item chaque  $F_i$  soit représentable par un schéma,
\item chaque  $F_i \subset F$  soit représentable par des immersions ouvertes sur la
catégorie des espaces localement annelés, et
\item la collection  $(F_i)_{i \in I}$  recouvre  $F$  en tant que foncteur sur la
catégorie des espaces localement annelés.
\end{enumerate}
\end{enumerate}
Nous laissons au lecteur le soin de développer cela davantage. Le résultat final
est alors que le foncteur  $F$  est représentable dans la catégorie des
espaces localement annelés et que l'objet représentant est un schéma.
\end{remark}



\section{Existence des produits fibrés de schémas}
\label{section-existence-fibre-products}

\noindent
Une question élémentaire est de savoir si les produits et les produits fibrés
existent dans la catégorie des schémas. Nous prouvons d'abord leur existence de
manière abstraite ; dans la section suivante, nous en donnons une description
concrète.

\begin{lemma}
\label{lemma-fibre-products}
La catégorie des schémas possède un objet final, des produits et des produits
fibrés. En d'autres termes, la catégorie des schémas a des limites finies, voir
Catégories, Lemme  \ref{categories-lemma-finite-limits-exist}.
\end{lemma}

\begin{proof}
On peut omettre cette démonstration. Il est plus important d'apprendre à
travailler avec le produit fibré, qui est expliqué dans la section suivante.

\medskip\noindent
D'après le Lemme  \ref{lemma-morphism-into-affine} , le schéma  $\Spec(\mathbf{Z})$  est un objet final dans la
catégorie des espaces localement annelés. Ainsi, il suffit de prouver que les
produits fibrés existent.

\medskip\noindent
Soient  $f : X \to S$  et  $g : Y \to S$  des morphismes de schémas. Nous devons montrer
que le foncteur
\begin{eqnarray*}
F : \Sch^{opp} & \longrightarrow & \textit{Ensembles} \\
T & \longmapsto &
\Mor(T, X) \times_{\Mor(T, S)} \Mor(T, Y)
\end{eqnarray*}
est représentable. Nous affirmons que le Lemme  \ref{lemma-glue-functors}  s'applique au foncteur
 $F$. Si nous prouvons cela, alors le lemme est prouvé.

\medskip\noindent
Tout d'abord, nous montrons que  $F$  satisfait la propriété de faisceau
dans la topologie de Zariski. À savoir, supposons que  $T$  soit un schéma,
 $T = \bigcup_{i \in I} U_i$  soit un recouvrement ouvert, et  $\xi_i \in F(U_i)$  tel que  $\xi_i|_{U_i \cap U_j} =  \xi_j|_{U_i \cap U_j}$ 
pour toutes les paires  $i, j$. Par définition,  $\xi_i$  correspond à une
paire  $(a_i, b_i)$  où  $a_i : U_i \to X$  et  $b_i : U_i \to Y$  sont des morphismes de schémas
tels que  $f \circ a_i = g \circ b_i$. La condition de recollement dit que  $a_i|_{U_i \cap U_j} =  a_j|_{U_i \cap U_j}$  et
 $b_i|_{U_i \cap U_j} =  b_j|_{U_i \cap U_j}$. Ainsi, en recollant les morphismes  $a_i$ , nous obtenons un
morphisme d'espaces localement annelés (c'est-à-dire un morphisme de schémas)
 $a : T \to X$  et de même  $b : T \to Y$  (voir par exemple la propriété universelle du
Lemme  \ref{lemma-glue}). De plus, sur les éléments d'un recouvrement ouvert, les
compositions  $f \circ a$  et  $g \circ b$  coïncident. Par conséquent,  $f \circ a = g \circ b$ 
et la paire  $(a, b)$  définit un élément de  $F(T)$  qui se restreint
aux paires  $(a_i, b_i)$  sur chaque  $U_i$. La condition de faisceau est
vérifiée.

\medskip\noindent
Ensuite, nous construisons la famille de sous-foncteurs. Choisissez un
recouvrement ouvert par des ouverts affines  $S = \bigcup\nolimits_{i \in I} U_i$. Pour chaque
 $i \in I$ , choisissez des recouvrements ouverts par des ouverts affines
 $f^{-1}(U_i) = \bigcup\nolimits_{j \in J_i} V_j$  et  $g^{-1}(U_i) = \bigcup\nolimits_{k \in K_i} W_k$. Notez que  $X = \bigcup_{i \in I} \bigcup_{j \in J_i} V_j$  est un recouvrement ouvert et
de même pour  $Y$. Pour tout  $i \in I$  et chaque paire  $(j, k) \in J_i \times K_i$ ,
nous avons un diagramme commutatif
$$
\xymatrix{
    & W_k \ar[d] \ar[rd] &   \\
V_j \ar[rd] \ar[r] & U_i \ar[rd] & Y \ar[d] \\
    & X \ar[r]  & S
}
$$
où toutes les flèches obliques sont des immersions ouvertes. Pour un tel triple,
nous obtenons un foncteur
\begin{eqnarray*}
F_{i, j, k} : \Sch^{opp} & \longrightarrow & \textit{Ensembles} \\
T & \longmapsto &
\Mor(T, V_j) \times_{\Mor(T, U_i)} \Mor(T, W_k).
\end{eqnarray*}
Il existe une transformation évidente de foncteurs  $F_{i, j, k} \to F$  (provenant du
grand diagramme commutatif ci-dessus) qui est injective, donc nous pouvons
considérer  $F_{i, j, k}$  comme un sous-foncteur de  $F$.

\medskip\noindent
Nous vérifions la condition (2)(a) du Lemme  \ref{lemma-glue-functors}. Cela découle directement
du Lemme  \ref{lemma-fibre-product-affine-schemes}. (Notez que nous utilisons ici le fait que les produits
fibrés dans la catégorie des schémas affines sont également des produits fibrés
dans la catégorie tout entière des espaces localement annelés.)

\medskip\noindent
Nous vérifions la condition (2)(b) du Lemme  \ref{lemma-glue-functors}. Soit  $T$  un
schéma et soit  $\xi \in F(T)$. En d'autres termes,  $\xi = (a, b)$  où  $a : T \to X$  et
 $b : T \to Y$  sont des morphismes de schémas tels que  $f \circ a = g \circ b$. Posons
 $V_{i, j, k} = a^{-1}(V_j) \cap b^{-1}(W_k)$. Pour tout autre morphisme  $h : T' \to T$  nous avons  $h^*\xi = (a \circ h, b \circ h)$. Par
conséquent, nous voyons que  $h^*\xi \in F_{i, j, k}(T')$  si et seulement si  $a(h(T')) \subset V_j$  et
 $b(h(T')) \subset W_k$. En d'autres termes, si et seulement si  $h(T') \subset V_{i, j, k}$. Cela prouve la
condition (2)(b).

\medskip\noindent
Nous vérifions la condition (2)(c) du Lemme  \ref{lemma-glue-functors}. Soit  $T$  un
schéma et soit  $\xi = (a, b) \in F(T)$  comme ci-dessus. Posons  $V_{i, j, k} = a^{-1}(V_j) \cap b^{-1}(W_k)$  comme ci-dessus.
La condition (2)(c) signifie simplement que  $T = \bigcup V_{i, j, k}$  ce qui est évident.
Ainsi, le lemme est prouvé et les produits fibrés existent.
\end{proof}

\begin{remark}
\label{remark-fibre-product-schemes-locally-ringed}
En utilisant la Remarque  \ref{remark-representable-locally-ringed} , on peut montrer que le produit fibré de
morphismes de schémas existe dans la catégorie des espaces localement annelés et
est un schéma.
\end{remark}




\section{Produits fibrés de schémas}
\label{section-fibre-products}

\noindent
Voici un rappel de la définition générale, même si nous avons déjà montré que les
produits fibrés de schémas existent.

\begin{definition}
\label{definition-fibre-product}
Étant donnés des morphismes de schémas $f : X \to S$  et  $g : Y \to S$, le
 {\it produit fibré}  est un schéma $X \times_S Y$ muni
des morphismes de projection $p : X \times_S Y \to X$  et  $q : X \times_S Y \to Y$ qui s'insèrent dans le
diagramme commutatif suivant
$$
\xymatrix{
X \times_S Y \ar[r]_q \ar[d]_p & Y \ar[d]^g \\
X \ar[r]^f & S
}
$$
et qui est universel parmi tous les diagrammes de ce type, voir Catégories,
Définition \ref{categories-definition-fibre-products}.
\end{definition}

\noindent
En d'autres termes, étant donné tout diagramme commutatif de morphismes de
schémas dont les flèches pleines sont fixées
$$
\xymatrix{
T \ar[rrrd] \ar@{-->}[rrd] \ar[rrdd]
&
&
\\
&
&
X \times_S Y \ar[d] \ar[r]
&
Y \ar[d]
\\
&
&
X \ar[r]
&
S
}
$$
il existe une flèche pointillée unique faisant commuter le diagramme. Nous allons
prouver quelques lemmes qui nous aideront à concevoir les produits fibrés.

\begin{lemma}
\label{lemma-fibre-product-affines}
Soient $f : X \to S$  et  $g : Y \to S$ des morphismes de schémas ayant la même
cible. Si $X, Y, S$ sont tous affines alors $X \times_S Y$  est affine.
\end{lemma}

\begin{proof}
Supposons que  $X = \Spec(A)$,  $Y = \Spec(B)$  et  $S = \Spec(R)$. D'après le Lemme
 \ref{lemma-fibre-product-affine-schemes} , le schéma affine  $\Spec(A \otimes_R B)$  est le produit fibré  $X \times_S Y$  dans
la catégorie des espaces localement annelés. Par conséquent, il est a fortiori le
produit fibré dans la catégorie des schémas.
\end{proof}

\begin{lemma}
\label{lemma-open-fibre-product}
Soient  $f : X \to S$  et  $g : Y \to S$  des morphismes de schémas ayant la même cible.
Soient  $X \times_S Y$,  $p$,  $q$  le produit fibré. Supposons que
 $U \subset S$,  $V \subset X$,  $W \subset Y$  soient des sous-schémas ouverts tels que
 $f(V) \subset U$  et  $g(W) \subset U$. Alors le morphisme canonique  $V \times_U W \to X \times_S Y$  est une
immersion ouverte qui identifie  $V \times_U W$  avec  $p^{-1}(V) \cap q^{-1}(W)$.
\end{lemma}

\begin{proof}
Soit  $T$  un schéma. Supposons que  $a : T \to V$  et  $b : T \to W$  soient des
morphismes tels que  $f \circ a = g \circ b$  comme morphismes dans  $U$. Alors ils
coïncident comme morphismes dans  $S$. Par la propriété universelle du
produit fibré, nous obtenons un morphisme unique  $T \to X \times_S Y$. Bien sûr, ce
morphisme a son image contenue dans l'ouvert  $p^{-1}(V) \cap q^{-1}(W)$. Ainsi  $p^{-1}(V) \cap q^{-1}(W)$  est
un produit fibré de  $V$  et  $W$  au-dessus de  $U$. Le
résultat découle de l'unicité des produits fibrés, voir Catégories, Section
 \ref{categories-section-fibre-products}.
\end{proof}

\noindent
En particulier, cela montre que  $V \times_U W = V \times_S W$  dans la situation du lemme.
De plus, si  $U, V, W$  sont tous affines, alors nous savons que  $V \times_U W$  est
affine. Et bien sûr, nous pouvons recouvrir  $X \times_S Y$  par de tels ouverts
affines  $V \times_U W$. Nous formulons cela comme un lemme.

\begin{lemma}
\label{lemma-affine-covering-fibre-product}
\begin{slogan}
Construction explicite des produits fibrés : un recouvrement ouvert affine d'un
produit fibré de schémas peut être assemblé à partir de recouvrements ouverts
affines compatibles des facteurs.
\end{slogan}
Soient  $f : X \to S$  et  $g : Y \to S$  des morphismes de schémas ayant la même cible.
Soit  $S = \bigcup U_i$  un recouvrement ouvert affine quelconque de  $S$. Pour
chaque  $i \in I$, soit  $f^{-1}(U_i) = \bigcup_{j \in J_i} V_j$  un recouvrement ouvert affine de
 $f^{-1}(U_i)$  et soit  $g^{-1}(U_i) = \bigcup_{k \in K_i} W_k$  un recouvrement ouvert affine de  $g^{-1}(U_i)$.
Alors
$$
X \times_S Y =
\bigcup\nolimits_{i \in I}
\bigcup\nolimits_{j \in J_i, \ k \in K_i}
V_j \times_{U_i} W_k
$$
est un recouvrement ouvert affine de  $X \times_S Y$.
\end{lemma}

\begin{proof}
Voir la discussion au-dessus du lemme.
\end{proof}

\noindent
En d'autres termes, nous aurions pu utiliser le lemme précédent pour construire le
produit fibré directement en recollant les schémas affines. (Ce qui est bien sûr
exactement ce que nous avons fait dans la preuve du Lemme  \ref{lemma-fibre-products}.) Voici une manière de décrire l'ensemble des points d'un produit fibré de
schémas.

\begin{lemma}
\label{lemma-points-fibre-product}
Soient  $f : X \to S$  et  $g : Y \to S$  des morphismes de schémas ayant la même cible.
Les points  $z$  de  $X \times_S Y$  sont en correspondance bijective avec les
quadruples
$$
(x, y, s, \mathfrak p)
$$
où  $x \in X$,  $y \in Y$,  $s \in S$  sont des points tels que  $f(x) = s$,
 $g(y) = s$  et  $\mathfrak p$  est un idéal premier de l'anneau  $\kappa(x) \otimes_{\kappa(s)} \kappa(y)$. Le
corps résiduel de  $z$  correspond au corps résiduel de l'idéal premier
 $\mathfrak p$.
\end{lemma}

\begin{proof}
Soit  $z$  un point de  $X \times_S Y$  et construisons un quadruple comme
ci-dessus. Rappelons que nous pouvons considérer  $z$  comme un morphisme
 $\Spec(\kappa(z)) \to X \times_S Y$, voir le Lemme  \ref{lemma-characterize-points}. Ce morphisme correspond aux morphismes
 $a : \Spec(\kappa(z)) \to X$  et  $b : \Spec(\kappa(z)) \to Y$  tels que  $f \circ a = g \circ b$. Par le même lemme, nous
obtenons à nouveau les points  $x \in X$,  $y \in Y$  au-dessus du même point
 $s \in S$  ainsi que les homomorphismes de corps  $\kappa(x) \to \kappa(z)$,  $\kappa(y) \to \kappa(z)$  tels
que les compositions  $\kappa(s) \to \kappa(x) \to \kappa(z)$  et  $\kappa(s) \to \kappa(y) \to \kappa(z)$  soient les mêmes. En d'autres
termes, nous obtenons un homomorphisme d'anneaux  $\kappa(x) \otimes_{\kappa(s)} \kappa(y) \to \kappa(z)$. Nous notons
 $\mathfrak p$  le noyau de cette application.

\medskip\noindent
Inversement, étant donné un quadruple  $(x, y, s, \mathfrak p)$ , nous obtenons un diagramme
commutatif aux flèches pleines
$$
\xymatrix{
X \times_S Y
\ar@/_/[dddr] \ar@/^/[rrrd]
& & & \\
&
\Spec(\kappa(x) \otimes_{\kappa(s)} \kappa(y)/\mathfrak p)
\ar[r] \ar[d] \ar@{-->}[lu]
&
\Spec(\kappa(y)) \ar[d] \ar[r] &
Y \ar[dd] \\
&
\Spec(\kappa(x)) \ar[r] \ar[d] &
\Spec(\kappa(s)) \ar[rd] &
\\
&
X \ar[rr] &
&
S
}
$$
voir la discussion dans la Section  \ref{section-points}. Ainsi, nous obtenons la flèche en
pointillés. Le point correspondant  $z$  de  $X \times_S Y$  est l'image du
point générique de  $\Spec(\kappa(x) \otimes_{\kappa(s)} \kappa(y)/\mathfrak p)$. Nous omettons la vérification que les deux
constructions sont réciproques l'une de l'autre.
\end{proof}

\begin{lemma}
\label{lemma-fibre-product-immersion}
Soient  $f : X \to S$  et  $g : Y \to S$  des morphismes de schémas ayant la même cible.
\begin{enumerate}
\item Si  $f : X \to S$  est une immersion fermée, alors  $X \times_S Y \to Y$  est une immersion
fermée. De plus, si  $X \to S$  correspond au faisceau quasi-cohérent d'idéaux
 $\mathcal{I} \subset \mathcal{O}_S$, alors  $X \times_S Y \to Y$  correspond au faisceau d'idéaux  $\Im(g^*\mathcal{I} \to \mathcal{O}_Y)$.
\item Si  $f : X \to S$  est une immersion ouverte, alors  $X \times_S Y \to Y$  est une immersion
ouverte.
\item Si  $f : X \to S$  est une immersion, alors  $X \times_S Y \to Y$  est une immersion.
\end{enumerate}
\end{lemma}

\begin{proof}
Supposons que  $X \to S$  soit une immersion fermée correspondant au faisceau
quasi-cohérent d'idéaux  $\mathcal{I} \subset \mathcal{O}_S$. D'après le Lemme  \ref{lemma-restrict-map-to-closed} , le
sous-espace fermé  $Z \subset Y$  défini par le faisceau d'idéaux  $\Im(g^*\mathcal{I} \to \mathcal{O}_Y)$  est le
produit fibré dans la catégorie des espaces localement annelés. D'après le Lemme
 \ref{lemma-closed-subspace-scheme} ,  $Z$  est un schéma. Par conséquent,  $Z = X \times_S Y$  et la
première affirmation en découle. La deuxième découle par exemple du Lemme
 \ref{lemma-open-fibre-product} . La troisième est une combinaison des deux premières.
\end{proof}

\begin{definition}
\label{definition-inverse-image-closed-subscheme}
Soit $f : X \to Y$ un morphisme de schémas. Soit $Z \subset Y$ un
sous-schéma fermé de $Y$. L'{\it image
réciproque $f^{-1}(Z)$  du sous-schéma fermé  $Z$} est le
sous-schéma fermé $Z \times_Y X$  de  $X$. Voir le lemme \ref{lemma-fibre-product-immersion} ci-dessus.
\end{definition}

\noindent
Nous pouvons également utiliser occasionnellement cette terminologie pour les
sous-schémas localement fermés et ouverts.









\section{Changement de base en géométrie algébrique}
\label{section-base-change}

\noindent
Une motivation pour l'introduction du langage des schémas est qu'il donne une
notion très précise de ce que signifie définir une variété sur un corps
particulier. Par exemple, une variété  $X$  sur  $\mathbf{Q}$  est, par définition
(Variétés, Définition  \ref{varieties-definition-variety}), un schéma muni d'un morphisme  $X \to \Spec(\mathbf{Q})$  de type fini et séparé,
et qui est lui-même irréductible et réduit\footnote{Bien sûr, les géomètres algébriques
discutent encore pour savoir s'il faut exiger que  $X$  soit
géométriquement irréductible sur  $\mathbf{Q}$.}. Dans tous les cas,
l'idée est plus généralement de travailler avec des schémas sur un
{\it schéma de base} donné, souvent noté
 $S$. Nous dirons « soit  $X$ un schéma
sur $S$ » pour signifier simplement que $X$ est muni d'un
morphisme $X \to S$. Dans les diagrammes, nous représenterons le
 {\it morphisme structural}   $X \to S$ par une
flèche descendante de $X$ vers $S$. Nous nous intéressons
souvent davantage aux propriétés de $X$ relativement à $S$ qu'à
la géométrie interne de $X$. Par exemple, nous aimerions connaître des
choses sur les fibres de  $X \to S$, ce qui se passe pour  $X$  après
changement de base, et ainsi de suite.

\medskip\noindent
Nous introduisons une partie du langage couramment utilisé. Bien sûr, ce
langage n'est qu'un cas particulier de réflexion sur la catégorie des objets
au-dessus d'un objet donné dans une catégorie, voir Catégories, Exemple
 \ref{categories-example-category-over-X}.

\begin{definition}
\label{definition-base-change}
Soit  $S$  un schéma.
\begin{enumerate}
\item Nous disons que  $X$  est un  {\it schéma
sur $S$} si $X$ est muni d'un morphisme
de schémas $X \to S$. Le morphisme  $X \to S$  est parfois appelé le
 {\it morphisme structural}.
\item Si  $R$  est un anneau, nous disons que  $X$  est un schéma
 {\it  sur  $R$}  pour dire que
 $X$  est un schéma sur  $\Spec(R)$.
\item Un  {\it  morphisme  $f : X \to Y$  de schémas sur
 $S$}  est un morphisme de schémas tel que la composition
 $X \to Y \to S$  de  $f$  avec le morphisme structural de  $Y$  soit
égale au morphisme structural de  $X$.
\item Nous notons  $\Mor_S(X, Y)$  l'ensemble de tous les morphismes de  $X$  vers
 $Y$  sur  $S$.
\item Soit  $X$  un schéma sur  $S$. Soit  $S' \to S$  un morphisme de
schémas. Le  {\it  changement de base}  de
 $X$  est le schéma  $X_{S'} = S' \times_S X$  sur  $S'$.
\item Soit  $f : X \to Y$  un morphisme de schémas sur  $S$. Soit  $S' \to S$  un
morphisme de schémas. Le  {\it  changement de base} 
de  $f$  est le morphisme induit  $f' : X_{S'} \to Y_{S'}$  (à savoir le morphisme
 $\text{id}_{S'} \times_{\text{id}_S} f$).
\item Soit  $R$  un anneau. Soit  $X$  un schéma sur  $R$. Soit
 $R \to R'$  un morphisme d'anneaux. Le
 {\it changement de base}   $X_{R'}$  est le schéma  $\Spec(R') \times_{\Spec(R)} X$  sur
 $R'$.
\end{enumerate}
\end{definition}

\noindent
Voici un résultat typique.

\begin{lemma}
\label{lemma-base-change-immersion}
Soit  $S$  un schéma. Soit  $f : X \to Y$  une immersion (resp. \ 
immersion fermée, resp. immersion ouverte) de schémas sur  $S$. Alors tout
changement de base de  $f$  est une immersion (resp. \  immersion
fermée, resp. immersion ouverte).
\end{lemma}

\begin{proof}
On peut penser au changement de base de  $f$  via le morphisme  $S' \to S$ 
comme la flèche verticale en haut à gauche dans le diagramme commutatif suivant:
$$
\xymatrix{
X_{S'} \ar[r] \ar[d] & X \ar[d] \ar@/^4ex/[dd] \\
Y_{S'} \ar[r] \ar[d] & Y \ar[d] \\
S' \ar[r] & S
}
$$
Le diagramme implique  $X_{S'} \cong Y_{S'} \times_Y X$, et le lemme découle du Lemme  \ref{lemma-fibre-product-immersion}.
\end{proof}

\noindent
En fait, ce type de résultat est tellement typique qu'il existe une expression
pour l'indiquer. La voici.

\begin{definition}
\label{definition-preserved-by-base-change}
Propriétés et changement de base.
\begin{enumerate}
\item Soit  $\mathcal{P}$  une propriété des schémas sur une base. On dit que  $\mathcal{P}$ 
est  {\it  préservée par un changement de base
arbitraire}, ou simplement que  $\mathcal{P}$  est
 {\it  préservée par changement de base}  si, chaque
fois que  $X/S$  possède  $\mathcal{P}$, tout changement de base  $X_{S'}/S'$  possède
 $\mathcal{P}$.
\item Soit  $\mathcal{P}$  une propriété des morphismes de schémas sur une base. On dit que
 $\mathcal{P}$  est  {\it  préservée par un changement de base
arbitraire}, ou simplement qu'elle est {\it  préservée par
changement de base}  si, chaque fois que  $f : X \to Y$  sur  $S$  possède
 $\mathcal{P}$, tout changement de base  $f' : X_{S'} \to Y_{S'}$  sur  $S'$  possède
 $\mathcal{P}$.
\end{enumerate}
\end{definition}

\noindent
À ce stade, nous pouvons dire que la propriété « être une immersion fermée » est préservée par
un changement de base arbitraire.

\begin{definition}
\label{definition-fibre}
Soit $f : X \to S$ un morphisme de schémas. Soit  $s \in S$ un point. La
 {\it fibre schématique $X_s$  de
 $f$ au-dessus de $s$}, ou simplement la
 {\it fibre de $f$ au-dessus de $s$},
est le schéma qui s'insère dans le diagramme de produit fibré suivant
$$
\xymatrix{
X_s = \Spec(\kappa(s)) \times_S X \ar[r] \ar[d] &
X \ar[d] \\
\Spec(\kappa(s)) \ar[r] &
S
}
$$
Nous considérons toujours la fibre  $X_s$  comme un schéma sur
 $\kappa(s)$.
\end{definition}

\begin{lemma}
\label{lemma-fibre-topological}
Soit  $f : X \to S$  un morphisme de schémas. Considérons les diagrammes
$$
\vcenter{
\xymatrix{
X_s \ar[r] \ar[d] & X \ar[d] \\
\Spec(\kappa(s)) \ar[r] & S
}
}
\quad\text{et}\quad
\vcenter{
\xymatrix{
\Spec(\mathcal{O}_{S, s}) \times_S X \ar[r] \ar[d] & X \ar[d] \\
\Spec(\mathcal{O}_{S, s}) \ar[r] & S
}
}
$$
Dans les deux cas, ces diagrammes induisent des carrés fibrés d'espaces
topologiques et en particulier, la flèche horizontale supérieure est un
homéomorphisme sur son image.
\end{lemma}

\begin{proof}
Choisissez un ouvert affine  $U \subset S$  qui contient  $s$. Les morphismes
horizontaux inférieurs se factorisent par  $U$, voir le Lemme  \ref{lemma-morphism-from-spec-local-ring} 
par exemple. Ainsi, nous pouvons supposer que  $S$  est affine. Si
 $X$  est également affine, alors le résultat découle de l'Algèbre,
Remarque  \ref{algebra-remark-fundamental-diagram}. Dans le cas général, le résultat découle en recouvrant
 $X$  par des ouverts affines.
\end{proof}

\begin{lemma}
\label{lemma-local-ring-fibre}
Soit  $f : X \to S$  un morphisme de schémas. Soit  $x \in X$  avec image
 $s \in S$. En considérant  $x$  comme un point de  $X_s$  (voir le
Lemme  \ref{lemma-fibre-topological}), nous avons
$$
\mathcal{O}_{X_s, x} \cong
\mathcal{O}_{X, x}/\mathfrak m_s\mathcal{O}_{X, x} \cong
\mathcal{O}_{X, x} \otimes_{\mathcal{O}_{S, s}} \kappa(s)
$$
\end{lemma}

\begin{proof}
Comme dans la preuve du Lemme  \ref{lemma-fibre-topological} , cela se réduit au cas où  $X$ 
et  $S$  sont affines. Dans le cas affine, l'énoncé se traduit par Algèbre,
Remarque  \ref{algebra-remark-local-ring-fibre}.
\end{proof}







\section{Morphismes quasi-compacts}
\label{section-quasi-compact}

\noindent
Un schéma est  {\it  quasi-compact}  si son espace
topologique sous-jacent est quasi-compact. Il existe une notion relative qui est
définie comme suit.

\begin{definition}
\label{definition-quasi-compact}
Un morphisme de schémas est dit  {\it  quasi-compact} 
si l'application sous-jacente des espaces topologiques est quasi-compacte, voir
Topologie, Définition  \ref{topology-definition-quasi-compact}.
\end{definition}

\begin{lemma}
\label{lemma-quasi-compact-affine}
Soit  $f : X \to S$  un morphisme de schémas. Les énoncés suivants sont équivalents
\begin{enumerate}
\item $f : X \to S$  est quasi-compact,
\item l'image réciproque de chaque ouvert affine est quasi-compacte, et
\item il existe un recouvrement ouvert affine  $S = \bigcup_{i \in I} U_i$  tel que  $f^{-1}(U_i)$  soit
quasi-compact pour tout  $i$.
\end{enumerate}
\end{lemma}

\begin{proof}
Supposons qu'on nous donne un recouvrement  $S = \bigcup_{i \in I} U_i$  comme dans (3). Tout
d'abord, soit  $U \subset S$  un ouvert affine quelconque. Pour tout  $u \in U$ ,
nous pouvons trouver un indice  $i(u) \in I$  tel que  $u \in U_{i(u)}$. Comme les
ouverts standards forment une base pour la topologie sur  $U_{i(u)}$ , nous
pouvons trouver  $W_u \subset U \cap U_{i(u)}$  qui est un ouvert standard dans  $U_{i(u)}$. Par
quasi-compacité, nous pouvons trouver un nombre fini de points  $u_1, \ldots, u_n \in U$  tels que
 $U = \bigcup_{j = 1}^n W_{u_j}$. Pour chaque  $j$ , écrivons  $f^{-1}U_{i(u_j)} = \bigcup_{k \in K_j} V_{jk}$  comme une union
finie d'ouverts affines. Puisque  $W_{u_j} \subset U_{i(u_j)}$  est un ouvert standard, nous voyons
que  $f^{-1}(W_{u_j}) \cap V_{jk}$  est un ouvert standard de  $V_{jk}$, voir Algèbre, Lemme
 \ref{algebra-lemma-spec-functorial}. Ainsi  $f^{-1}(W_{u_j}) \cap V_{jk}$  est affine, et donc  $f^{-1}(W_{u_j})$  est une union
finie d'ouverts affines. Cela prouve que l'image réciproque de tout ouvert affine est
une union finie d'ouverts affines.

\medskip\noindent
Ensuite, supposons que l'image réciproque de tout ouvert affine soit une union
finie d'ouverts affines. Soit  $K \subset S$  un ouvert quasi-compact quelconque.
Comme  $S$  a une base de la topologie constituée d'ouverts affines, nous
voyons que  $K$  est une union finie d'ouverts affines. Par conséquent,
l'image réciproque de  $K$  est une union finie d'ouverts affines. Par
conséquent,  $f$  est quasi-compact.

\medskip\noindent
Enfin, supposons que  $f$  soit quasi-compact. Dans ce cas, l'argument du
paragraphe précédent montre que l'image réciproque de tout ouvert affine est une
union finie d'ouverts affines.
\end{proof}

\begin{lemma}
\label{lemma-quasi-compact-preserved-base-change}
Être quasi-compact est une propriété des morphismes de schémas sur une base qui
est préservée par tout changement de base.
\end{lemma}

\begin{proof}
Omis.
\end{proof}

\begin{lemma}
\label{lemma-composition-quasi-compact}
La composition de morphismes quasi-compacts est quasi-compacte.
\end{lemma}

\begin{proof}
Cela découle des définitions et de Topologie, Lemme  \ref{topology-lemma-composition-quasi-compact}.
\end{proof}

\begin{lemma}
\label{lemma-closed-immersion-quasi-compact}
Une immersion fermée est quasi-compacte.
\end{lemma}

\begin{proof}
Cela découle des définitions et de Topologie, Lemme  \ref{topology-lemma-closed-in-quasi-compact}.
\end{proof}

\begin{example}
\label{example-open-immersion-not-quasi-compact}
Une immersion ouverte n'est généralement pas quasi-compacte. L'exemple standard de
ceci est le sous-espace ouvert  $U \subset X$, où  $X = \Spec(k[x_1, x_2, x_3, \ldots])$, où  $U$  est
 $X \setminus \{0\}$, et où  $0$  est le point de  $X$  correspondant à
l'idéal maximal  $(x_1, x_2, x_3, \ldots)$.
\end{example}

\begin{lemma}
\label{lemma-image-quasi-compact-closed}
Soit  $f : X \to S$  un morphisme quasi-compact de schémas. Les affirmations
suivantes sont équivalentes
\begin{enumerate}
\item $f(X) \subset S$  est fermé, et
\item $f(X) \subset S$  est stable par spécialisation.
\end{enumerate}
\end{lemma}

\begin{proof}
Nous avons (1)  $\Rightarrow$  (2) par Topologie, Lemme  \ref{topology-lemma-open-closed-specialization}. Supposons (2).
Soit  $U \subset S$  un ouvert affine. Il suffit de prouver que  $f(X) \cap U$  est
fermé. Puisque  $U \cap f(X)$  est stable par spécialisations dans  $U$, nous
avons réduit au cas où  $S$  est affine. Comme  $f$  est
quasi-compact, nous en déduisons que  $X = f^{-1}(S)$  est quasi-compact puisque
 $S$  est affine. Ainsi nous pouvons écrire  $X = \bigcup_{i = 1}^n U_i$  avec  $U_i \subset X$ 
ouvert affine. Disons  $S = \Spec(R)$  et  $U_i = \Spec(A_i)$  pour une certaine
 $R$-algèbre  $A_i$. Alors  $f(X) = \Im(\Spec(A_1 \times \ldots \times A_n)
\to \Spec(R))$. Le résultat découle donc d'Algèbre, Lemme
 \ref{algebra-lemma-image-stable-specialization-closed}.
\end{proof}

\begin{lemma}
\label{lemma-quasi-compact-closed}
Soit  $f : X \to S$  un morphisme quasi-compact de schémas. Alors  $f$  est
fermé si et seulement si les spécialisations se relèvent le long de  $f$,
voir Topologie, Définition  \ref{topology-definition-lift-specializations}.
\end{lemma}

\begin{proof}
D'après Topologie, Lemme  \ref{topology-lemma-closed-open-map-specialization} , si  $f$  est fermé, alors les
spécialisations se relèvent le long de  $f$. Inversement, supposons que
les spécialisations se relèvent le long de  $f$. Soit  $Z \subset X$  un
sous-ensemble fermé. Nous pouvons considérer  $Z$  comme un schéma avec la
structure de schéma réduite induite, voir Définition  \ref{definition-reduced-induced-scheme}. Puisque
 $Z \subset X$  est fermé, la restriction de  $f$  à  $Z$  est toujours
quasi-compacte. De plus, les spécialisations se relèvent également le long de
 $Z \to S$ , voir Topologie, Lemme  \ref{topology-lemma-lift-specialization-composition}. Par conséquent, il suffit de
prouver que  $f(X)$  est fermé si les spécialisations se relèvent le long de
 $f$. En particulier,  $f(X)$  est stable par spécialisation, voir
Topologie, Lemme  \ref{topology-lemma-lift-specializations-images}. Ainsi  $f(X)$  est fermé par le Lemme
 \ref{lemma-image-quasi-compact-closed}.
\end{proof}










\section{Critère valuatif de fermeture universelle}
\label{section-valuative-criterion-universal-closedness}

\noindent
En topologie, la section  \ref{topology-section-proper}  étudie les applications propres comme des
applications fermées entre espaces topologiques dont toutes les fibres sont
quasi-compactes, ou comme des applications telles que tous les changements de base
sont des applications fermées. Voici la notion correspondante en géométrie
algébrique.

\begin{definition}
\label{definition-universally-closed}
Un morphisme de schémas  $f : X \to S$  est dit  {\it 
universellement fermé}  si chaque changement de base  $f' : X_{S'} \to S'$  est
fermé.
\end{definition}

\noindent
En fait, l’adjectif « universel » est souvent utilisé de cette manière. En
d’autres termes, étant donné une propriété  $\mathcal{P}$  des morphismes, on dit que
« $X \to S$  est universellement  $\mathcal{P}$ » si et seulement si chaque
changement de base  $X_{S'} \to S'$  possède  $\mathcal{P}$.

\medskip\noindent
Veuillez consulter Morphismes, Section  \ref{morphisms-section-proper}  pour une discussion plus
détaillée des propriétés des morphismes universellement fermés. Dans cette
section, nous limitons la discussion à la relation entre les morphismes
universellement fermés et les morphismes satisfaisant la partie existence du
critère valuatif.

\begin{lemma}
\label{lemma-specializations-lift}
Soit  $f : X \to S$  un morphisme de schémas.
\begin{enumerate}
\item Si  $f$  est universellement fermé, alors les spécialisations se relèvent
le long de tout changement de base de  $f$, voir Topologie, Définition
 \ref{topology-definition-lift-specializations}.
\item Si  $f$  est quasi-compact et que les spécialisations se relèvent le long
de tout changement de base de  $f$, alors  $f$  est universellement
fermé.
\end{enumerate}
\end{lemma}

\begin{proof}
La partie (1) est une conséquence directe de Topologie, Lemme  \ref{topology-lemma-closed-open-map-specialization}. La
partie (2) découle des lemmes  \ref{lemma-quasi-compact-closed}  et  \ref{lemma-quasi-compact-preserved-base-change}.
\end{proof}

\begin{definition}
\label{definition-valuative-criterion}
Soit  $f : X \to S$  un morphisme de schémas. Nous disons que  $f$ 
 {\it  satisfait la partie existence du critère
valuatif}  si, étant donné tout diagramme commutatif dont les flèches pleines sont données
$$
\xymatrix{
\Spec(K) \ar[r] \ar[d] & X \ar[d] \\
\Spec(A) \ar[r] \ar@{-->}[ru] & S
}
$$
où  $A$  est un anneau de valuation avec corps des fractions  $K$,
la flèche pointillée existe. Nous disons que  $f$   {\it 
satisfait la partie unicité du critère valuatif}  s'il existe au plus
une flèche pointillée pour tout diagramme comme ci-dessus (sans exiger l'existence
bien sûr).
\end{definition}

\noindent
Un {\it anneau de valuation} est un anneau local intègre
maximal pour la relation de domination dans son corps des fractions, voir Algèbre,
Définition \ref{algebra-definition-valuation-ring}. Par conséquent, le spectre d'un anneau de valuation
possède un point générique unique $\eta$  et un point fermé
unique $0$, et bien sûr nous avons la spécialisation $\eta \leadsto 0$.
L'importance des anneaux de valuation réside dans le fait que toute spécialisation
de points dans tout schéma est l'image de $\eta \leadsto 0$ sous un certain morphisme du
spectre d'un certain anneau de valuation. Voici le résultat précis.

\begin{lemma}
\label{lemma-points-specialize}
\begin{slogan}
Les spécialisations se réalisent au moyen d'anneaux de valuation.
\end{slogan}
Soit $S$ un schéma. Soit $s' \leadsto s$ une spécialisation de
points de $S$. Alors
\begin{enumerate}
\item il existe un anneau de valuation  $A$  et un morphisme  $f : \Spec(A) \to S$  tels
que le point générique  $\eta$  de  $\Spec(A)$  ait pour image  $s'$  et le
point fermé pour image  $s$, et
\item étant donnée une extension de corps  $K/\kappa(s')$, on peut faire en sorte que
l'extension  $\kappa(\eta)/\kappa(s')$  induite par  $f$  soit isomorphe à l'extension
donnée.
\end{enumerate}
\end{lemma}

\begin{proof}
Soit  $s' \leadsto s$  une spécialisation dans  $S$, et soit  $K/\kappa(s')$  une
extension de corps. D'après le Lemme  \ref{lemma-specialize-points}  et la discussion suivant le
Lemme  \ref{lemma-characterize-points} , cela conduit à des homomorphismes d'anneaux  $\mathcal{O}_{S, s} \to \kappa(s') \to K$. Soit
 $A \subset K$  un quelconque anneau de valuation dont le corps des fractions est
 $K$  et qui domine l'image de  $\mathcal{O}_{S, s} \to K$, voir Algèbre, Lemme
 \ref{algebra-lemma-dominate}. L'homomorphisme d'anneaux  $\mathcal{O}_{S, s} \to A$  induit le morphisme
 $f : \Spec(A) \to S$, voir Lemme  \ref{lemma-morphism-from-spec-local-ring}. Ce morphisme possède toutes les propriétés
souhaitées par construction.
\end{proof}

\begin{lemma}
\label{lemma-lift-specializations-valuative}
Soit  $f : X \to S$  un morphisme de schémas. Les énoncés suivants sont équivalents:
\begin{enumerate}
\item Les spécialisations se relèvent par tout changement de base de  $f$.
\item Le morphisme  $f$  satisfait la partie existence du critère valuatif.
\end{enumerate}
\end{lemma}

\begin{proof}
Supposons que (1) soit vrai. Considérons un diagramme dont les flèches pleines sont données, comme dans la Définition
 \ref{definition-valuative-criterion}. Pour trouver la flèche en pointillés, nous pouvons remplacer
 $X \to S$  par  $X_{\Spec(A)} \to \Spec(A)$  puisque, après tout, l’hypothèse est stable par
changement de base. Ainsi, nous pouvons supposer  $S = \Spec(A)$. Soit  $x' \in X$ 
l’image de  $\Spec(K) \to X$, de sorte que nous avons  $\kappa(x') \subset K$, voir le Lemme
 \ref{lemma-characterize-points}. D’après l’hypothèse, il existe une spécialisation  $x' \leadsto x$  dans
 $X$  telle que  $x$  ait pour image le point fermé de
 $S = \Spec(A)$. Nous obtenons un morphisme d’anneaux locaux  $A \to \mathcal{O}_{X, x}$  et un
morphisme d’anneaux  $\mathcal{O}_{X, x} \to \kappa(x')$, voir le Lemme  \ref{lemma-specialize-points}  et la discussion
suivant le Lemme  \ref{lemma-characterize-points}. La composition  $A \to \mathcal{O}_{X, x} \to \kappa(x') \to K$  est l'injection donnée
 $A \to K$. Comme  $A \to \mathcal{O}_{X, x}$  est local, l'image de  $\mathcal{O}_{X, x} \to K$  domine
 $A$  et est donc égale à  $A$, d'après Algèbre, Définition
 \ref{algebra-definition-valuation-ring}. Ainsi, nous obtenons un homomorphisme d'anneaux  $\mathcal{O}_{X, x} \to A$  et donc
un morphisme  $\Spec(A) \to X$  (voir le Lemme  \ref{lemma-morphism-from-spec-local-ring}  et la discussion qui suit).
Cela prouve (2).

\medskip\noindent
En revanche, supposons que (2) soit vérifié. Il est immédiat que la partie
existence du critère valuatif est vérifiée pour tout changement de base
 $X_{S'} \to S'$  de  $f$  en considérant le diagramme commutatif suivant
$$
\xymatrix{
\Spec(K) \ar[r] \ar[d] & X_{S'} \ar[r] \ar[d] & X \ar[d] \\
\Spec(A) \ar[r] \ar@{-->}[ru] \ar@{-->}[rru] & S' \ar[r] & S
}
$$
En d'autres termes, la flèche pointillée plus horizontale mènera à l'autre par
définition du produit fibré. Il suffit donc de montrer que
les spécialisations se relèvent le long de  $f$. Soit  $s' \leadsto s$  une
spécialisation dans  $S$, et soit  $x' \in X$  un point situé au-dessus de
 $s'$. Appliquons le Lemme  \ref{lemma-points-specialize}  à  $s' \leadsto s$  et à l'extension de
corps  $K = \kappa(x')/\kappa(s')$. Nous obtenons un diagramme commutatif
$$
\xymatrix{
\Spec(K) \ar[rr] \ar[d] & & X \ar[d] \\
\Spec(A) \ar[r] \ar@{-->}[rru] &
\Spec(\mathcal{O}_{S, s}) \ar[r] & S
}
$$
et, selon la condition (2), nous obtenons la flèche pointillée. L'image
 $x$  du point fermé de  $\Spec(A)$  dans  $X$  sera une solution à
notre problème, c'est-à-dire que  $x$  est une spécialisation de
 $x'$  et a pour image  $s$.
\end{proof}

\begin{proposition}[ Critère valuatif de fermeture universelle]
\label{proposition-characterize-universally-closed}
Soit  $f$  un morphisme quasi-compact de schémas. Alors  $f$  est
universellement fermé si et seulement si  $f$  satisfait la partie
existence du critère valuatif.
\end{proposition}

\begin{proof}
Ceci est une conséquence formelle des lemmes  \ref{lemma-specializations-lift}  et  \ref{lemma-lift-specializations-valuative} 
ci-dessus.
\end{proof}

\begin{example}
\label{example-projective-line-universally-closed}
Soit  $k$  un corps. Considérons le morphisme structural  $p : \mathbf{P}^1_k \to \Spec(k)$  de
la droite projective sur  $k$, voir l'exemple  \ref{example-projective-line}. Utilisons le
critère valuatif ci-dessus pour prouver que  $p$  est universellement
fermé. Par construction,  $\mathbf{P}^1_k$  est recouvert par deux ouverts affines et
donc  $p$  est quasi-compact. Supposons qu'un diagramme commutatif
$$
\xymatrix{
\Spec(K) \ar[r]_\xi \ar[d] & \mathbf{P}^1_k \ar[d] \\
\Spec(A) \ar[r]^\varphi & \Spec(k)
}
$$
soit donné, où  $A$  est un anneau de valuation et  $K$  est son
corps des fractions. Rappelons que  $\mathbf{P}^1_k$  est obtenu en recollant  $\Spec(k[x])$ 
à  $\Spec(k[y])$  en identifiant  $D(x)$  à  $D(y)$  via  $x = y^{-1}$  (ou plus
symétriquement  $xy = 1$). Pour montrer qu'il existe un morphisme  $\Spec(A) \to \mathbf{P}^1_k$ 
s'insérant en diagonale dans le diagramme ci-dessus, nous pouvons supposer que
 $\xi$  ait son image dans l'ouvert  $\Spec(k[x])$  (par symétrie). Cela donne le
diagramme commutatif d'anneaux suivant
$$
\xymatrix{
K & k[x] \ar[l]^{\xi^\sharp} \\
A \ar[u] & k \ar[u] \ar[l]_{\varphi^\sharp}
}
$$
Par Algèbre, Lemme  \ref{algebra-lemma-valuation-ring-x-or-x-inverse} , nous voyons que soit  $\xi^\sharp(x) \in A$ , soit
 $\xi^\sharp(x)^{-1} \in A$. Dans le premier cas, nous obtenons un morphisme d'anneaux
$$
k[x] \to A,
\ \lambda \mapsto \varphi^\sharp(\lambda),
\  x \mapsto \xi^\sharp(x)
$$
s'insérant dans le diagramme d'anneaux ci-dessus, et nous réussissons. Dans le
deuxième cas, nous voyons que nous obtenons un morphisme d'anneaux
$$
k[y] \to A,
\ \lambda \mapsto \varphi^\sharp(\lambda),
\ y \mapsto \xi^\sharp(x)^{-1}.
$$
Cela donne un morphisme  $\Spec(A) \to \Spec(k[y]) \to \mathbf{P}^1_k$  qui s'insère en diagonale dans le diagramme
commutatif initial de cet exemple (vérification omise).
\end{example}





























\section{Axiomes de séparation}
\label{section-separation-axioms}

\noindent
Un espace topologique  $X$  est de Hausdorff si et seulement si la
diagonale  $\Delta \subset X \times X$  est un sous-ensemble fermé. L'analogue en géométrie
algébrique est, étant donné un schéma  $X$  sur un schéma de base
 $S$, de considérer le morphisme diagonal
$$
\Delta_{X/S} : X \longrightarrow X \times_S X.
$$
C'est le morphisme unique de schémas tel que  $\text{pr}_1 \circ \Delta_{X/S} = \text{id}_X$  et  $\text{pr}_2 \circ \Delta_{X/S} = \text{id}_X$  (il
existe dans toute catégorie avec produits fibrés).

\begin{lemma}
\label{lemma-diagonal-affines-closed}
Le morphisme diagonal d'un morphisme entre affines est fermé.
\end{lemma}

\begin{proof}
Le morphisme diagonal associé au morphisme  $\Spec(S) \to \Spec(R)$  est le morphisme entre
spectres correspondant à l'homomorphisme d'anneaux  $S \otimes_R S \to S$,  $a \otimes b \mapsto ab$.
Cette application est clairement surjective, donc  $S \cong S \otimes_R S/J$  pour un certain
idéal  $J \subset S \otimes_R S$. Par conséquent,  $\Delta$  est une immersion fermée selon
l'exemple  \ref{example-closed-immersion-affines}.
\end{proof}

\begin{lemma}
\label{lemma-diagonal-immersion}
\begin{slogan}
Le morphisme diagonal pour les schémas relatifs est une immersion.
\end{slogan}
Soit  $X$  un schéma sur  $S$. Le morphisme diagonal  $\Delta_{X/S}$ 
est une immersion.
\end{lemma}

\begin{proof}
Rappelons que si  $V \subset X$  est ouvert affine et se projette dans  $U \subset S$ 
ouvert affine, alors  $V \times_U V$  est ouvert affine dans  $X \times_S X$, voir les
lemmes  \ref{lemma-fibre-product-affines}  et  \ref{lemma-open-fibre-product}. Considérons le sous-schéma ouvert
 $W$  de  $X \times_S X$  qui est l'union de ces ouverts affines  $V \times_U V$.
D'après le lemme  \ref{lemma-closed-local-target} , il suffit de montrer que chaque morphisme
 $\Delta_{X/S}^{-1}(V \times_U V) \to V \times_U V$  est une immersion fermée. Comme  $V = \Delta_{X/S}^{-1}(V \times_U V)$ , nous vérifions
simplement que  $\Delta_{V/U}$  est une immersion fermée, ce qui est le lemme
 \ref{lemma-diagonal-affines-closed}.
\end{proof}

\begin{definition}
\label{definition-separated}
Soit  $f : X \to S$  un morphisme de schémas.
\begin{enumerate}
\item Nous disons que  $f$  est  {\it  séparé}  si le
morphisme diagonal  $\Delta_{X/S}$  est une immersion fermée.
\item On dit que  $f$  est  {\it  quasi-séparé}  si
le morphisme diagonal  $\Delta_{X/S}$  est un morphisme quasi-compact.
\item On dit qu'un schéma  $Y$  est  {\it  séparé} 
si le morphisme  $Y \to \Spec(\mathbf{Z})$  est séparé.
\item On dit qu'un schéma  $Y$  est  {\it 
quasi-séparé}  si le morphisme  $Y \to \Spec(\mathbf{Z})$  est quasi-séparé.
\end{enumerate}
\end{definition}

\noindent
D'après les lemmes  \ref{lemma-diagonal-immersion}  et  \ref{lemma-immersion-when-closed} , nous voyons que  $\Delta_{X/S}$  est
une immersion fermée si et seulement si  $\Delta_{X/S}(X) \subset X \times_S X$  est un sous-ensemble fermé.
De plus, d'après le lemme  \ref{lemma-closed-immersion-quasi-compact} , nous voyons qu'un morphisme séparé est
quasi-séparé. La raison d'introduire les morphismes quasi-séparés est que les
morphismes non séparés apparaissent naturellement dans l'étude des variétés
algébriques (en particulier lorsqu'on étudie les problèmes de modules, les champs
algébriques, etc.). Mais le plus souvent, ils restent tout de même quasi-séparés.

\begin{example}
\label{example-not-quasi-separated}
Voici un exemple de morphisme non quasi-séparé. Supposons  $X = X_1 \cup X_2 \to S = \Spec(k)$  avec
 $X_1 = X_2 = \Spec(k[t_1, t_2, t_3, \ldots])$  recollés le long du complément de  $\{0\} = \{(t_1, t_2, t_3, \ldots)\}$  (recollés comme dans
l'exemple  \ref{example-affine-space-zero-doubled}). Dans ce cas, l'image inverse du schéma affine
 $X_1 \times_S X_2$  sous  $\Delta_{X/S}$  est le schéma  $\Spec(k[t_1, t_2, t_3, \ldots]) \setminus \{0\}$  qui n'est pas
quasi-compact.
\end{example}

\begin{lemma}
\label{lemma-where-are-they-equal}
Soient  $X$,  $Y$  des schémas sur  $S$. Soient
 $a, b : X \to Y$  des morphismes de schémas sur  $S$. Il existe un plus grand
sous-schéma localement fermé  $Z \subset X$  tel que  $a|_Z = b|_Z$. En fait,
 $Z$  est l'égalisateur de  $(a, b)$. De plus, si  $Y$  est séparé
sur  $S$, alors  $Z$  est un sous-schéma fermé.
\end{lemma}

\begin{proof}
L'égalisateur de  $(a, b)$  est pour des raisons catégoriques le produit fibré
 $Z$  dans le diagramme suivant.
$$
\xymatrix{
Z = Y \times_{(Y \times_S Y)} X \ar[r] \ar[d] &
 X \ar[d]^{(a , b)} \\
Y \ar[r]^-{\Delta_{Y/S}} & Y \times_S Y
}
$$
Ainsi, le lemme découle des lemmes  \ref{lemma-base-change-immersion},  \ref{lemma-diagonal-immersion}  et de la Définition
 \ref{definition-separated}.
\end{proof}

\begin{lemma}
\label{lemma-characterize-quasi-separated}
Soit  $f : X \to S$  un morphisme de schémas. Les affirmations suivantes sont
équivalentes:
\begin{enumerate}
\item Le morphisme  $f$  est quasi-séparé.
\item Pour chaque paire d'ouverts affines  $U, V \subset X$  qui se projettent dans un ouvert
affine commun de  $S$ , l'intersection  $U \cap V$  est une union finie
d'ouverts affines de  $X$.
\item Il existe un recouvrement par ouverts affines  $S = \bigcup_{i \in I} U_i$  et pour chaque
 $i$  un recouvrement par ouverts affines  $f^{-1}U_i = \bigcup_{j \in I_i} V_j$  tel que pour chaque
 $i$  et chaque paire  $j, j' \in I_i$  l'intersection  $V_j \cap V_{j'}$  soit une
union finie d'ouverts affines de  $X$.
\end{enumerate}
\end{lemma}

\begin{proof}
Prouvons que (3) implique (1). D'après le Lemme  \ref{lemma-affine-covering-fibre-product} , le recouvrement
 $X \times_S X = \bigcup_i \bigcup_{j, j'} V_j \times_{U_i} V_{j'}$  est un recouvrement par ouverts affines de  $X \times_S X$. De plus,
 $\Delta_{X/S}^{-1}(V_j \times_{U_i} V_{j'}) = V_j \cap V_{j'}$. Ainsi, l'implication découle du Lemme  \ref{lemma-quasi-compact-affine}.

\medskip\noindent
L'implication (1)  $\Rightarrow$  (2) découle du fait que, sous les hypothèses de
(2), le produit fibré  $U \times_S V$  est un ouvert affine de  $X \times_S X$.
L'implication (2)  $\Rightarrow$  (3) est triviale.
\end{proof}

\begin{lemma}
\label{lemma-characterize-separated}
Soit  $f : X \to S$  un morphisme de schémas.
\begin{enumerate}
\item Si  $f$  est séparé alors pour chaque paire d'ouverts affines  $(U, V)$ 
de  $X$  qui se projettent dans un ouvert affine commun de  $S$ 
nous avons
\begin{enumerate}
\item l'intersection  $U \cap V$  est affine.
\item l'homomorphisme d'anneaux  $\mathcal{O}_X(U) \otimes_{\mathbf{Z}} \mathcal{O}_X(V)
\to \mathcal{O}_X(U \cap V)$  est surjectif.
\end{enumerate}
\item Si toute paire de points  $x_1, x_2 \in X$  au-dessus d'un point commun  $s \in S$  est
contenue dans des ouverts affines  $x_1 \in U$,  $x_2 \in V$  qui se projettent
dans un ouvert affine commun de  $S$  tels que (a), (b) sont vérifiés,
alors  $f$  est séparé.
\end{enumerate}
\end{lemma}

\begin{proof}
Supposons  $f$  séparé. Supposons que  $(U, V)$  soit une paire comme
dans (1). Soit  $W = \Spec(R)$  un ouvert affine de  $S$  contenant à la fois
 $f(U)$  et  $f(V)$. Notons  $U = \Spec(A)$  et  $V = \Spec(B)$  pour les
 $R$-algèbres  $A$  et  $B$. D'après le Lemme  \ref{lemma-open-fibre-product} ,
on voit que  $U \times_S V = U \times_W V = \Spec(A \otimes_R B)$  est un ouvert affine de  $X \times_S X$. Par conséquent,
d'après le Lemme  \ref{lemma-closed-subspace-scheme} , on voit que  $\Delta^{-1}(U \times_S V) \to U \times_S V$  peut être identifié à
 $\Spec((A \otimes_R B)/J) \to \Spec(A \otimes_R B)$  pour un certain idéal  $J \subset A \otimes_R B$. Ainsi  $U \cap V = \Delta^{-1}(U \times_S V)$  est affine.
L'assertion (1)(b) est vérifiée car  $A \otimes_{\mathbf{Z}} B \to (A \otimes_R B)/J$  est surjective.

\medskip\noindent
Supposons que l'hypothèse formulée en (2) soit vérifiée. Il est clair que la
collection d'ouverts affines  $U \times_S V$  pour les paires  $(U, V)$  comme en
(2) forme un recouvrement affine ouvert de  $X \times_S X$  (voir par exemple
\ Lemme  \ref{lemma-affine-covering-fibre-product}). Il suffit donc de montrer que chaque morphisme
 $U \cap V = \Delta_{X/S}^{-1}(U \times_S V) \to U \times_S V$  est une immersion fermée, voir le Lemme  \ref{lemma-closed-local-target}. D'après
l'hypothèse (a), nous avons  $U \cap V = \Spec(C)$  pour un certain anneau  $C$.
Après avoir choisi un ouvert affine  $W = \Spec(R)$  de  $S$  dans lequel
les images de  $U$  et de  $V$  sont contenues et en écrivant  $U = \Spec(A)$,  $V = \Spec(B)$ ,
nous voyons que l'hypothèse (b) signifie que la composition
$$
A \otimes_{\mathbf{Z}} B \to A \otimes_R B \to C
$$
est surjective. Ainsi  $A \otimes_R B \to C$  est surjective et nous concluons que
 $\Spec(C) \to \Spec(A \otimes_R B)$  est une immersion fermée.
\end{proof}

\begin{example}
\label{example-projective-line-separated}
Soit  $k$  un corps. Considérons le morphisme structural  $p : \mathbf{P}^1_k \to \Spec(k)$  de
la droite projective sur  $k$, voir l'exemple  \ref{example-projective-line}. Utilisons le
lemme ci-dessus pour prouver que  $p$  est séparé. Par construction,
 $\mathbf{P}^1_k$  est recouvert par deux ouverts affines  $U = \Spec(k[x])$  et  $V = \Spec(k[y])$ 
avec intersection  $U \cap V = \Spec(k[x, y]/(xy - 1))$  (en utilisant la notation évidente). Il suffit donc
de vérifier que les conditions (2)(a) et (2)(b) du Lemme  \ref{lemma-characterize-separated}  sont
vérifiées pour les paires d'ouverts affines  $(U, U)$,  $(U, V)$,
 $(V, U)$  et  $(V, V)$. Pour les paires  $(U, U)$  et  $(V, V)$ , c'est
trivial. Pour la paire  $(U, V)$ , cela revient à prouver que  $U \cap V$  est
affine, ce qui est vrai, et que l'homomorphisme d'anneaux
$$
k[x] \otimes_{\mathbf{Z}} k[y] \longrightarrow k[x, y]/(xy - 1)
$$
est surjectif. C'est évident car tout élément du côté droit peut être écrit comme
la somme d'un polynôme en  $x$  et d'un polynôme en  $y$.
\end{example}

\begin{lemma}
\label{lemma-fibre-product-after-map}
Soient  $f : X \to T$  et  $g : Y \to T$  des morphismes de schémas ayant la même cible.
Soit  $h : T \to S$  un morphisme de schémas. Alors le morphisme induit  $i : X \times_T Y \to X \times_S Y$ 
est une immersion. Si  $T \to S$  est séparé, alors  $i$  est une
immersion fermée. Si  $T \to S$  est quasi-séparé, alors  $i$  est un
morphisme quasi-compact.
\end{lemma}

\begin{proof}
Par la théorie générale des catégories, le diagramme suivant
$$
\xymatrix{
X \times_T Y \ar[r] \ar[d] & X \times_S Y \ar[d] \\
T \ar[r]^{\Delta_{T/S}} \ar[r] & T \times_S T
}
$$
est un diagramme de produit fibré. Le lemme découle des lemmes  \ref{lemma-diagonal-immersion},
 \ref{lemma-fibre-product-immersion}  et  \ref{lemma-quasi-compact-preserved-base-change}.
\end{proof}

\begin{lemma}
\label{lemma-semi-diagonal}
Soit  $g : X \to Y$  un morphisme de schémas sur  $S$. Le morphisme
 $i : X \to X \times_S Y$  est une immersion. Si  $Y$  est séparé sur  $S$ , c'est
une immersion fermée. Si  $Y$  est quasi-séparé sur  $S$ , il est
quasi-compact.
\end{lemma}

\begin{proof}
Il s'agit d'un cas particulier du lemme  \ref{lemma-fibre-product-after-map}  appliqué au morphisme
 $X = X \times_Y Y \to X \times_S Y$.
\end{proof}

\begin{lemma}
\label{lemma-section-immersion}
Soit  $f : X \to S$  un morphisme de schémas. Soit  $s : S \to X$  une section de
 $f$  (dans une formule  $f \circ s = \text{id}_S$). Alors  $s$  est une
immersion. Si  $f$  est séparé, alors  $s$  est une immersion
fermée. Si  $f$  est quasi-séparé, alors  $s$  est quasi-compact.
\end{lemma}

\begin{proof}
Il s'agit d'un cas particulier du Lemme  \ref{lemma-semi-diagonal}  appliqué à  $g =s$  donc
le morphisme  $i = s : S \to S \times_S X$.
\end{proof}

\begin{lemma}
\label{lemma-separated-permanence}
Propriétés de permanence.
\begin{enumerate}
\item Une composition de morphismes séparés est séparée.
\item Une composition de morphismes quasi-séparés est quasi-séparée.
\item Le changement de base d'un morphisme séparé est séparé.
\item Le changement de base d'un morphisme quasi-séparé est quasi-séparé.
\item Un produit (fibré) de morphismes séparés est séparé.
\item Un produit (fibré) de morphismes quasi-séparés est quasi-séparé.
\end{enumerate}
\end{lemma}

\begin{proof}
Soient  $X \to Y \to Z$  des morphismes. Supposons que  $X \to Y$  et  $Y \to Z$ 
soient séparés. La composition
$$
X \to X \times_Y X \to X \times_Z X
$$
est fermée parce que le premier morphisme l'est par hypothèse et le second d'après le Lemme
 \ref{lemma-fibre-product-after-map}. Le même argument fonctionne pour « quasi-séparé » (avec les mêmes
références).

\medskip\noindent
Soit  $f : X \to Y$  un morphisme de schémas sur une base  $S$. Soit
 $S' \to S$  un morphisme de schémas. Soit  $f' : X_{S'} \to Y_{S'}$  le changement de base de
 $f$. Alors le morphisme diagonal de  $f'$  est un morphisme
$$
\Delta_{f'} :
X_{S'} = S' \times_S X
\longrightarrow
X_{S'} \times_{Y_{S'}} X_{S'} = S' \times _S (X \times_Y X)
$$
qui est facilement vu comme étant le changement de base de  $\Delta_f$. Ainsi (3)
et (4) découlent du fait que les immersions fermées et les morphismes
quasi-compacts sont préservés par tout changement de base (Lemmes  \ref{lemma-fibre-product-immersion}  et
 \ref{lemma-quasi-compact-preserved-base-change}).

\medskip\noindent
Si  $f : X \to Y$  et  $g : U \to V$  sont des morphismes de schémas sur une base
 $S$, alors  $f \times g$  est la composition de  $X \times_S U \to X \times_S V$  (un changement
de base de  $g$) et  $X \times_S V \to Y \times_S V$  (un changement de base de  $f$).
Ainsi, (5) et (6) suivent de (1)  --  (4).
\end{proof}

\begin{lemma}
\label{lemma-compose-after-separated}
\begin{slogan}
Les morphismes séparés et quasi-séparés satisfont la propriété de simplification.
\end{slogan}
Soient  $f : X \to Y$  et  $g : Y \to Z$  des morphismes de schémas. Si  $g \circ f$  est
séparé, alors  $f$ l'est aussi. Si  $g \circ f$  est quasi-séparé, alors
 $f$ l'est également.
\end{lemma}

\begin{proof}
Supposons que  $g \circ f$  soit séparé. Considérons la factorisation  $X \to X \times_Y X \to X \times_Z X$ 
du morphisme diagonal de  $g \circ f$. D'après le Lemme  \ref{lemma-fibre-product-after-map} , le dernier
morphisme est une immersion. D'après l'hypothèse, l'image de  $X$  dans
 $X \times_Z X$  est fermée. Par conséquent, elle est aussi fermée dans  $X \times_Y X$.
Ainsi, nous voyons que  $X \to X \times_Y X$  est une immersion fermée selon le Lemme
 \ref{lemma-immersion-when-closed}.

\medskip\noindent
Supposons que  $g \circ f$  soit quasi-séparé. Soit  $V \subset Y$  un ouvert affine
dont l'image est contenue dans un ouvert affine de  $Z$. Soient  $U_1, U_2 \subset X$  des
ouverts affines dont les images sont contenues dans  $V$. Alors  $U_1 \cap U_2$  est une
union finie d'ouverts affines car  $U_1, U_2$  ont leurs images dans un ouvert affine
commun de  $Z$. Puisque nous pouvons recouvrir  $Y$  par des
ouverts affines comme  $V$ , nous déduisons le lemme du Lemme  \ref{lemma-characterize-quasi-separated}.
\end{proof}

\begin{lemma}
\label{lemma-quasi-compact-permanence}
Soient  $f : X \to Y$  et  $g : Y \to Z$  des morphismes de schémas. Si  $g \circ f$  est
quasi-compact et  $g$  est quasi-séparé, alors  $f$  est
quasi-compact.
\end{lemma}

\begin{proof}
Ceci est vrai parce que  $f$  est égal à la composition  $(1, f) : X \to X \times_Z Y \to Y$. La
première application est quasi-compacte d'après le Lemme  \ref{lemma-section-immersion}  car elle est
une section du morphisme quasi-séparé  $X \times_Z Y \to X$  (un changement de base de
 $g$, voir le Lemme  \ref{lemma-separated-permanence}). La deuxième application est
quasi-compacte car elle est le changement de base de  $g \circ f$, voir le Lemme
 \ref{lemma-quasi-compact-preserved-base-change}. Et les compositions de morphismes quasi-compacts sont
quasi-compactes, voir le Lemme  \ref{lemma-composition-quasi-compact}.
\end{proof}

\begin{lemma}
\label{lemma-affine-separated}
Un schéma affine est séparé. Un morphisme d'un schéma affine vers un autre schéma
est séparé.
\end{lemma}

\begin{proof}
Soit  $U = \Spec(A)$  un schéma affine. Alors  $U \to \Spec(\mathbf{Z})$  a une diagonale fermée
d'après le Lemme  \ref{lemma-diagonal-affines-closed}. Ainsi  $U$  est séparé selon la Définition
 \ref{definition-separated}. Si  $U \to X$  est un morphisme de schémas, alors nous pouvons
appliquer le Lemme  \ref{lemma-compose-after-separated}  aux morphismes  $U \to X \to \Spec(\mathbf{Z})$  pour conclure que
 $U \to X$  est séparé.
\end{proof}

\noindent
Vous vous êtes peut-être demandé si la condition de ne considérer que des paires
d'ouverts affines dont l'image est contenue dans un ouvert affine est vraiment
nécessaire pour pouvoir conclure que leur intersection est affine. Souvent, ce
n'est pas le cas!

\begin{lemma}
\label{lemma-curiosity}
Soit  $f : X \to S$  un morphisme. Supposons que  $f$  soit séparé et que
 $S$  soit un schéma séparé. Supposons que  $U \subset X$  et  $V \subset X$ 
soient des ouverts affines. Alors  $U \cap V$  est affine (et un sous-schéma fermé
de  $U \times V$).
\end{lemma}

\begin{proof}
Dans ce cas,  $X$  est séparé par le Lemme  \ref{lemma-separated-permanence}. Ainsi,
 $U \cap V$  est affine en appliquant le Lemme  \ref{lemma-characterize-separated}  au morphisme
 $X \to \Spec(\mathbf{Z})$.
\end{proof}

\noindent
D'autre part, l'exemple suivant montre que nous ne pouvons pas nous attendre à ce
que l'image d'un ouvert affine soit contenue dans un ouvert affine.

\begin{example}
\label{example-image-affine-projective}
Considérons le schéma non affine  $U = \Spec(k[x, y]) \setminus \{(x, y)\}$  de l'Exemple  \ref{example-not-affine}. D'autre
part, considérons le schéma
$$
\mathbf{GL}_{2, k} = \Spec(k[a, b, c, d, 1/ad - bc]).
$$
Il existe un morphisme  $\mathbf{GL}_{2, k} \to U$  correspondant à l'homomorphisme d'anneaux
 $x \mapsto a$,  $y \mapsto b$. Il est facile de voir qu'il s'agit d'un morphisme
surjectif, et donc l'image n'est contenue dans aucun ouvert affine de
 $U$. En fait, le schéma affine  $\mathbf{GL}_{2, k}$  surjecte également sur
 $\mathbf{P}^1_k$, et  $\mathbf{P}^1_k$  n'admet même d'immersion dans
 {\it aucun} schéma affine.
\end{example}

\begin{remark}
\label{remark-quasi-compact-and-quasi-separated}
Soit  $P$  l'une des  $4$  propriétés suivantes : «
quasi-séparé », « séparé », « quasi-compact et quasi-séparé » ou « quasi-compact
et séparé ». Alors les énoncés suivants sont vrais
\begin{enumerate}
\item Tout schéma affine est  $P$.
\item Étant donné un morphisme  $f : X \to Y$  de schémas où  $Y$  est  $P$,
alors  $f$  est  $P$  si et seulement si  $X$  est
 $P$.
\item Si  $X \to Y$  et  $Z \to Y$  sont des morphismes de schémas et  $X$,
 $Y$,  $Z$  ont  $P$, alors  $X \times_Y Z$  a  $P$.
\end{enumerate}
L'énoncé (1) est clair. Soit  $f : X \to Y$  comme dans (2). Si  $f$  a
 $P$, alors  $X$  l'a aussi par les lemmes  \ref{lemma-separated-permanence}  et
 \ref{lemma-composition-quasi-compact}  appliqués à  $X \to Y \to \Spec(\mathbf{Z})$. Si  $X$  a  $P$, alors
 $f$  l'a aussi par les lemmes  \ref{lemma-compose-after-separated}  et  \ref{lemma-quasi-compact-permanence}  appliqués à
 $X \to Y \to \Spec(\mathbf{Z})$. Soient  $X \to Y$  et  $Z \to Y$  comme dans (3). Alors la
projection  $X \times_Y Z \to Z$  a  $P$  comme changement de base de la flèche
 $X \to Y$  qui a  $P$  par (2), voir les lemmes  \ref{lemma-separated-permanence}  et
 \ref{lemma-quasi-compact-preserved-base-change}. Par conséquent,  $X \times_Y Z$  a  $P$  par (2).
\end{remark}





\section{Critère valuatif de séparation}
\label{section-valuative-separatedness}

\begin{lemma}
\label{lemma-separated-implies-valuative}
Soit  $f : X \to S$  un morphisme de schémas. Si  $f$  est séparé, alors
 $f$  satisfait à la partie unicité du critère valuatif.
\end{lemma}

\begin{proof}
Soit un diagramme comme dans la Définition  \ref{definition-valuative-criterion}  donné. Supposons qu'il
existe deux morphismes  $a, b : \Spec(A) \to X$  s'insérant dans le diagramme. Soit
 $Z \subset \Spec(A)$  l'égalisateur de  $a$  et  $b$. D'après le Lemme
 \ref{lemma-where-are-they-equal} , ceci est un sous-schéma fermé de  $\Spec(A)$. Par hypothèse, il
contient le point générique de  $\Spec(A)$. Comme  $A$  est un anneau intègre,
cela implique  $Z = \Spec(A)$. Par conséquent,  $a = b$  comme voulu.
\end{proof}

\begin{lemma}[ Critère valuatif de séparation]
\label{lemma-valuative-criterion-separatedness}
\begin{reference}
\cite[II Proposition 7.2.3]{EGA}
\end{reference}
Soit  $f : X \to S$  un morphisme. Supposons
\begin{enumerate}
\item le morphisme  $f$  est quasi-séparé, et
\item le morphisme  $f$  satisfait la partie unicité du critère valuatif.
\end{enumerate}
Alors  $f$  est séparé.
\end{lemma}

\begin{proof}
Par l'hypothèse (1), la Proposition  \ref{proposition-characterize-universally-closed}, et les lemmes  \ref{lemma-diagonal-immersion}  et
 \ref{lemma-immersion-when-closed} , nous voyons qu'il suffit de prouver que le morphisme  $\Delta_{X/S} : X \to X \times_S X$ 
satisfait la partie existence du critère valuatif. Considérons le diagramme commutatif
$$
\xymatrix{
\Spec(K) \ar[r] \ar[d] & X \ar[d] \\
\Spec(A) \ar[r] \ar@{-->}[ru] & X \times_S X
}
$$
dont les flèches pleines sont données. La flèche en bas à droite correspond à une paire de morphismes  $a, b : \Spec(A) \to X$ 
au-dessus de  $S$. Par (2), nous voyons que  $a = b$. Ainsi,
 $a$  fournit la flèche pointillée cherchée.
\end{proof}





\section{Monomorphismes}
\label{section-monomorphisms}

\begin{definition}
\label{definition-monomorphism}
Un morphisme de schémas est appelé  {\it  monomorphisme
}  s'il est un monomorphisme dans la catégorie des schémas, voir
Catégories, Définition  \ref{categories-definition-mono-epi}.
\end{definition}

\begin{lemma}
\label{lemma-monomorphism}
\begin{slogan}
Un morphisme de schémas est un monomorphisme si et seulement si sa diagonale est un
isomorphisme.
\end{slogan}
Soit  $j : X \to Y$  un morphisme de schémas. Alors  $j$  est un monomorphisme
si et seulement si le morphisme diagonal  $\Delta_{X/Y} : X \to X \times_Y X$  est un isomorphisme.
\end{lemma}

\begin{proof}
Ceci est vrai dans toute catégorie avec des produits fibrés.
\end{proof}

\begin{lemma}
\label{lemma-monomorphism-separated}
Un monomorphisme de schémas est séparé.
\end{lemma}

\begin{proof}
Cela découle du fait qu'un isomorphisme est une immersion fermée et du Lemme
 \ref{lemma-monomorphism}  ci-dessus.
\end{proof}

\begin{lemma}
\label{lemma-composition-monomorphism}
Une composition de monomorphismes est un monomorphisme.
\end{lemma}

\begin{proof}
Vrai dans toute catégorie.
\end{proof}

\begin{lemma}
\label{lemma-base-change-monomorphism}
Le changement de base d'un monomorphisme est un monomorphisme.
\end{lemma}

\begin{proof}
Vrai dans toute catégorie avec des produits fibrés.
\end{proof}

\begin{lemma}
\label{lemma-injective-points}
Soit  $j : X \to Y$  un morphisme de schémas. Si  $j$  est injectif sur les
points, alors  $j$  est séparé.
\end{lemma}

\begin{proof}
Soit  $z$  un point de  $X \times_Y X$. Alors  $x = \text{pr}_1(z)$  et  $\text{pr}_2(z)$ 
sont les mêmes parce que  $j$  envoie ces points sur le même point
 $y$  de  $Y$. Alors nous pouvons choisir un voisinage ouvert
affine  $V \subset Y$  de  $y$  et un voisinage ouvert affine  $U \subset X$  de
 $x$  avec  $j(U) \subset V$. Alors  $z \in U \times_V U \subset X \times_Y X$. Ainsi  $X \times_Y X$  est l'union
des ouverts affines  $U \times_V U$. Puisque  $\Delta_{X/Y}^{-1}(U \times_V U) = U$  et puisque  $U \to U \times_V U$  est
une immersion fermée, nous concluons que  $\Delta_{X/Y}$  est une immersion fermée
(voir l'argument dans la démonstration du Lemme  \ref{lemma-diagonal-immersion}).
\end{proof}

\begin{lemma}
\label{lemma-injective-points-surjective-stalks}
Soit  $j : X \to Y$  un morphisme de schémas. Si
\begin{enumerate}
\item $j$  est injectif sur les points, et
\item pour tout  $x \in X$  l'homomorphisme d'anneaux  $j^\sharp_x : \mathcal{O}_{Y, j(x)} \to \mathcal{O}_{X, x}$  est surjectif,
\end{enumerate}
alors  $j$  est un monomorphisme.
\end{lemma}

\begin{proof}
Soient  $a, b : Z \to X$  deux morphismes de schémas tels que  $j \circ a  = j \circ b$. Alors (1)
implique  $a = b$  en tant qu'applications sous-jacentes des espaces
topologiques. Pour tout  $z \in Z$ , nous avons  $a^\sharp_z \circ j^\sharp_{a(z)} = b^\sharp_z \circ j^\sharp_{b(z)}$  en tant
qu'applications  $\mathcal{O}_{Y, j(a(z))} \to \mathcal{O}_{Z, z}$. La surjectivité des applications  $j^\sharp_x$  impose
 $a^\sharp_z = b^\sharp_z$,  $\forall z \in Z$. Cela implique que  $a^\sharp = b^\sharp$. Ainsi, nous concluons
 $a = b$  en tant que morphismes de schémas comme voulu.
\end{proof}

\begin{lemma}
\label{lemma-immersions-monomorphisms}
Une immersion de schémas est un monomorphisme. En particulier, toute immersion est
séparée.
\end{lemma}

\begin{proof}
Nous pouvons le vérifier en constatant que le critère du Lemme  \ref{lemma-injective-points-surjective-stalks} 
s'applique. Plus élégamment peut-être, nous pouvons utiliser que les lemmes
 \ref{lemma-restrict-map-to-opens}  et  \ref{lemma-characterize-closed-subspace}  impliquent que les immersions ouvertes et fermées
sont des monomorphismes et donc que toute immersion (qui est une composition de
celles-ci) est un monomorphisme.
\end{proof}

\begin{lemma}
\label{lemma-subscheme-of-separated-scheme}
Soit  $f : X \to S$  un morphisme séparé. Tout sous-schéma localement fermé
 $Z \subset X$  est séparé sur  $S$.
\end{lemma}

\begin{proof}
Résulte du Lemme  \ref{lemma-immersions-monomorphisms}  et du fait qu'une composition de morphismes séparés
est séparée (Lemme  \ref{lemma-separated-permanence}).
\end{proof}

\begin{example}
\label{example-Q-over-Z}
Le morphisme  $\Spec(\mathbf{Q}) \to \Spec(\mathbf{Z})$  est un monomorphisme. Cela est vrai parce que
 $\mathbf{Q} \otimes_{\mathbf{Z}} \mathbf{Q} = \mathbf{Q}$. Plus généralement, pour tout schéma  $S$  et tout point
 $s \in S$ , le morphisme canonique
$$
\Spec(\mathcal{O}_{S, s}) \longrightarrow S
$$
est un monomorphisme.
\end{example}

\begin{lemma}
\label{lemma-mono-towards-spec-field}
Soient  $k_1, \ldots, k_n$  des corps. Pour tout monomorphisme de schémas  $X \to \Spec(k_1 \times \ldots \times k_n)$ , il
existe un sous-ensemble  $I \subset \{1, \ldots, n\}$  tel que  $X \cong \Spec(\prod_{i \in I} k_i)$  en tant que schémas sur
 $\Spec(k_1 \times \ldots \times k_n)$. Plus généralement, si  $X = \coprod_{i \in I} \Spec(k_i)$  est une union disjointe de
spectres de corps et que  $Y \to X$  est un monomorphisme, alors il existe un
sous-ensemble  $J \subset I$  tel que  $Y = \coprod_{i \in J} \Spec(k_i)$.
\end{lemma}

\begin{proof}
Commencez par réduire au cas  $n = 1$  (ou  $\# I = 1$) en prenant les images
inverses des sous-schémas ouverts et fermés  $\Spec(k_i)$. Dans ce cas,
 $X$  n'a qu'un seul point, donc il est affine. Le problème algébrique
correspondant est le suivant: si  $k \to R$  est un homomorphisme d'algèbres avec
 $R \otimes_k R \cong R$, alors  $R \cong k$  ou  $R = 0$. Ceci est vrai pour des raisons
de dimension. Voir également Algèbre, Lemme  \ref{algebra-lemma-epimorphism-over-field}
\end{proof}










\section{Fonctorialité pour les modules quasi-cohérents}
\label{section-quasi-coherent}

\noindent
Soit  $X$  un schéma. Nous notons  $\QCoh(\mathcal{O}_X)$  la catégorie des
 $\mathcal{O}_X$-modules quasi-cohérents comme défini dans Modules, Définition
 \ref{modules-definition-quasi-coherent}. Nous avons vu dans la Section  \ref{section-quasi-coherent-affine}  que la catégorie
 $\QCoh(\mathcal{O}_X)$  a beaucoup de bonnes propriétés lorsque  $X$  est affine.
Puisque la propriété d'être quasi-cohérent est locale sur  $X$, ces
propriétés sont héritées par la catégorie des faisceaux quasi-cohérents sur tout
schéma  $X$. Nous les énumérons ici.
\begin{enumerate}
\item Un faisceau de  $\mathcal{O}_X$-modules  $\mathcal{F}$  est quasi-cohérent si et
seulement si la restriction de  $\mathcal{F}$  à chaque ouvert affine  $U = \Spec(R)$ 
est de la forme  $\widetilde M$  pour un certain  $R$-module  $M$.
\item Un faisceau de  $\mathcal{O}_X$-modules  $\mathcal{F}$  est quasi-cohérent si et
seulement si la restriction de  $\mathcal{F}$  à chacun des membres d'un recouvrement
affine ouvert est quasi-cohérente.
\item Toute somme directe de faisceaux quasi-cohérents est quasi-cohérente.
\item Toute colimite de faisceaux quasi-cohérents est quasi-cohérente.
\item
\label{item-kernel}
Le noyau et le conoyau d'un morphisme de faisceaux quasi-cohérents sont
quasi-cohérents.
\item Étant donné une suite exacte courte de  $\mathcal{O}_X$-modules  $0 \to \mathcal{F}_1 \to \mathcal{F}_2 \to \mathcal{F}_3 \to 0$ , si deux
des trois sont quasi-cohérents, le troisième l'est aussi.
\item Étant donné un morphisme de schémas  $f : Y \to X$ , l'image inverse d'un
 $\mathcal{O}_X$-module quasi-cohérent est un  $\mathcal{O}_Y$-module quasi-cohérent. Voir
Modules, Lemme  \ref{modules-lemma-pullback-quasi-coherent}.
\item Étant donné deux  $\mathcal{O}_X$-modules quasi-cohérents, le produit tensoriel est
quasi-cohérent, voir Modules, Lemme  \ref{modules-lemma-tensor-product-permanence}.
\item Étant donné un  $\mathcal{O}_X$-module quasi-cohérent  $\mathcal{F}$ , les algèbres
tensorielle, symétrique et extérieure sur  $\mathcal{F}$  sont quasi-cohérentes, voir
Modules, Lemme  \ref{modules-lemma-whole-tensor-algebra-permanence}.
\item Étant donné deux  $\mathcal{O}_X$-modules quasi-cohérents  $\mathcal{F}$ et  $\mathcal{G}$ 
tels que  $\mathcal{F}$  soit de présentation finie, alors le faisceau Hom interne  $\SheafHom_{\mathcal{O}_X}(\mathcal{F}, \mathcal{G})$ 
est quasi-cohérent, voir Modules, Lemme  \ref{modules-lemma-internal-hom-locally-kernel-direct-sum}  et (\ref{item-kernel}) ci-dessus.
\end{enumerate}
En revanche, l'image directe d'un module quasi-cohérent n'est pas
quasi-cohérente en général. Voici un cas où elle l'est.

\begin{lemma}
\label{lemma-push-forward-quasi-coherent}
Soit  $f : X \to S$  un morphisme de schémas. Si  $f$  est quasi-compact et
quasi-séparé, alors  $f_*$  transforme les  $\mathcal{O}_X$-modules
quasi-cohérents en  $\mathcal{O}_S$-modules quasi-cohérents.
\end{lemma}

\begin{proof}
La question est locale sur  $S$  et donc nous pouvons supposer que
 $S$  est affine. Comme  $X$  est quasi-compact, nous pouvons écrire
 $X = \bigcup_{i = 1}^n U_i$  avec chaque  $U_i$  ouvert affine. Comme  $f$  est
quasi-séparé, nous pouvons écrire  $U_i \cap U_j = \bigcup_{k = 1}^{n_{ij}} U_{ijk}$  pour des ouverts affines
 $U_{ijk}$, voir le Lemme  \ref{lemma-characterize-quasi-separated}. On note  $f_i : U_i \to S$  et  $f_{ijk} : U_{ijk} \to S$  les
restrictions de  $f$. Pour tout ouvert  $V$  de  $S$  et
tout faisceau  $\mathcal{F}$  sur  $X$ , nous avons
\begin{eqnarray*}
f_*\mathcal{F}(V) & = & \mathcal{F}(f^{-1}V) \\
& = &
\Ker\left(
\bigoplus\nolimits_i \mathcal{F}(f^{-1}V \cap U_i)
\to
\bigoplus\nolimits_{i, j, k} \mathcal{F}(f^{-1}V \cap U_{ijk})\right) \\
& = &
\Ker\left(
\bigoplus\nolimits_i f_{i, *}(\mathcal{F}|_{U_i})(V)
\to
\bigoplus\nolimits_{i, j, k} f_{ijk, *}(\mathcal{F}|_{U_{ijk}})(V)\right) \\
& = &
\Ker\left(
\bigoplus\nolimits_i f_{i, *}(\mathcal{F}|_{U_i})
\to
\bigoplus\nolimits_{i, j, k} f_{ijk, *}(\mathcal{F}|_{U_{ijk}})\right)(V)
\end{eqnarray*}
En d'autres termes, il existe une suite exacte de faisceaux
$$
0 \to f_*\mathcal{F}
\to \bigoplus f_{i, *}\mathcal{F}_i
\to \bigoplus f_{ijk, *}\mathcal{F}_{ijk}
$$
où  $\mathcal{F}_i, \mathcal{F}_{ijk}$  désignent les restrictions de  $\mathcal{F}$  à l'ouvert correspondant.
Si  $\mathcal{F}$  est un  $\mathcal{O}_X$-module quasi-cohérent alors  $\mathcal{F}_i$  est
un  $\mathcal{O}_{U_i}$-module quasi-cohérent et  $\mathcal{F}_{ijk}$  est un  $\mathcal{O}_{U_{ijk}}$-module
quasi-cohérent. Ainsi, par le Lemme  \ref{lemma-widetilde-pullback} , nous voyons que le deuxième et
le troisième terme de la suite exacte sont des  $\mathcal{O}_S$-modules
quasi-cohérents. Nous concluons donc que  $f_*\mathcal{F}$  est un  $\mathcal{O}_S$-module
quasi-cohérent.
\end{proof}

\noindent
En utilisant cela, nous pouvons caractériser les immersions (fermées) de schémas
comme suit.

\begin{lemma}
\label{lemma-characterize-closed-immersions}
Soit  $f : X \to Y$  un morphisme de schémas. Supposons que
\begin{enumerate}
\item $f$  induise un homéomorphisme de  $X$  avec un sous-ensemble fermé
de  $Y$, et
\item $f^\sharp : \mathcal{O}_Y \to f_*\mathcal{O}_X$  soit surjectif.
\end{enumerate}
Alors  $f$  est une immersion fermée de schémas.
\end{lemma}

\begin{proof}
Supposons (1) et (2). D'après (1), le morphisme  $f$  est quasi-compact
(voir Topologie, Lemme  \ref{topology-lemma-closed-in-quasi-compact}). Les conditions (1) et (2) impliquent les
conditions (1) et (2) du Lemme  \ref{lemma-injective-points-surjective-stalks}. Par conséquent,  $f : X \to Y$  est un
monomorphisme. En particulier,  $f$  est séparé, voir Lemme  \ref{lemma-monomorphism-separated}.
Ainsi, le Lemme  \ref{lemma-push-forward-quasi-coherent}  ci-dessus s'applique et nous concluons que
 $f_*\mathcal{O}_X$  est un  $\mathcal{O}_Y$-module quasi-cohérent. Par conséquent, le noyau
de  $\mathcal{O}_Y \to f_*\mathcal{O}_X$  est quasi-cohérent selon le Lemme  \ref{lemma-extension-quasi-coherent}. Comme un faisceau
quasi-cohérent est localement engendré par des sections (voir Modules, Définition
 \ref{modules-definition-quasi-coherent}), cela implique que  $f$  est une immersion fermée, voir
Définition  \ref{definition-closed-immersion-locally-ringed-spaces}.
\end{proof}

\noindent
Nous pouvons utiliser ce lemme pour démontrer le lemme suivant.

\begin{lemma}
\label{lemma-composition-immersion}
Une composition d'immersions de schémas est une immersion, une composition
d'immersions fermées de schémas est une immersion fermée, et une composition
d'immersions ouvertes de schémas est une immersion ouverte.
\end{lemma}

\begin{proof}
Cela est évident dans le cas des immersions ouvertes puisqu'un sous-espace ouvert
d'un sous-espace ouvert est également un sous-espace ouvert.

\medskip\noindent
Supposons que  $a : Z \to Y$  et  $b : Y \to X$  soient des immersions fermées de
schémas. Nous allons vérifier que  $c = b \circ a$  est également une immersion fermée.
L'hypothèse implique que  $a$  et  $b$  sont des homéomorphismes sur
des sous-ensembles fermés, et donc  $c = b \circ a$  est aussi un homéomorphisme sur un
sous-ensemble fermé. De plus, l'application  $\mathcal{O}_X \to c_*\mathcal{O}_Z$  est surjective
puisqu'elle se factorise comme la composition des applications surjectives
 $\mathcal{O}_X \to b_*\mathcal{O}_Y$  et  $b_*\mathcal{O}_Y \to b_*a_*\mathcal{O}_Z$  (surjectives car  $b_*$  est exact, voir
Modules, Lemme  \ref{modules-lemma-i-star-exact}). Ainsi, d'après le Lemme  \ref{lemma-characterize-closed-immersions}  ci-dessus,
 $c$  est une immersion fermée.

\medskip\noindent
Enfin, nous arrivons au cas des immersions. Supposons que  $a : Z \to Y$  et
 $b : Y \to X$  soient des immersions de schémas. Cela signifie qu’il existe des
sous-schémas ouverts  $V \subset Y$  et  $U \subset X$  tels que  $a(Z) \subset V$,
 $b(Y) \subset U$  et  $a : Z \to V$  et  $b : Y \to U$  soient des immersions fermées.
Puisque la topologie sur  $Y$  est induite à partir de la topologie sur
 $U$ , nous pouvons trouver un ouvert  $U' \subset U$  tel que  $V = b^{-1}(U')$.
Ensuite, nous voyons que  $Z \to V = b^{-1}(U') \to U'$  est une composition d’immersions fermées et
est donc une immersion fermée. Cela prouve que  $Z \to X$  est une immersion, ce qui achève la démonstration.
\end{proof}











\begin{multicols}{2}[\section{Autres chapitres}]
\noindent
Préliminaires
\begin{enumerate}
\item \hyperref[introduction-section-phantom]{Introduction}
\item \hyperref[conventions-section-phantom]{Conventions}
\item \hyperref[sets-section-phantom]{Théorie des ensembles}
\item \hyperref[categories-section-phantom]{Catégories}
\item \hyperref[topology-section-phantom]{Topologie}
\item \hyperref[sheaves-section-phantom]{Faisceaux sur les espaces}
\item \hyperref[sites-section-phantom]{Sites et faisceaux}
\item \hyperref[stacks-section-phantom]{Champs}
\item \hyperref[fields-section-phantom]{Corps}
\item \hyperref[algebra-section-phantom]{Algèbre commutative}
\item \hyperref[brauer-section-phantom]{Groupes de Brauer}
\item \hyperref[homology-section-phantom]{Algèbre homologique}
\item \hyperref[derived-section-phantom]{Catégories dérivées}
\item \hyperref[simplicial-section-phantom]{Méthodes simpliciales}
\item \hyperref[more-algebra-section-phantom]{Compléments d'algèbre}
\item \hyperref[smoothing-section-phantom]{Lissage des morphismes d'anneaux}
\item \hyperref[modules-section-phantom]{Faisceaux de modules}
\item \hyperref[sites-modules-section-phantom]{Modules sur les sites}
\item \hyperref[injectives-section-phantom]{Objets injectifs}
\item \hyperref[cohomology-section-phantom]{Cohomologie des faisceaux}
\item \hyperref[sites-cohomology-section-phantom]{Cohomologie sur les sites}
\item \hyperref[dga-section-phantom]{Algèbre différentielle graduée}
\item \hyperref[dpa-section-phantom]{Algèbre à puissances divisées}
\item \hyperref[sdga-section-phantom]{Faisceaux différentiels gradués}
\item \hyperref[hypercovering-section-phantom]{Hyperrecouvrements}
\end{enumerate}
Schémas
\begin{enumerate}
\setcounter{enumi}{25}
\item \hyperref[schemes-section-phantom]{Schémas}
\item \hyperref[constructions-section-phantom]{Constructions de schémas}
\item \hyperref[properties-section-phantom]{Propriétés des schémas}
\item \hyperref[morphisms-section-phantom]{Morphismes de schémas}
\item \hyperref[coherent-section-phantom]{Cohomologie des schémas}
\item \hyperref[divisors-section-phantom]{Diviseurs}
\item \hyperref[limits-section-phantom]{Limites de schémas}
\item \hyperref[varieties-section-phantom]{Variétés}
\item \hyperref[topologies-section-phantom]{Topologies sur les schémas}
\item \hyperref[descent-section-phantom]{Descente}
\item \hyperref[perfect-section-phantom]{Catégories dérivées de schémas}
\item \hyperref[more-morphisms-section-phantom]{Compléments sur les morphismes}
\item \hyperref[flat-section-phantom]{Compléments sur la platitude}
\item \hyperref[groupoids-section-phantom]{Schémas en groupoïdes}
\item \hyperref[more-groupoids-section-phantom]{Compléments sur les schémas en groupoïdes}
\item \hyperref[etale-section-phantom]{Morphismes étales de schémas}
\end{enumerate}
Thèmes de théorie des schémas
\begin{enumerate}
\setcounter{enumi}{41}
\item \hyperref[chow-section-phantom]{Homologie de Chow}
\item \hyperref[intersection-section-phantom]{Théorie de l'intersection}
\item \hyperref[pic-section-phantom]{Schémas de Picard des courbes}
\item \hyperref[weil-section-phantom]{Théories cohomologiques de Weil}
\item \hyperref[adequate-section-phantom]{Modules adéquats}
\item \hyperref[dualizing-section-phantom]{Complexes dualisants}
\item \hyperref[duality-section-phantom]{Dualité pour les schémas}
\item \hyperref[discriminant-section-phantom]{Discriminants et différentes}
\item \hyperref[derham-section-phantom]{Cohomologie de de Rham}
\item \hyperref[local-cohomology-section-phantom]{Cohomologie locale}
\item \hyperref[algebraization-section-phantom]{Géométrie algébrique et formelle}
\item \hyperref[curves-section-phantom]{Courbes algébriques}
\item \hyperref[resolve-section-phantom]{Résolution des surfaces}
\item \hyperref[models-section-phantom]{Réduction semi-stable}
\item \hyperref[functors-section-phantom]{Foncteurs et morphismes}
\item \hyperref[equiv-section-phantom]{Catégories dérivées de variétés}
\item \hyperref[pione-section-phantom]{Groupes fondamentaux de schémas}
\item \hyperref[etale-cohomology-section-phantom]{Cohomologie étale}
\item \hyperref[crystalline-section-phantom]{Cohomologie cristalline}
\item \hyperref[proetale-section-phantom]{Cohomologie pro-étale}
\item \hyperref[relative-cycles-section-phantom]{Cycles relatifs}
\item \hyperref[more-etale-section-phantom]{Compléments de cohomologie étale}
\item \hyperref[trace-section-phantom]{La formule des traces}
\end{enumerate}
Espaces algébriques
\begin{enumerate}
\setcounter{enumi}{64}
\item \hyperref[spaces-section-phantom]{Espaces algébriques}
\item \hyperref[spaces-properties-section-phantom]{Propriétés des espaces algébriques}
\item \hyperref[spaces-morphisms-section-phantom]{Morphismes d'espaces algébriques}
\item \hyperref[decent-spaces-section-phantom]{Espaces algébriques décents}
\item \hyperref[spaces-cohomology-section-phantom]{Cohomologie des espaces algébriques}
\item \hyperref[spaces-limits-section-phantom]{Limites d'espaces algébriques}
\item \hyperref[spaces-divisors-section-phantom]{Diviseurs sur les espaces algébriques}
\item \hyperref[spaces-over-fields-section-phantom]{Espaces algébriques sur des corps}
\item \hyperref[spaces-topologies-section-phantom]{Topologies sur les espaces algébriques}
\item \hyperref[spaces-descent-section-phantom]{Descente et espaces algébriques}
\item \hyperref[spaces-perfect-section-phantom]{Catégories dérivées d'espaces}
\item \hyperref[spaces-more-morphisms-section-phantom]{Compléments sur les morphismes d'espaces}
\item \hyperref[spaces-flat-section-phantom]{Platitude sur les espaces algébriques}
\item \hyperref[spaces-groupoids-section-phantom]{Groupoïdes dans les espaces algébriques}
\item \hyperref[spaces-more-groupoids-section-phantom]{Compléments sur les groupoïdes dans les espaces}
\item \hyperref[bootstrap-section-phantom]{Amorçage}
\item \hyperref[spaces-pushouts-section-phantom]{Sommes amalgamées d'espaces algébriques}
\end{enumerate}
Thèmes de géométrie
\begin{enumerate}
\setcounter{enumi}{81}
\item \hyperref[spaces-chow-section-phantom]{Groupes de Chow des espaces}
\item \hyperref[groupoids-quotients-section-phantom]{Quotients de groupoïdes}
\item \hyperref[spaces-more-cohomology-section-phantom]{Compléments sur la cohomologie des espaces}
\item \hyperref[spaces-simplicial-section-phantom]{Espaces simpliciaux}
\item \hyperref[spaces-duality-section-phantom]{Dualité pour les espaces}
\item \hyperref[formal-spaces-section-phantom]{Espaces algébriques formels}
\item \hyperref[restricted-section-phantom]{Algébrisation des espaces formels}
\item \hyperref[spaces-resolve-section-phantom]{Résolution des surfaces revisitée}
\end{enumerate}
Théorie des déformations
\begin{enumerate}
\setcounter{enumi}{89}
\item \hyperref[formal-defos-section-phantom]{Théorie formelle des déformations}
\item \hyperref[defos-section-phantom]{Théorie des déformations}
\item \hyperref[cotangent-section-phantom]{Le complexe cotangent}
\item \hyperref[examples-defos-section-phantom]{Problèmes de déformation}
\end{enumerate}
Champs algébriques
\begin{enumerate}
\setcounter{enumi}{93}
\item \hyperref[algebraic-section-phantom]{Champs algébriques}
\item \hyperref[examples-stacks-section-phantom]{Exemples de champs}
\item \hyperref[stacks-sheaves-section-phantom]{Faisceaux sur les champs algébriques}
\item \hyperref[criteria-section-phantom]{Critères de représentabilité}
\item \hyperref[artin-section-phantom]{Axiomes d'Artin}
\item \hyperref[quot-section-phantom]{Espaces de Quot et de Hilbert}
\item \hyperref[stacks-properties-section-phantom]{Propriétés des champs algébriques}
\item \hyperref[stacks-morphisms-section-phantom]{Morphismes de champs algébriques}
\item \hyperref[stacks-limits-section-phantom]{Limites de champs algébriques}
\item \hyperref[stacks-cohomology-section-phantom]{Cohomologie des champs algébriques}
\item \hyperref[stacks-perfect-section-phantom]{Catégories dérivées de champs}
\item \hyperref[stacks-introduction-section-phantom]{Introduction aux champs algébriques}
\item \hyperref[stacks-more-morphisms-section-phantom]{Compléments sur les morphismes de champs}
\item \hyperref[stacks-geometry-section-phantom]{La géométrie des champs}
\end{enumerate}
Thèmes de théorie des espaces de modules
\begin{enumerate}
\setcounter{enumi}{107}
\item \hyperref[moduli-section-phantom]{Champs de modules}
\item \hyperref[moduli-curves-section-phantom]{Espaces de modules de courbes}
\end{enumerate}
Divers
\begin{enumerate}
\setcounter{enumi}{109}
\item \hyperref[examples-section-phantom]{Exemples}
\item \hyperref[exercises-section-phantom]{Exercices}
\item \hyperref[guide-section-phantom]{Guide bibliographique}
\item \hyperref[desirables-section-phantom]{Desiderata}
\item \hyperref[coding-section-phantom]{Style de programmation}
\item \hyperref[obsolete-section-phantom]{Obsolète}
\item \hyperref[fdl-section-phantom]{Licence de documentation libre GNU}
\item \hyperref[index-section-phantom]{Index généré automatiquement}
\end{enumerate}
\end{multicols}


\bibliography{my}
\bibliographystyle{amsalpha}

\end{document}
